% AUTO-GENERATED by scripts/build_lectures.py — do not edit.
% Edit lectures/bodies/*.tex or lectures/preamble.tex instead.
\documentclass[11pt,openany]{report}
\usepackage[a4paper, margin=1in, top=1.1in, bottom=1.1in]{geometry}
% ---------------------------------------------------------------------------
%  Shared preamble for the "Вероятности и Статистика" lecture notes.
%
%  Used by both builds:
%    * lecture_NN.tex   — standalone article, one lecture   (\section = top level)
%    * lectures_full.tex — collected book, report class     (\chapter = per lecture)
%
%  The driver sets \numberwithin{thm}{section|chapter} after \input-ing this.
% ---------------------------------------------------------------------------

\usepackage{fontspec}
\setmainfont{Times New Roman}

% Bulgarian: loads hyphenation patterns (without these TeX cannot break a single
% Cyrillic word, which is what produced the loose lines and overfull boxes) and
% translates the auto-generated captions ("Contents" -> "Съдържание").
\usepackage{polyglossia}
\setmainlanguage{bulgarian}
\setotherlanguage{english}

\usepackage{amsmath}
\usepackage{amssymb}
\usepackage{amsfonts}
\usepackage{amsthm}
\usepackage[most]{tcolorbox}
\usepackage{microtype}
\usepackage{fancyhdr}

% "г.", "т.е.", "п.с." are abbreviations, not sentence ends — suppress the
% extra inter-sentence space LaTeX would otherwise insert after them.
\frenchspacing

\linespread{1.08}
\setlength{\parskip}{0.4em}
\setlength{\parindent}{0pt}
\clubpenalty=10000       % no orphans
\widowpenalty=10000      % no widows
\displaywidowpenalty=10000

% Long unbreakable inline-math runs (\{1,2,\dots,n\}, \text{...} spans) cannot be
% hyphenated. Let TeX stretch a line as a last resort rather than let it stick
% out into the margin.
\setlength{\emergencystretch}{3em}

% --- palette ---------------------------------------------------------------
\definecolor{defframe}{RGB}{38,90,160}
\definecolor{defback}{RGB}{234,241,250}
\definecolor{thmframe}{RGB}{34,120,70}
\definecolor{thmback}{RGB}{233,246,238}

% --- numbered statement environments ---------------------------------------
% One shared counter for all five, so a reader scanning for "Твърдение 4.7"
% never has to guess which of five sequences it belongs to.
\newtheoremstyle{lecstyle}
  {\topsep}{\topsep}     % space above / below
  {\normalfont}          % body font — upright; italic Cyrillic at 11pt is hard to read
  {0pt}                  % indent
  {\bfseries}            % head font
  {.}                    % punctuation after head
  {0.5em}                % space after head
  {}                     % default head spec
\theoremstyle{lecstyle}

\newtheorem{thm}{Теорема}
\newtheorem{defn}[thm]{Дефиниция}
\newtheorem{prop}[thm]{Твърдение}
\newtheorem{lem}[thm]{Лема}
\newtheorem{cor}[thm]{Следствие}

\newtheoremstyle{lecremark}
  {\topsep}{\topsep}{\normalfont}{0pt}{\itshape}{.}{0.5em}{}
\theoremstyle{lecremark}
\newtheorem{remark}[thm]{Забележка}
\newtheorem{example}[thm]{Пример}

% Definitions in blue, everything provable in green. Proofs stay outside.
\tcbset{lecbox/.style={enhanced, breakable, arc=2.4mm, boxrule=0.7pt,
    left=2.6mm, right=2.6mm, top=1.8mm, bottom=1.8mm, before skip=0.6em,
    after skip=0.6em}}
\tcolorboxenvironment{defn}{lecbox, colback=defback, colframe=defframe}
\tcolorboxenvironment{thm} {lecbox, colback=thmback, colframe=thmframe}
\tcolorboxenvironment{prop}{lecbox, colback=thmback, colframe=thmframe}
\tcolorboxenvironment{lem} {lecbox, colback=thmback, colframe=thmframe}
\tcolorboxenvironment{cor} {lecbox, colback=thmback, colframe=thmframe}

% Proofs: unboxed, with an automatic flush-right QED on every one of them.
\renewcommand{\proofname}{Доказателство}
\renewcommand{\qedsymbol}{$\square$}

% --- links ------------------------------------------------------------------
\usepackage{hyperref}   % load last
\hypersetup{
  colorlinks=true, linkcolor=defframe, urlcolor=defframe, citecolor=defframe,
  linktoc=all, bookmarksnumbered=true, bookmarksopen=true, unicode=true,
  bookmarksdepth=3,   % printed TOC stays lean; the PDF outline keeps full depth
  pdfauthor={}}


%% (Caption names are set at the top of the document body — polyglossia installs
%%  its Bulgarian captions at \begin{document} and would overwrite preamble
%%  \renewcommand declarations.)
\numberwithin{thm}{chapter}
\numberwithin{equation}{chapter}

%% Master contents lists lectures and their sections; the PDF bookmark tree
%% still carries the full subsection depth.
\setcounter{tocdepth}{1}
\setcounter{secnumdepth}{2}

\hypersetup{pdftitle={Вероятности и Статистика — лекционни записки}}

\pagestyle{fancy}
\fancyhf{}
\fancyhead[L]{\small\itshape\nouppercase{\leftmark}}
\fancyhead[R]{\small\thepage}
\renewcommand{\headrulewidth}{0.4pt}
\fancypagestyle{plain}{\fancyhf{}\fancyhead[R]{\small\thepage}%
  \renewcommand{\headrulewidth}{0pt}}

\begin{document}

%% Lectures are the chapters of this book.
\renewcommand{\chaptername}{Лекция}
\renewcommand{\proofname}{Доказателство}
\renewcommand{\crosslecture}[2]{\hyperref[#1]{#2}}
\hypersetup{pageanchor=false}

\begin{titlepage}
\centering
\vspace*{5cm}
{\Huge\bfseries Вероятности и Статистика\par}
\vspace{1.2cm}
{\Large Лекционни записки\par}
\vspace{0.5cm}
{\large Специалност „Софтуерно инженерство“, 2021/2022\par}
\vspace{2.5cm}
{\large 14 лекции и едно упражнение\par}
\vfill
\end{titlepage}

\hypersetup{pageanchor=true}
\pagenumbering{roman}
\tableofcontents
\chapter*{Организационна бележка}
\addcontentsline{toc}{chapter}{Организационна бележка}

Материалът по комбинаторика в лекция 15 е преговор на знания, които се
предполагат познати от курса по дискретна математика.

Изискванията за изпита, независимо дали е онлайн или присъствен, не са твърде
високи, но се очаква базова грамотност: да не се допускат груби грешки като
отрицателни дисперсии и непознаване на основни понятия като математическо
очакване и плътност. Консултация преди изпита ще бъде организирана, ако е
възможно.


\chapter*{Основни означения}
\addcontentsline{toc}{chapter}{Основни означения}

\[
\begin{array}{c|l}
\text{Означение} & \text{Значение} \\ \hline
\E X & \text{математическо очакване на }X \\
\Var X & \text{дисперсия на }X \\
\ind_A & \text{индикатор на събитието }A \\
\mathbb{P}(A\given B) & \text{условна вероятност на }A\text{ при }B \\
\Ber(p) & \text{разпределение на Бернули с параметър }p \\
\Bin(n,p) & \text{биномно разпределение с параметри }n,p \\
\Poi(\lambda) & \text{Поасоново разпределение с параметър }\lambda
\end{array}
\]

\clearpage
\pagenumbering{arabic}

\chapter{Въведение. Случайни експерименти, събития и \texorpdfstring{$\sigma$}{σ}-алгебри}
\section{Въведение и мотивация за курса}

Този курс има за цел да изгради у вас така нареченото стохастично (вероятностно) мислене. Въпреки че курсът е чисто математически и ще изисква усвояване на строги дефиниции и теоретичен апарат, този поглед към света е изключително полезен и намира приложение на най-високо ниво във фундаменталната наука, в съвременните технологии (като изкуствен интелект и машинно обучение), както и в напълно ежедневни, злободневни ситуации.

\subsection{Примери за приложение на стохастиката}

За да се мотивира необходимостта от такъв апарат, ще разгледаме няколко примера, които показват мощта на вероятностните модели.

\textbf{Брауново движение.} През 1827 г. ботаникът Робърт Браун (Р. Браун 1827; 1905) наблюдава хаотичното движение на частица полен, поставена на водна повърхност. На онзи етап от развитието на науката атомарната теория все още не е била утвърдена и не е било ясно дали това движение се дължи на вътрешна енергия на частицата, или на външната среда. Едва през 1905 г. Алберт Айнщайн успява да опише математически това движение, стъпвайки на кинетичната теория. Резултатът постулира, че в резултат от ударите на водните молекули от всички посоки, частицата приема хаотична траектория, която постоянно мени посоката си поради моментния дисбаланс на силите.
Любопитно е, че ако разгледаме Брауновото движение като чисто стохастичен обект, траекторията е непрекъсната, но никъде диференцируема и има безкрайна дължина. Въпреки тази физическа „нерационалност“, моделът е изключително добро приближение на реалността. Дори днес стохастичните диференциални уравнения са в основата на моделирането на сложни физични и финансови феномени.


\textbf{Тотализаторът и шансът.} На 6 септември 2009 г. в България се случва събитие, което предизвиква огромен медиен скандал: в два последователни тиража на тотализатора „6 от 49“ се падат едни и същи числа. Примерът е взет от книга на колеги от Imperial College, които го използват, за да покажат как дори елементарни вероятностни разсъждения помагат при анализа на реални събития. Обикновено първосигналната реакция е да се търси измама и манипулация. Въпреки това, ако се направи елементарен статистически анализ (взимайки предвид, че до онзи момент са изиграни около 10 000 тиража), се оказва, че вероятността за подобно събитие, макар и малка, съвсем не е пренебрежима. Разбирането на тези базови вероятности ни позволява да оценяваме обективно фактите, вместо да се поддаваме на емоционални спекулации. Допълнителен аргумент против хипотезата за измама е фактът, че във втория тираж печалбата е била разделена между 14 души — едва ли някой би организирал измама, за да дели печалбата с още толкова хора.


\textbf{Статистика и актуални проблеми.} В контекста на пандемията от COVID-19 се появяват множество спекулации. Например, когато се съобщи, че около 4-5\% от починалите за деня са били ваксинирани, веднага започват конспиративни теории за фалшиви сертификати или за пълна неефективност на ваксините. Отново, един елементарен статистически анализ на базата на съотношението между ваксинирани и неваксинирани в обществото би показал реалната ефективност на ваксините, която, макар и не 100\%, съществува.


\textbf{Фундаментална математика.} Преди няколко години известният математик Терънс Тао доказа в голяма степен хипотезата на българския академик Благовест Сендов. Ключов момент от неговата изключително сложна работа отново беше базиран на стохастиката.


\section{Елементарни събития и пространство от елементарни събития}

За да формализираме тези идеи и да започнем да градим математическата теория, трябва да въведем основните понятия, свързани със случайните явления.

\begin{defn}
Под \emph{случаен експеримент} разбираме изследване, опит или наблюдение, чийто точен изход не е предварително известен (не е предопределен). Например, при хвърляне на монета не знаем дали ще се падне „ези“ или „тура“, защото не разполагаме с цялата необходима информация (сили, въздушно съпротивление, начален импулс и др.).
\end{defn}

\begin{defn}
Възможните изходи от даден случаен експеримент се наричат \emph{елементарни събития} и обикновено се обозначават с малката гръцка буква $\omega$.
\end{defn}

\begin{defn}
За даден случаен експеримент, множеството от всички възможни елементарни събития се нарича \emph{пространство от елементарни събития} и се означава с главно $\Omega$.
\end{defn}

Нека разгледаме няколко конкретни примера за пространства от елементарни събития:

\textbf{Хвърляне на монета.}
Експеримент с бинарен изход.
\[
\Omega = \{ \text{'ези'}; \text{'тура'} \}
\]
За удобство и компютърна обработка, много често тези два изхода се кодират с 0 и 1:
\[
\Omega = \{ 0; 1 \}
\]


\textbf{Хвърляне на зарче.}
Експериментът се състои в хвърляне на стандартен шестстенен зар, независимо дали е честен или не. Възможните изходи са броят точки на горната стена:
\[
\Omega = \{ 1, 2, 3, 4, 5, 6 \}
\]
Едно конкретно елементарно събитие, например падането на числото 5, се записва като $\omega_5 = \{ 5 \}$.


\textbf{Един тираж на тото „6 от 49“.}
Тук експериментът е изтеглянето на шест числа от общо 49. Всяко елементарно събитие е една конкретна (неподредена) шесторка от числа.
\[
\Omega = \{ \omega_1, \omega_2, \dots, \omega_N \}
\]
Броят на възможните изходи (кардиналността на множеството) е комбинация от 49 елемента 6-ти клас:
\[
N = \binom{49}{6} = \underline{\underline{13983816}}
\]


\textbf{Неизброимо множество — време на живот.}
Да си представим, че изследваме времето на живот на един процесор. То може да бъде всяко неотрицателно реално число (дори 0, ако е фабрично дефектен). В този случай пространството е непрекъснато (неизброимо):
\[
\Omega = \mathbb{R}_+ = [0, \infty)
\]


\textbf{Брауново движение.}
За полена върху водната повърхност, $\Omega$ се състои от всички възможни непрекъснати криви (функции), които стартират от началото на координатната система и описват траектория. Това е изключително сложно и голямо функционално пространство. За щастие, в повечето практически задачи не е нужно да описваме експлицитно цялата тази структура.


\section{Събития и операции с тях}

В реалните проблеми рядко се интересуваме от това дали се е случило точно едно конкретно елементарно събитие. Много по-често въпросите ни касаят \emph{група} от елементарни събития, които споделят общ признак.

\begin{defn}
Всяко подмножество $A$ на пространството от елементарни събития $\Omega$ ($A \subseteq \Omega$) се нарича \emph{събитие}.
\end{defn}

Например, ако разглеждаме клиентите, влизащи в един магазин, $\Omega$ може да съдържа всички възможни хора. Събитието $A = \{\text{жени}\}$ обединява всички елементарни събития, които отговарят на признака пол. Магазинът обикновено не се интересува дали точно Иван или Деница е влязъл, а по-скоро от общите характеристики на клиентите (напр. мъж/жена).
Ако хвърляме зарче, събитието $A = \{ 1, 3, 5 \}$ представлява падането на нечетно число. Групирали сме изходите по признака „нечетност“.

\subsection{Операции със събития}

Тъй като събитията са множества, с тях можем да извършваме стандартните теоретико-множествени операции, които имат ясен вероятностен смисъл. Нека $A, B \subseteq \Omega$.

\textbf{Влагане (Подмножество):}
\[
A \subseteq B
\]
Казваме, че $A$ е подмножество на $B$, ако настъпването на $A$ гарантира настъпването на $B$. Формално: ако $\omega \in A \implies \omega \in B$. (Събитията са равни, $A = B$, ако $A \subseteq B$ и $B \subseteq A$).

\textbf{Обединение (Логическо „ИЛИ“):}
\[
C = A \cup B
\]
Това е събитието $C$, такова че ако $\omega \in C$, то $\omega \in A$ \textbf{или} $\omega \in B$.
Например, ако $A = \{\text{гласували за партия } X\}$, а $B = \{\text{хора на възраст 19-21 год.}\}$, то обединението им $A \cup B$ съдържа всички избиратели, които са или гласували за партия $X$, или са на възраст между 19 и 21 години (или и двете едновременно).

\textbf{Сечение (Логическо „И“):}
\[
C = A \cap B
\]
Това е събитието $C$, такова че ако $\omega \in C$, то $\omega \in A$ \textbf{и} $\omega \in B$.
В същия пример, сечението $A \cap B$ ще бъдат точно онези избиратели на възраст между 19 и 21 години, които са гласували за партия $X$. За една партия е важно да анализира тези сечения, за да знае какви са исканията на тази конкретна демографска група.

\textbf{Допълнение (Отрицание):}
\[
A^c = \Omega \setminus A
\]
Събитието $A^c$ се случва тогава и само тогава, когато събитието $A$ не се случва. Формално: ако $\omega \in A^c$, то $\omega \notin A$.
Разбира се, в сила е:
\[
\Omega = A \cup A^c
\]
Примери за допълнение:
Ако $A = \{\text{жени}\}$, то $A^c = \{\text{мъже}\}$.
Ако $A = \{1, 3, 5\}$, то $A^c = \{2, 4, 6\}$.
Ако $A = \{\text{гласували за } X\}$, то $A^c = \{\text{негласували за } X\}$.

\subsection{Свойства на операциите}

Свойствата на тези операции са интуитивни и добре познати от дискретната математика:

\textbf{а) Комутативност:}
\[
A \cup B = B \cup A; \quad A \cap B = B \cap A
\]

\textbf{б) Асоциативност:}
\[
A \cup (B \cup C) = (A \cup B) \cup C = A \cup B \cup C
\]
\[
A \cap (B \cap C) = (A \cap B) \cap C = A \cap B \cap C
\]

\textbf{в) Дистрибутивност:}
\[
A \cup (B \cap C) = (A \cup B) \cap (A \cup C)
\]

\textbf{г) Инволюция на допълнението:}
Ако вземем допълнението на допълнението, се връщаме в изходното множество:
\[
(A^c)^c = A
\]

\textbf{д) Безкрайни обединения и сечения:}
Ако $A_i$ за $i \ge 1$ е редица от събития (изброимо много), можем да дефинираме тяхното обединение и сечение:
\[
\bigcup_{i=1}^{\infty} A_i = \{ \omega \in \Omega : \omega \text{ принадлежи на поне едно } A_i \}
\]
\[
\bigcap_{i=1}^{\infty} A_i = \{ \omega \in \Omega : \omega \text{ принадлежи на } \forall A_i \}
\]
Тези безкрайни операции са фундаментални за аксиоматичното изграждане на теорията на вероятностите (подобно на границите в математическия анализ), въпреки че в практиката най-често работим с крайни множества.

\textbf{е) Закони на Де Морган (De Morgan):}
Законите на Де Морган свързват обединението, сечението и допълнението:
За две събития:
\[
(A \cup B)^c = A^c \cap B^c
\]
За изброимо много събития:
\[
\left(\bigcup_{i=1}^\infty A_i\right)^c = \bigcap_{i=1}^\infty A_i^c
\]
\[
\left(\bigcap_{i=1}^{\infty} A_i\right)^c = \bigcup_{i=1}^{\infty} A_i^c
\]
С думи: допълнението на обединението е сечение от допълненията, а допълнението на сечението е обединение от допълненията.

\section{Сигма-алгебра ($\sigma$-алгебра)}

За да дефинираме коректно вероятностна мярка, не винаги можем просто да разглеждаме всички възможни подмножества на $\Omega$. Необходима ни е математическа структура, която да съдържа събитията, които можем да „измерим“, и тази структура трябва да бъде затворена относно определени операции (за да не се окаже, че ако знаем вероятността на събитията $A$ и $B$, не можем да сметнем вероятността на тяхното обединение). Тази структура се нарича $\sigma$-алгебра.

\begin{keydefn}[$\sigma$-алгебра]\label{def:sigma-algebra}
Нека $\Omega$ е пространство от елементарни събития. Нека $\mathcal{A}$ е колекция от подмножества на $\Omega$. Тогава $\mathcal{A}$ се нарича \emph{$\sigma$-алгебра}, ако са изпълнени следните три условия:
\begin{itemize}
    \item[(i)] $\emptyset \in \mathcal{A}$ \quad (Празното множество е част от колекцията.)
    \item[(ii)] Ако $A \in \mathcal{A}$, то $A^c \in \mathcal{A}$ \quad (Колекцията е затворена относно операцията допълнение.)
    \item[(iii)] Ако $A_i \in \mathcal{A}$ за $\forall i \ge 1$, то $\bigcup_{i=1}^{\infty} A_i \in \mathcal{A}$ \quad (Колекцията е затворена относно изброими обединения.)
\end{itemize}
\end{keydefn}

Затвореността за обединение на две множества $A \cup B \in \mathcal{A}$ следва от (iii), като приемем, че останалите множества от безкрайната редица са празни.

От тези три аксиоми можем да изведем важно следствие за операцията сечение, която не фигурира експлицитно в дефиницията.

\begin{cor}
Нека $\mathcal{A}$ е $\sigma$-алгебра и $A_i \in \mathcal{A}$ за $\forall i \ge 1$. Тогава:
\[
B = \bigcap_{i=1}^{\infty} A_i \in \mathcal{A}
\]
Тоест, $\sigma$-алгебрата е затворена и относно изброими сечения.
\end{cor}

\begin{proof}
Искаме да покажем, че $\mathcal{A} \ni B = \bigcap_{i=1}^{\infty} A_i$.
От свойство (ii) знаем, че ако допълнението $B^c \in \mathcal{A}$, то и самото множество $B = (B^c)^c \in \mathcal{A}$. Затова ще изследваме допълнението на $B$. Използвайки законите на Де Морган, обръщаме сечението в обединение:
\[
B^c = \left( \bigcap_{i=1}^{\infty} A_i \right)^c = \bigcup_{i=1}^{\infty} A_i^c
\]
От свойство (ii) на $\sigma$-алгебрата, тъй като $A_i \in \mathcal{A}$, следва че и $A_i^c \in \mathcal{A}$.
Тогава, съгласно свойство (iii), изброимото обединение на множества, които принадлежат на $\mathcal{A}$, също принадлежи на $\mathcal{A}$. Следователно:
\[
\bigcup_{i=1}^{\infty} \underbrace{A_i^c}_{\substack{\in \mathcal{A} \\ \text{от (ii)}}} \stackrel{\text{(iii)}}{\in} \mathcal{A}
\]
Така показахме, че $B^c \in \mathcal{A}$. Прилагайки отново свойство (ii), заключаваме, че $B = (B^c)^c \in \mathcal{A}$. С това следствието е доказано.
\end{proof}

\subsection{Примери за $\sigma$-алгебри}

\begin{example}[Тривиална $\sigma$-алгебра]\label{ex:01-10}
За експеримент с $\Omega = \{0, 1\}$ (хвърляне на монета), най-простата $\sigma$-алгебра е:
\[
\mathcal{A} = \{\emptyset, \Omega\}
\]
Тя отговаря на условията, но е безполезна за практически изчисления, тъй като не съдържа информация за самите изходи (0 или 1).
\end{example}

\begin{example}[Множество от всички подмножества: Булеан]\label{ex:01-11}
За същото $\Omega = \{0, 1\}$, естествената $\sigma$-алгебра е колекцията от всички възможни подмножества:
\[
\mathcal{A} = 2^\Omega = \{\emptyset, \Omega, \{0\}, \{1\}\}
\]
Ако $\Omega$ има $n$ елемента, то $\mathcal{A} = 2^\Omega$ има $2^n$ елемента. За дискретни, крайни или изброимо безкрайни пространства, обикновено работим точно с тази най-богата $\sigma$-алгебра, защото върху нея винаги може коректно да се дефинира вероятност.
\end{example}

\begin{example}[Непрекъснати пространства и Борелова $\sigma$-алгебра]\label{ex:01-12}
Когато разглеждаме неизброими пространства, например $\Omega = \mathbb{R}_+ = [0, \infty)$, колекцията от всички подмножества $2^\Omega$ става „твърде голяма“. Математически е невъзможно да се зададе вероятностна мярка, която да измерва \emph{всяко едно} подмножество на реалните числа. За щастие, в практиката ние никога не се интересуваме от екзотични и „изкривени“ множества. Вълнуваме се от интервали от вида $A = (a, b)$ — например дали температурата е между 10 и 12 градуса. $\sigma$-алгебрата, породена от всички такива интервали, се нарича \emph{Борелова $\sigma$-алгебра} и върху нея може коректно да се зададе вероятност.
\end{example}

\begin{example}[Редуциране на неизброимо пространство]\label{ex:01-13}
Представете си, че стреляме с лък по мишена за дартс. Точната точка на попадение е от неизброимото множество на двумерната равнина. Ние обаче се интересуваме само от това колко точки ще получим, т.е. в кой цветен сектор е попаднала стрелата (например: 1 — черно, 2 — синьо, 3 — червено). Така ние свеждаме сложния неизброим експеримент до просто пространство с три елементарни изхода:
\[
\Omega = \{\omega_1, \omega_2, \omega_3\}
\]
Съответната $\sigma$-алгебра ще бъде $2^\Omega$ и ще съдържа точно $2^3 = 8$ елемента:
\[
\mathcal{A} = 2^\Omega = \{\emptyset, \Omega, \{\omega_1\}, \{\omega_2\}, \{\omega_3\}, \{\omega_1, \omega_2\}, \{\omega_1, \omega_3\}, \{\omega_2, \omega_3\}\}
\]
Това улеснява значително анализа и елиминира нуждата от сложен математически апарат там, където практическият проблем не го изисква.
\end{example}

\section{Математическо моделиране на проблема с тотализатора}

Нека се върнем на медийния скандал с тотализатора „6 от 49“. Искаме да моделираме събитието: \emph{„Паднали са се едни и същи числа в два последователни тиража в рамките на 10 000 изиграни тиража.“}

Пространството от елементарни събития за всички 10 000 тиража заедно може да се опише като Декартово произведение на индивидуалните пространства на всеки тираж:
\[
\Omega = \bigotimes_{i=1}^{10000} \Omega_i
\]
Едно конкретно елементарно събитие (история на всички тегления) изглежда така:
\[
\omega = (\omega^{(1)}, \omega^{(2)}, \dots, \omega^{(10000)})
\]
където всяко $\omega^{(i)}$ е шесторка числа от $i$-тия тираж.
Мощността на това пространство е колосална:
\[
|\Omega| = (\underbrace{13983816}_{N})^{10000} = N^{10000}
\]
Съответната $\sigma$-алгебра $\mathcal{A} = 2^\Omega$ ще има $2^{N^{10000}}$ елемента, което е необозримо голямо число, но важното е, че е \emph{крайно}, следователно нямаме теоретични проблеми с измерването.

За да пресметнем вероятността на разглежданото скандално събитие $A$, ние го разбиваме на по-прости, елементарни компоненти. Дефинираме $A_i$ като събитието: \emph{„шесторката в $i$-тия тираж е същата като шесторката в $(i+1)$-вия тираж“}.
Формално:
\[
A_i = \{ \omega^{(i)} = \omega^{(i+1)} \} \subseteq \Omega, \quad A_i \in 2^\Omega
\]
Тогава търсеното събитие $A$ („едни и същи числа в два последователни тиража от всички 10 000“) е просто обединението на всички възможни такива съвпадения за $i$ от 1 до 9999:
\[
A = \bigcup_{i=1}^{9999} A_i
\]
Превеждайки словесното описание на събитието в стриктна теоретико-множествена форма, ние създаваме солидна основа за пресмятане на неговата вероятност. В следващите лекции ще видим как да оценим вероятността на подобно обединение. Макар тя да не е голяма, оценката показва, че не е толкова малка, че сама по себе си да предизвика наистина сериозно съмнение.

\section{Задачи}

\begin{enumerate}
\item Проверете закона на Де Морган за две множества,
\[
(A \cup B)^c = A^c \cap B^c,
\]
а по желание проверете и другия закон на Де Морган.

\item \textbf{Парадоксът с двата плика.}

В края на лекцията оставяме една класическа, интригуваща задача.
Представете си, че имате два затворени плика. Във всеки плик има лист с написано реално число, съответно $a$ и $b$, като $a < b$ (двете числа са различни). Организаторът на играта знае числата, но вие нямате абсолютно никаква предварителна (априорна) информация в какъв диапазон са те (може да са 1 и 2, може да са 1 и 10 000, може да са отрицателни).
Избирате случайно един от пликовете и го отваряте. Виждате числото в него. Целта ви е да познаете дали това е по-голямото число, или по-малкото (или респективно да решите дали да смените плика, за да вземете по-голямата сума).

\textbf{Въпросът е:} Съществува ли стратегия, която да ви гарантира вероятност строго по-голяма от $\frac{1}{2}$ (тоест по-добра от чисто налучкване), че ще познаете/изберете по-голямото число?
Помислете върху този проблем, без да спекулирате произволно, а търсейки строга математическа логика.

\begin{supp}[Решение на задачата с двата плика]\label{supp:envelopes}
Отговорът е \emph{да}, и ключът е стратегията да бъде \emph{случайна}.

Нека $C$ е събитието „сменям плика“, а $A$ и $B$ — съответно събитията „видял съм
по-малкото число $a$“ и „видял съм по-голямото число $b$“, всяко с вероятност $\tfrac12$.
Печелим, ако при $A$ сменим, а при $B$ не сменим, тоест вероятността за успех е
\[
\mathbb{P}(\text{успех})
 = \tfrac12\,\mathbb{P}(C \given A) + \tfrac12\big(1 - \mathbb{P}(C \given B)\big)
 = \tfrac12 + \tfrac12\big(\mathbb{P}(C \given A) - \mathbb{P}(C \given B)\big).
\]
Следователно всяка стратегия, при която сменяме \emph{по-охотно} след малко число,
отколкото след голямо — тоест $\mathbb{P}(C \given A) > \mathbb{P}(C \given B)$ — печели с
вероятност строго над $\tfrac12$.

Такава стратегия съществува, при това без никакво предварително знание за $a$ и $b$:
след като видим число $x$, сменяме с вероятност $e^{-x}$ (или с която и да е строго
намаляваща функция със стойности в $(0,1)$). Тогава
\[
\mathbb{P}(C \given A) - \mathbb{P}(C \given B) = e^{-a} - e^{-b} > 0,
\]
понеже $a < b$. Печалбата е малка, ако $a$ и $b$ са близки, но е строго положителна винаги.
\end{supp}
\end{enumerate}


\chapter{Аксиоми на вероятността. Дискретна и геометрична вероятност}
\section{Въведение. Дефиниция за вероятност}

След като въведохме пространството от елементарни събития и $\sigma$-алгебрата, разполагаме с всичко необходимо, за да дефинираме строго що е то вероятност. Нека имаме едно множество от елементарни събития $\Omega$ — всичките възможни изходи на някакъв експеримент — и към него една $\sigma$-алгебра $\mathcal{A}$ (колекция от подмножества на $\Omega$). Най-често, ако говорим за дискретни вероятности, това може да са всички подмножества на $\Omega$.

\begin{keydefn}[вероятностна функция]\label{def:probability}
Функцията $\mathbb{P}: \mathcal{A} \to [0, 1]$ се нарича \emph{вероятностна функция} (или просто \emph{вероятност}), ако са изпълнени следните три условия (аксиоми):
\begin{enumerate}
    \item[(A1)] \textbf{Вероятност на достоверното събитие:}
    \[ \mathbb{P}(\Omega) = 1 \]

    \item[(A2)] \textbf{Вероятност на противоположното събитие:} ако $A \in \mathcal{A}$, то $A^c \in \mathcal{A}$ и
    \[ \mathbb{P}(A^c) = 1 - \mathbb{P}(A) \]

    \item[(A3)] \textbf{Изброима адитивност:} ако $A_1, A_2, \dots \in \mathcal{A}$ са две по две несечещи се (т.е. $A_i \cap A_j = \emptyset$ за $i \neq j$), то
    \[ \mathbb{P}\left(\bigcup_{i=1}^{\infty} A_i\right) = \sum_{i=1}^{\infty} \mathbb{P}(A_i) \]
\end{enumerate}
\end{keydefn}

Аксиома (A1) по съвсем естествен начин ни казва, че вероятността да се падне някое от всички възможни елементарни събития (т.е. експериментът да завърши с някакъв резултат) е точно $1$.

В този смисъл вероятността отговаря на класическото разбиране, което имаме за мярка. Ако искаме да пресметнем общото тегло на всички нас, можем да сумираме теглата на всеки от нас индивидуално, защото телата ни не се пресичат. Аналогично, ако искаме да пресметнем лицето на част от пода, можем да сумираме плочките, които попадат в тази част, стига те да не се припокриват. Грубо казано, вероятността на изброимо много несечещи се събития е сумата от индивидуалните вероятности. Това е вероятността чисто като мярка — тя удовлетворява тези свойства, които я правят вероятност.

Преди да преминем към следствията, нека въведем една операция над множества. Ако $A$ и $B$ са две събития от $\mathcal{A}$, разликата $A \setminus B$ разбираме като:
\[ A \setminus B = A \cap B^c \]
Това са всички елементарни събития, които попадат в $A$, но не попадат в $B$. Геометрично, ако начертаем множествата $A$ и $B$, $A \setminus B$ е сечението на $A$ с допълнението на $B$.

\section{Следствия от дефиницията за вероятност}

Нека имаме $\sigma$-алгебра $\mathcal{A}$ и вероятност $\mathbb{P}$. В сила са следните твърдения:

\begin{prop}[Вероятност на празното множество]\label{prop:empty}
\[ \mathbb{P}(\emptyset) = 0 \]
\end{prop}
\begin{proof}
Празното множество е допълнението на пълното множество $\Omega$ (т.е. $\emptyset = \Omega^c$). Използвайки второто свойство:
\[ \mathbb{P}(\emptyset) = \mathbb{P}(\Omega^c) = 1 - \mathbb{P}(\Omega) = 1 - 1 = 0 \]
Това е интуитивно ясно — вероятността да не се случи нищо е нула, тъй като все нещо трябва да се случи при експеримента.
\end{proof}

\begin{prop}[Разбиване на събитие]\label{prop:split}
За всеки $A, B \in \mathcal{A}$:
\[ \mathbb{P}(B) = \mathbb{P}(A \cap B) + \mathbb{P}(A^c \cap B) \]
\end{prop}
\begin{proof}
Всяко множество $B$ може да се представи като обединение на две несечещи се множества: частта на $B$, която попада в $A$, и частта на $B$, която е извън $A$. Тоест:
\[ B = (A \cap B) \cup (A^c \cap B) \]
Тъй като $(A \cap B)$ и $(A^c \cap B)$ не се пресичат, можем да приложим аксиомата за събираемост, получавайки точно търсената формула.
\end{proof}

\begin{prop}[Монотонност на вероятността]\label{prop:monotone}
Ако $A \subseteq B$, то $\mathbb{P}(A) \le \mathbb{P}(B)$.
\end{prop}
\begin{proof}
Щом $A \subseteq B$, можем да представим $B$ като:
\[ B = A \cup (A^c \cap B) \]
Като приложим твърдение~\ref{prop:split} и вземем предвид, че сечението $A \cap B = A$ (защото $A$ е подмножество на $B$), получаваме:
\[ \mathbb{P}(B) = \mathbb{P}(A) + \mathbb{P}(A^c \cap B) \]
Тъй като $\mathbb{P}(A^c \cap B) \ge 0$, следва, че $\mathbb{P}(B) \ge \mathbb{P}(A)$.
\end{proof}

\begin{cor}[Вероятност на обединение на две събития]\label{cor:union2}
За всеки $A, B \in \mathcal{A}$:
\[ \mathbb{P}(A \cup B) = \mathbb{P}(A) + \mathbb{P}(B) - \mathbb{P}(A \cap B) \]
\end{cor}
\begin{proof}
Тъй като $A \cup B = A \cup (A^c \cap B)$, където $A$ и $A^c \cap B$ не се пресичат, от адитивността получаваме:
\[ \mathbb{P}(A \cup B) = \mathbb{P}(A) + \mathbb{P}(A^c \cap B) \]
Замествайки $\mathbb{P}(A^c \cap B) = \mathbb{P}(B) - \mathbb{P}(A \cap B)$ от твърдение~\ref{prop:split}, получаваме търсеното равенство.
\end{proof}

Ако събитията са взаимно изключващи се ($A \cap B = \emptyset$), членът $\mathbb{P}(A \cap B) = 0$ и вероятността на обединението е просто сумата им.

\begin{prop}[Непрекъснатост в нулата]\label{prop:cont0}
Ако имаме намаляваща редица от събития $A_1 \supseteq A_2 \supseteq \dots \supseteq A_n \supseteq \dots$, такива че тяхното сечение е празното множество ($\bigcap_{j=1}^\infty A_j = \emptyset$), то границата на техните вероятности е нула:
\[ \lim_{j \to \infty} \mathbb{P}(A_j) = 0 \]
\end{prop}
\begin{proof}
Идеята е да изразим множеството $A_j$ като изброимо обединение на несечещи се „пръстени“ (венци). Можем да представим $A_j$ като:
\[ A_j = (A_j \setminus A_{j+1}) \cup (A_{j+1} \setminus A_{j+2}) \cup \dots = \bigcup_{n=j}^{\infty} (A_n \setminus A_{n+1}) \]
Тези разлики $A_n \setminus A_{n+1}$ са взаимно несечещи се по построение. Поради адитивността на вероятността имаме:
\[ \mathbb{P}(A_j) = \sum_{n=j}^{\infty} \mathbb{P}(A_n \setminus A_{n+1}) \]
За $j=1$ получаваме:
\[ \mathbb{P}(A_1) = \sum_{n=1}^{\infty} \mathbb{P}(A_n \setminus A_{n+1}) \le 1 \]
Тъй като редът с общ член $\mathbb{P}(A_n \setminus A_{n+1})$ е сходящ (сумата му е ограничена от $1$), то остатъчната сума на този сходящ ред трябва да клони към нула, когато $j$ клони към безкрайност. Именно тази остатъчна сума е равна на $\mathbb{P}(A_j)$. Следователно $\lim_{j \to \infty} \mathbb{P}(A_j) = 0$. Този резултат е много важен и се нарича „непрекъснатост в нулата“.
\end{proof}

\begin{prop}[Полуадитивност (субадитивност)]\label{prop:subadd}
Ако $A_1, A_2, \dots$ са произволни събития (не задължително несечещи се), то вероятността на тяхното обединение е по-малка или равна на сумата от техните вероятности:
\[ \mathbb{P}\left(\bigcup_{j=1}^{\infty} A_j\right) \le \sum_{j=1}^{\infty} \mathbb{P}(A_j) \]
\end{prop}
\begin{proof}
Идеята е да конструираме нова редица от непресичащи се събития $C_1, C_2, \dots$, които покриват същото обединение. Полагаме $C_1 = A_1$, $C_2 = A_2 \setminus A_1$, $C_3 = A_3 \setminus (A_1 \cup A_2)$ и общо:
\[ C_n = A_n \setminus \bigcup_{i=1}^{n-1} A_i \]
Тези $C_j$ са непресичащи се и $\bigcup C_j = \bigcup A_j$. Тъй като $C_j \subseteq A_j$, от монотонността (твърдение~\ref{prop:monotone}) имаме $\mathbb{P}(C_j) \le \mathbb{P}(A_j)$. Тогава:
\[ \mathbb{P}\left(\bigcup A_j\right) = \mathbb{P}\left(\bigcup C_j\right) = \sum \mathbb{P}(C_j) \le \sum \mathbb{P}(A_j) \]
\end{proof}

\section{Дискретна вероятност}

Нека сега разгледаме как конкретно се задава вероятност, като започнем с дискретната.
Най-простият прототип е експеримент с два възможни изхода, например 0 и 1 (ези/тура, успех/неуспех).
Тук $\Omega = \{0, 1\}$. Дефинираме вероятността да се падне $1$ като $\mathbb{P}(\{1\}) = p$ (за някакво число $0 \le p \le 1$), а вероятността да се падне $0$ е $\mathbb{P}(\{0\}) = 1 - p$. Често бележим $q = 1 - p$.

\subsection{Краен брой елементарни събития}

\begin{defn}[дискретна вероятност върху крайно пространство]\label{def:discrete}
Нека $\Omega = \{1, 2, \dots, N\}$ и нека са дадени $N$ на брой неотрицателни числа $p_1, p_2, \dots, p_N$, такива че
\[ \sum_{i=1}^N p_i = 1. \]
За всяко събитие $A \subseteq \Omega$ дефинираме
\[ \mathbb{P}(A) := \sum_{i \in A} p_i . \]
\end{defn}

Това кодира всеки експеримент с $N$ възможни изхода; в частност вероятността на единичен изход е $\mathbb{P}(\{i\}) = p_i$.
Тази функция реално е вероятност. Сумата $\sum_{i \in A} p_i$ винаги е между 0 и 1. Също така, аксиомата за противоположното събитие е изпълнена:
\[ \mathbb{P}(A^c) = \sum_{i \in A^c} p_i = 1 - \sum_{i \in A} p_i = 1 - \mathbb{P}(A) \]

\subsection{Равномерна вероятност}
Равномерната вероятност е частен случай, при който всички $p_i$ са равни:
\[ p_i = \frac{1}{N}, \quad 1 \le i \le N \]
Този модел се използва за честни игри, например при хвърляне на зар или рулетка, или в байесовата статистика като неинформативно априорно разпределение, когато нямаме допълнителна информация за това кой изход е по-вероятен.

\begin{example}[Тото 6/49]\label{ex:02-1}
Продължаваме \crosslecture{ex:01-3}{модела на един тираж от предходната лекция}. Броят на възможните шестици е $N = 13983816$. Ако играта е честна, вероятността да се падне коя да е конкретна шестица (например $\{1, 2, 3, 4, 5, 6\}$) е точно същата като всяка друга:
\[ \mathbb{P}(\{\text{конкретна шестица}\}) = \frac{1}{13983816} \]
Много хора играят с числа от 1 до 31 (поради дати на раждане и т.н.). Вероятността тези числа да се паднат не е по-голяма, но ако се паднат „по-големи“ числа, има по-малко участници, с които ще се раздели печалбата. Самата вероятност за печалба обаче е напълно фиксирана и униформна.

\end{example}

\subsection{Изброимо множество от елементарни събития}
Често е удобно пространството $\Omega$ да бъде безкрайно, но изброимо, например $\Omega = \{1, 2, 3, \dots\}$ или $\{0, 1, 2, \dots\}$. Това се прилага, когато моделираме брой събития, време до повреда (дискретизирано) и др. Не е нужно да слагаме изкуствена горна граница (например 10 000 дни за процесор), тъй като математически е по-удобно да допуснем потенциална безкрайност.

Тук имаме редица от вероятности $p_i \ge 0$, за които:
\[ \sum_{i=1}^{\infty} p_i = 1 \]
Важен факт: \textbf{Не можем да дефинираме равномерна вероятност върху изброимо безкрайно множество.} Ако всички $p_i$ са равни на някакво $p > 0$, сумата им ще бъде безкрайност. Ако $p = 0$, сумата им ще е $0$. И в двата случая сумата не е $1$.

\begin{example}[Разпределение с „тежки опашки“]\label{ex:02-2}
Можем да изберем:
\[ p_i = \frac{6}{\pi^2} \frac{1}{i^2}, \quad \text{за } i \ge 1 \]
Това е валидно вероятностно разпределение, защото от математическия анализ знаем, че $\sum_{i=1}^\infty \frac{1}{i^2} = \frac{\pi^2}{6}$, така че константата $\frac{6}{\pi^2}$ нормализира сумата до 1. Такива модели се ползват за екстремни събития, които макар и редки, имат значим ефект (тежки опашки).

\end{example}

\begin{example}[Геометрично разпределение]\label{ex:02-3}
Моделираме броя неуспехи преди първия успех. Нека $p \in (0, 1)$ и $q = 1 - p$. Тогава за $\Omega = \{0, 1, 2, 3, \dots\}$ полагаме:
\[ p_i = q^i p, \quad i \ge 0 \]
Това също сумира до 1 благодарение на формулата за сума на безкрайна геометрична прогресия:
\[ \sum_{i=0}^\infty q^i p = p \sum_{i=0}^\infty q^i = p \frac{1}{1-q} = \frac{p}{p} = 1 \]
Геометричното разпределение е интересно със свойството си за „безпаметност“ — много подходящо за моделиране на живот на машинни елементи (но не и на хора, защото ние стареем).

\end{example}

\section{Геометрична вероятност}

Освен дискретното пространство, можем да имаме и непрекъснато.

\begin{defn}[геометрична вероятност]\label{def:geometric}
Нека $\Omega \subset \mathbb{R}^d$ е област в равнината (или пространството) с краен обем (или площ):
\[ \infty > \operatorname{Vol}(\Omega) = |\Omega| = \int_{\Omega} dx . \]
\emph{Геометрична вероятност} наричаме равномерната вероятност върху това непрекъснато множество: за всяко подмножество $A \subseteq \Omega$
\[ \mathbb{P}(A) := \frac{|A|}{|\Omega|} . \]
\end{defn}

Тук множеството от елементарни събития има неизброимо много точки. Ние не ползваме кардиналност (брой точки), а понятието обем (или площ). Всяка конкретна точка $x \in \Omega$ има обем $0$, следователно вероятността да се падне точно дадена конкретна точка е:
\[ \mathbb{P}(\{x\}) = \frac{0}{|\Omega|} = 0 \]
\phantomsection\label{ex:monte-carlo-volume}Това е фундаментално — при непрекъснатите разпределения вероятността за една конкретна точка е нула, въпреки че точката е възможен изход. Измерват се вероятности за попадане в области с положителен обем. Това свойство се използва например при методите Монте Карло, където генерираме случайни точки и оценяваме площи чрез съотношението на броя попаднали в $A$ точки към общия брой точки.

\section{Вероятностно пространство}

Дотук въведохме поотделно трите обекта, които са ни нужни. Остава да ги съберем в едно
понятие.

\begin{keydefn}[вероятностно пространство]\label{def:prob-space}
Нека $\Omega$ е пространство от елементарни събития, $\mathcal{A}$ е $\sigma$-алгебра от
подмножества на $\Omega$, а $\mathbb{P}$ е вероятностна функция върху $\mathcal{A}$ по
дефиниция~\ref{def:probability}. Тогава тройката
\[ (\Omega, \mathcal{A}, \mathbb{P}) \]
се нарича \emph{вероятностно пространство}.
\end{keydefn}

С това разполагаме с всички елементи, които ни позволяват да дефинираме вероятността
строго: множеството от изходи, колекцията от събития, на които можем да придаваме
вероятност, и самата вероятност. Най-простият пример е крайното равномерно пространство:
$\Omega = \{1, \dots, N\}$, $\mathcal{A} = 2^\Omega$ и $\mathbb{P}(\{i\}) = 1/N$ за всяко $i$.

Няма да разработваме това понятие в дълбочина, но е важно да се знае с какъв обект се
работи: всяко твърдение по-нататък мълчаливо предполага фиксирано вероятностно
пространство.

\section{Условна вероятност}

Накрая ще въведем едно много важно понятие, което отчита как се променят вероятностите при постъпване на нова информация. Да предположим, че имаме вероятностно пространство $(\Omega, \mathcal{A}, \mathbb{P})$. Наблюдаваме експеримента и научаваме, че се е случило събитието $A$.

Нека $A, B \in \mathcal{A}$ и $\mathbb{P}(A) > 0$. \emph{Условна вероятност} на $B$ при условие $A$ наричаме
\[ \mathbb{P}(B \given A) = \frac{\mathbb{P}(A \cap B)}{\mathbb{P}(A)} . \]
Формалната дефиниция и подробното ѝ обсъждане следват в началото на
\crosslecture{def:cond3}{лекция 3}; тук само въвеждаме означението.

Интуитивно, тъй като вече знаем, че сме попаднали в $A$, цялото ни пространство от възможни изходи се стеснява от $\Omega$ до $A$. Всички елементарни събития извън $A$ стават невъзможни и биват „изхвърлени“. От събитието $B$ оцелява само тази част, която се пресича с $A$ (т.е. $A \cap B$).
За да получим коректна нова вероятност, която да сумира до 1, трябва да разделим (нормализираме) на вероятността на новото „пълно“ пространство, което е $\mathbb{P}(A)$.

Например, в геометричната вероятност, ако първоначалната вероятност е била $\frac{|B|}{|\Omega|}$, след като знаем, че сме в $A$, новата вероятност става:
\[ \mathbb{P}(B \given A) = \frac{|A \cap B| / |\Omega|}{|A| / |\Omega|} = \frac{|A \cap B|}{|A|} \]
Това преизчисляване е в основата на това как ние променяме своите очаквания въз основа на допълнителна информация. В следващата лекция ще разгледаме конкретни примери и приложения на условната вероятност.

\section{Задачи}

\begin{enumerate}
    \item За дискретната вероятност от дефиниция~\ref{def:discrete} проверете аксиомата за адитивност: ако $A_1,\dots,A_k \subseteq \Omega$ са две по две непресичащи се, докажете, че
    \[
    \mathbb{P}\left(\bigcup_{j=1}^k A_j\right)
    = \sum_{j=1}^k \mathbb{P}(A_j).
    \]
    \item Нека равномерният експеримент има $N$ изхода и $A$ е събитието, че се е паднал нечетен номер. Пресметнете $\mathbb{P}(A)$, когато $N$ е нечетно.
\end{enumerate}


\chapter{Условна вероятност, независимост и формула на Бейс}
\section{Условна вероятност}

\subsection{Въведение и интуиция}

Когато наблюдаваме някакъв случаен експеримент, ние работим с предварително зададени вероятности. Но какво се случва, когато получим частична информация за резултата от експеримента? Нашите очаквания, и съответно вероятностите на събитията, се променят. За да формализираме това, въвеждаме понятието за условна вероятност.

Да разгледаме един съвсем прост и класически пример – хвърляне на правилно (симетрично) зарче. Множеството от възможни изходи е $\Omega = \{1, 2, 3, 4, 5, 6\}$. При липса на каквато и да е друга информация, вероятността да се падне което и да е число $i \in \Omega$ е точно една шеста:
\[
\mathbb{P}(\{i\}) = \frac{1}{6}.
\]

Представете си сега, че имате допълнителната информация, че се е паднало нечетно число. Нека обозначим събитието „паднало се е нечетно число“ с $A = \{1, 3, 5\}$. Тогава вероятностите за отделните изходи драстично се променят. Ясно е, че ако числото е нечетно, то не може да бъде $2$. Следователно, вероятността да се е паднало $2$ при положение, че $A$ се е случило, е $0$. От друга страна, възможностите остават само три ($1, 3$ и $5$), и тъй като зарчето е симетрично, те са равновероятни. Следователно, вероятността да се е паднало $1$ при условие, че е паднало нечетно число, е вече $1/3$. 

Формално това може да се пресметне чрез отношението:
\[
\mathbb{P}(\{i\} \mid A) = \frac{\mathbb{P}(\{i\} \cap A)}{\mathbb{P}(A)}
\]
Например, ако $i = 1$, то $\{1\} \cap \{1, 3, 5\} = \{1\}$. Тогава $\mathbb{P}(\{1\} \cap A) = \frac{1}{6}$, а $\mathbb{P}(A) = \frac{3}{6} = \frac{1}{2}$. Резултатът е:
\[
\mathbb{P}(\{1\} \mid A) = \frac{1/6}{1/2} = \frac{1}{3}.
\]
Ако пък $i = 2$ (което е четно число), сечението $\{2\} \cap \{1, 3, 5\}$ е празното множество $\emptyset$, което има вероятност $0$. В този случай условната вероятност е $\frac{0}{1/2} = 0$. Този прост пример нагледно илюстрира как появата на нова информация (че се е случило събитието $A$) преизчислява нашите вероятности.

\subsection{Формална дефиниция}

\begin{defn}[условна вероятност]\label{def:cond3}
Нека е дадено вероятностно пространство $(\Omega, \mathcal{F}, \mathbb{P})$ и събитие $A \in \mathcal{F}$ такова, че $\mathbb{P}(A) > 0$. \emph{Условна вероятност на събитието $B$ при условие $A$} наричаме
\[
\mathbb{P}(B \mid A) = \frac{\mathbb{P}(A \cap B)}{\mathbb{P}(A)} .
\]
\end{defn}

Изискването $\mathbb{P}(A) > 0$ е необходимо, за да можем да извършим делението. Ако събитието $A$ има нулева вероятност (мярка $0$), в определени случаи (като при формулата за пълната вероятност) за удобство можем да приемем условната вероятност за нула, за да не усложняваме записите, макар че строго погледнато тя не е дефинирана.

Важно е да се отбележи, че функцията $B \mapsto \mathbb{P}(B \mid A)$ сама по себе си задава нова вероятностна мярка върху същото измеримо пространство $(\Omega, \mathcal{F})$. Тя притежава всички свойства на една класическа вероятност.

\subsection{Пример с Тото 6 от 49}

За да видим как частичната информация променя вероятностите в един по-реалистичен случай, да разгледаме играта Тото 6 от 49. Имате пуснат фиш с една конкретна шесторка, например $\{1, 2, 3, 4, 5, 6\}$. Нека обозначим нашата шесторка с $W_1$.
Вероятността тя да бъде изтеглена първоначално е:
\[
\mathbb{P}(W_1) = \frac{1}{\binom{49}{6}} = \frac{1}{13\,983\,816} \approx \frac{1}{N}
\]
където с $N$ сме обозначили общия брой възможни шесторки.

Представете си, че докато сте слушали тегленето по радиото, сте чули само част от числата и знаете, че в печелившата шесторка със сигурност присъстват числата $1$ и $2$. Нека това събитие да го наречем $A$. Търсим новата (условна) вероятност вашата шесторка $W_1$ да е печеливша: $\mathbb{P}(W_1 \mid A)$.

Тъй като вашата шесторка $W_1 = \{1, 2, 3, 4, 5, 6\}$ съдържа числата $1$ и $2$, то събитието $W_1$ е подмножество на събитието $A$. Тогава сечението $W_1 \cap A$ е просто $W_1$.
По дефиницията за условна вероятност:
\[
\mathbb{P}(W_1 \mid A) = \frac{\mathbb{P}(W_1 \cap A)}{\mathbb{P}(A)} = \frac{\mathbb{P}(W_1)}{\mathbb{P}(A)}
\]
Вероятността на събитието $A$ (да се изтеглят $1$ и $2$) може да се намери чрез комбинаторика. Ние вече сме фиксирали двете числа (1 и 2), така че остава да изберем останалите 4 числа от възможните 47. Броят на шесторките, съдържащи $1$ и $2$, е $\binom{47}{4}$.
Следователно $\mathbb{P}(A) = \frac{\binom{47}{4}}{\binom{49}{6}}$. Като заместим в условната вероятност:
\[
\mathbb{P}(W_1 \mid A) = \frac{1 / \binom{49}{6}}{\binom{47}{4} / \binom{49}{6}} = \frac{1}{\binom{47}{4}}
\]
Стойността $\binom{47}{4}$ е около $178\,000$. Виждаме как при наличието само на тази малка частична информация вероятността за печалба се увеличава драстично – от порядъка на $1$ към $14$ милиона на $1$ към $178$ хиляди.

\subsection{Прост модел за времето (Марковски вериги)}

Условните вероятности са в основата на стохастични процеси като Марковските вериги, които моделират системи, чието бъдещо състояние зависи само от настоящето.
Нека разгледаме един изключително опростен модел за прогноза на времето, при който дните могат да бъдат само два вида: „дъжд“ (0) и „слънце“ (1).
Можем да зададем модела чрез условни вероятности, или така наречените преходни вероятности (стохастична матрица):
\begin{align*}
\mathbb{P}(\text{утре е дъжд} \mid \text{днес е дъжд}) &= p_1 \\
\mathbb{P}(\text{утре е слънце} \mid \text{днес е дъжд}) &= 1 - p_1 \\
\mathbb{P}(\text{утре е слънце} \mid \text{днес е слънце}) &= p_2 \\
\mathbb{P}(\text{утре е дъжд} \mid \text{днес е слънце}) &= 1 - p_2
\end{align*}
Тук вероятностите за времето утре се преизчисляват при условие какво е времето днес. Ако пространството на елементарните събития за два последователни дни се зададе като декартово произведение $\{0, 1\} \times \{0, 1\}$, то $p_1$ е всъщност условната вероятност $\mathbb{P}(\text{утре е 0} \mid \text{днес е 0}) = \frac{\mathbb{P}(\{0, 0\})}{\mathbb{P}(\{0, 0\} \cup \{0, 1\})}$. Този пример илюстрира, че често на практика ние директно моделираме и задаваме условните вероятности за дадена система.

\subsection{Пример с гласоподаватели}

Друг интересен пример е свързан със социологически проучвания и избори.
Нека имаме общо гласоподаватели, от които:
\begin{itemize}
    \item $N_1$ гласуват за Партия 1.
    \item $N_2$ гласуват за Партия 2.
\end{itemize}
Нека измежду тях имаме определен брой млади гласоподаватели:
\begin{itemize}
    \item $m_1$ млади гласуват за Партия 1 ($m_1 \le N_1$).
    \item $m_2$ млади гласуват за Партия 2 ($m_2 \le N_2$).
\end{itemize}

Ако изберем случаен човек от \emph{цялата} популация, вероятността събитието $A_1$ („човекът гласува за Партия 1“) да се случи е:
\[
\mathbb{P}(A_1) = \frac{N_1}{N_1 + N_2}
\]
Ако обаче имаме информацията, че случайно избраният човек е млад (събитие $M$), каква е условната вероятност той да е гласувал за Партия 1?
Младите хора, гласуващи за Партия 1, са подмножество на всички млади хора. Следователно:
\[
\mathbb{P}(A_1 \mid M) = \frac{\mathbb{P}(A_1 \cap M)}{\mathbb{P}(M)} = \frac{\frac{m_1}{N_1 + N_2}}{\frac{m_1 + m_2}{N_1 + N_2}} = \frac{m_1}{m_1 + m_2}
\]
Забележете, че в крайния резултат $\frac{m_1}{m_1 + m_2}$ изобщо не участват $N_1$ и $N_2$. Тази пропорция може да бъде коренно различна от общата пропорция $N_1 / (N_1 + N_2)$. Този пример нагледно показва защо при статистически анализи трябва внимателно да се работи с условни вероятности, когато се ограничаваме до някакво подмножество (например определена възрастова група).

\section{Независимост на събития}

\subsection{Независимост на две събития}

Понятието за независимост е директно свързано с условната вероятност. 
\begin{defn}[независими събития]\label{def:independent}
Две събития $A$ и $B$ се наричат \emph{независими} (ще го бележим понякога с $A \perp B$), ако вероятността на тяхното сечение е равна на произведението от индивидуалните им вероятности:
\[
\mathbb{P}(A \cap B) = \mathbb{P}(A)\mathbb{P}(B) .
\]
\end{defn}

Защо това се нарича независимост? Ако приемем, че $\mathbb{P}(A) > 0$, можем да пресметнем условната вероятност на $B$ при условие $A$:
\[
\mathbb{P}(B \mid A) = \frac{\mathbb{P}(A \cap B)}{\mathbb{P}(A)} = \frac{\mathbb{P}(A)\mathbb{P}(B)}{\mathbb{P}(A)} = \mathbb{P}(B)
\]
Това означава, че информацията за настъпването на събитието $A$ не променя по никакъв начин първоначалната вероятност на събитието $B$. Събитието $A$ не носи никаква допълнителна информация за $B$. 

Ако $\mathbb{P}(A) = 0$, то сечението им автоматично има нулева вероятност, и равенството $\mathbb{P}(A \cap B) = \mathbb{P}(A)\mathbb{P}(B)$ е тривиално изпълнено ($0 = 0 \cdot \mathbb{P}(B)$).

\subsection{Независимост в съвкупност}

Често се налага да работим с повече от две събития, за които знаем (или допускаме в модела си), че са независими. Ако събитията са независими, вероятността на сеченията им се пресмята изключително лесно чрез произведения, което значително опростява изчисленията.

\begin{defn}[независимост в съвкупност]\label{def:mutual-independent}
Събитията $A_1, A_2, \dots, A_n$ се наричат \emph{независими в съвкупност}, ако за всяко подмножество от индекси $I \subseteq \{1, 2, \dots, n\}$ с поне два елемента е изпълнено:
\[
\mathbb{P}\left(\bigcap_{i \in I} A_i\right) = \prod_{i \in I} \mathbb{P}(A_i) .
\]
\end{defn}

Например, ако имаме три събития $A_1, A_2, A_3$, трябва да проверим четирите възможни подмножества от индекси с поне 2 елемента: $\{1, 2\}$, $\{1, 3\}$, $\{2, 3\}$ и $\{1, 2, 3\}$. Това означава, че трябва да са изпълнени едновременно следните условия:
\begin{align*}
\mathbb{P}(A_1 \cap A_2) &= \mathbb{P}(A_1)\mathbb{P}(A_2) \\
\mathbb{P}(A_1 \cap A_3) &= \mathbb{P}(A_1)\mathbb{P}(A_3) \\
\mathbb{P}(A_2 \cap A_3) &= \mathbb{P}(A_2)\mathbb{P}(A_3) \\
\mathbb{P}(A_1 \cap A_2 \cap A_3) &= \mathbb{P}(A_1)\mathbb{P}(A_2)\mathbb{P}(A_3)
\end{align*}

Полезна интуиция дава пресмятането на лице на правоъгълник: лицето е произведението от дължините на двете му страни $a$ и $b$. Тук страните са съответните „мерки“, а правоъгълникът е „сечението“. Тази хубава структура съществува само защото дължината и ширината на правоъгълника са ортогонални (независими) една спрямо друга.

\subsection{Приложение: Еднакви тиражи в Тотото}

Нека се върнем към един реален казус, възникнал в България през 2009 г., когато в два последователни тиража на Тото 6 от 49 се паднаха едни и същи числа. Хората тогава са се запитали: каква е вероятността за подобно нещо в историята на тотализатора?
Нека имаме 10 000 тиража. Търсим вероятността на събитието $A$: „паднали са се еднакви числа в поне два последователни тиража от общо 10 000 тиража“.
Можем да дефинираме $A_i$ като събитието „шесторката на $i$-тия тираж съвпада с шесторката на $(i+1)$-вия тираж“.
Тогава събитието $A$ е обединението на тези събития за $i$ от 1 до $9999$:
\[
A = \bigcup_{i=1}^{9999} A_i
\]
Изчисляването на вероятност на голямо обединение е трудно и би изисквало сложния принцип на включване и изключване. Много по-лесно е, когато работим с независими събития и техните сечения. Ето защо преминаваме към допълнението на $A$, използвайки формулите на Де Морган:
\[
\mathbb{P}(A) = 1 - \mathbb{P}(A^c) = 1 - \mathbb{P}\left(\left(\bigcup_{i=1}^{9999} A_i\right)^c\right) = 1 - \mathbb{P}\left(\bigcap_{i=1}^{9999} A_i^c\right)
\]
Събитията $A_i^c$ означават, че $i$-тият и $(i+1)$-вият тираж \emph{не} съвпадат. Може да се докаже, че ако събитията $A_i$ са независими в съвкупност, то и техните допълнения $A_i^c$ също са независими в съвкупност. Тогава вероятността на сечението се разпада на произведение:
\[
\mathbb{P}(A) = 1 - \prod_{i=1}^{9999} \mathbb{P}(A_i^c)
\]
Тъй като правилата на играта не се променят с времето, вероятностите $\mathbb{P}(A_i^c)$ са едни и същи за всички $1 \le i \le 9999$. Тоест $\mathbb{P}(A_i^c) = \mathbb{P}(A_1^c)$. Тогава получаваме:
\[
\mathbb{P}(A) = 1 - (\mathbb{P}(A_1^c))^{9999}
\]
Тъй като $\mathbb{P}(A_1)$ е $1 / N$ (където $N \approx 14$ милиона), то $\mathbb{P}(A_1^c) = 1 - 1/N$. Така формулата става:
\[
\mathbb{P}(A) = 1 - \left(1 - \frac{1}{N}\right)^{9999}
\]
Това приближение се пресмята изключително лесно и дори на ум може да се докаже (използвайки известни приближения от анализа като $(1-x)^n \approx 1-nx$ за малки $x$), че вероятността за това „чудо“ не е чак толкова колосално малка, колкото изглежда на пръв поглед, и истерията не винаги е математически обоснована.

\section{Пълна група от събития и Формула за пълната вероятност}

\subsection{Пълна група от събития (Хипотези)}

Когато разглеждаме някакво множество от елементарни събития $\Omega$, често е много естествено то да се раздели на групи въз основа на определена характеристика. Например, клиентите в магазин могат да се разделят на „жени“ и „мъже“. Гласоподавателите могат да се разделят на гласуващи за Партия 1, Партия 2, ..., Партия $N$ и негласуващи.

\begin{defn}[пълна група от събития]\label{def:partition}
Казваме, че събитията $H_1, H_2, \dots, H_n$ (или изброимо безкрайно много $H_1, H_2, \dots$) образуват \emph{пълна група от събития}, ако:
\begin{enumerate}
    \item Те са взаимно изключващи се (непресичащи се): $H_i \cap H_j = \emptyset$ за $i \neq j$.
    \item Те изчерпват цялото пространство: $\bigcup_{i} H_i = \Omega$.
\end{enumerate}
\end{defn}

Често тези събития $H_i$ се наричат \textbf{хипотези}, тъй като при наблюдаване на някакво следствие (данни), ние се питаме от коя точно хипотеза е породено то.

\subsection{Формула за пълната вероятност}

\begin{thm}[Формула за пълната вероятност]\label{thm:total-probability}
Нека е дадено вероятностно пространство $(\Omega, \mathcal{F}, \mathbb{P})$, нека $H_1, H_2, \dots, H_n$ е пълна група от събития и нека $A$ е някакво събитие. Тогава:
\[
\mathbb{P}(A) = \sum_{i=1}^n \mathbb{P}(A \cap H_i) = \sum_{i=1}^n \mathbb{P}(A \mid H_i)\mathbb{P}(H_i) .
\]
\end{thm}

Ако имаме изброимо безкрайно много хипотези, сумата просто продължава до безкрайност.
Ако случайно някое от събитията $H_i$ има вероятност нула, се възприема конвенцията, че съответното събираемо е нула (условната вероятност $\mathbb{P}(A \mid H_i)$ условно се приема за нула), за да не се усложнява сумата.

\begin{proof}
Тъй като $A = A \cap \Omega$, а $\Omega = \bigcup_{i=1}^n H_i$, по дистрибутивния закон имаме:
\[
A = A \cap \left( \bigcup_{i=1}^n H_i \right) = \bigcup_{i=1}^n (A \cap H_i)
\]
Тъй като $H_i$ са взаимно изключващи се, сеченията $A \cap H_i$ също са взаимно изключващи се. Следователно вероятността на тяхното обединение е сумата от вероятностите им:
\[
\mathbb{P}(A) = \sum_{i=1}^n \mathbb{P}(A \cap H_i)
\]
Като приложим дефиницията за условна вероятност $\mathbb{P}(A \cap H_i) = \mathbb{P}(A \mid H_i)\mathbb{P}(H_i)$, получаваме желаната формула.
\end{proof}

\textbf{Пример с болница:}
Представете си човек, който влиза в тежко състояние в болница. Каква е вероятността човекът да почине (събитие $A$)? Често тази обща статистика липсва или е трудна за намиране директно, но лесно може да се изведе. Нека $H_i$ са различните възможни тежки болести (хипотези). От медицинската статистика лекарите знаят честотата (вероятността) човек да се разболее от съответната болест $\mathbb{P}(H_i)$, както и смъртността при всяка болест $\mathbb{P}(A \mid H_i)$. Тогава общата вероятност се пресмята лесно чрез формулата за пълната вероятност, сумирайки по всички болести.

\section{Формула на Бейс}

Формулата на Бейс (от името на Томас Бейс) е едно от най-важните твърдения в теорията на вероятностите, стоящо в основата на Бейсовата статистика и машинното самообучение. Тя отговаря на въпроса: \emph{как да обновим вероятностите на нашите хипотези $H_k$, след като сме наблюдавали събитие (данни) $A$?}

\begin{thm}[Формула на Бейс]\label{thm:bayes}
Нека $H_1, H_2, \dots$ е пълна група от събития (хипотези) и нека $A$ е събитие с $\mathbb{P}(A) > 0$. Тогава вероятността на хипотезата $H_k$ при условие $A$ е:
\[
\mathbb{P}(H_k \mid A) = \frac{\mathbb{P}(A \mid H_k)\mathbb{P}(H_k)}{\sum_{i} \mathbb{P}(A \mid H_i)\mathbb{P}(H_i)} .
\]
\end{thm}

\begin{proof}
По дефиницията за условна вероятност:
\[
\mathbb{P}(H_k \mid A) = \frac{\mathbb{P}(H_k \cap A)}{\mathbb{P}(A)}
\]
Чрез условната вероятност числителят може да се запише като $\mathbb{P}(H_k \cap A) = \mathbb{P}(A \mid H_k)\mathbb{P}(H_k)$. Знаменателят $\mathbb{P}(A)$ се замества с разложението му от формулата за пълната вероятност. С това формулата е доказана.
\end{proof}

\subsection{Отношение между хипотези и обновяване на вероятностите}

Пресмятането на знаменателя във формулата на Бейс често пъти е компютърно много тежко (особено ако преминем към непрекъснати разпределения, където сумата става интеграл). Затова в практиката често се разглежда съотношението между две хипотези. 

Ако разгледаме две хипотези $H_k$ и $H_i$ при наличие на данни $D$, имаме:
\[
\frac{\mathbb{P}(H_k \mid D)}{\mathbb{P}(H_i \mid D)} = \frac{\frac{\mathbb{P}(D \mid H_k)\mathbb{P}(H_k)}{\mathbb{P}(D)}}{\frac{\mathbb{P}(D \mid H_i)\mathbb{P}(H_i)}{\mathbb{P}(D)}} = \frac{\mathbb{P}(D \mid H_k)}{\mathbb{P}(D \mid H_i)} \cdot \frac{\mathbb{P}(H_k)}{\mathbb{P}(H_i)}
\]
Това ни показва как апостериорното (обновеното) съотношение на хипотезите се получава чрез умножаване на априорното (първоначалното) им съотношение $\frac{\mathbb{P}(H_k)}{\mathbb{P}(H_i)}$ с фактора $\frac{\mathbb{P}(D \mid H_k)}{\mathbb{P}(D \mid H_i)}$. Това напълно елиминира трудната за пресмятане вероятност $\mathbb{P}(D)$ от знаменателя.

\subsection{Пример 1: Честна или измамна монета?}

Човек твърди, че хвърля честна монета.
\begin{itemize}
    \item $H_1$: монетата е честна (шанс за ези $1/2$).
    \item $H_2$: монетата е нечестна (за простота да допуснем, че шансът за ези е точно $1/3$).
\end{itemize}
Вие имате голямо доверие на този човек, затова вашето априорно (първоначално) субективно вярване е $\mathbb{P}(H_1) = 0{,}99$ и $\mathbb{P}(H_2) = 0{,}01$. Априорното съотношение е:
\[
\frac{\mathbb{P}(H_1)}{\mathbb{P}(H_2)} = \frac{0{,}99}{0{,}01} = 99
\]
Сега хвърляте монетата 100 пъти и събирате данни $D$: паднали са се 30 езита. Използвайки биномно разпределение, условните вероятности за тези данни са:
\begin{align*}
\mathbb{P}(D \mid H_1) &= \binom{100}{30} \left(\frac{1}{2}\right)^{30} \left(\frac{1}{2}\right)^{70} \\
\mathbb{P}(D \mid H_2) &= \binom{100}{30} \left(\frac{1}{3}\right)^{30} \left(\frac{2}{3}\right)^{70}
\end{align*}
Как се променя нашето съотношение след наблюдаването на данните $D$?
\[
\frac{\mathbb{P}(H_1 \mid D)}{\mathbb{P}(H_2 \mid D)} = \frac{\mathbb{P}(D \mid H_1)}{\mathbb{P}(D \mid H_2)} \cdot \frac{\mathbb{P}(H_1)}{\mathbb{P}(H_2)} = \frac{\mathbb{P}(D \mid H_1)}{\mathbb{P}(D \mid H_2)} \cdot 99
\]
Ако се пресметне точното отношение на биномните вероятности, крайният резултат се оказва приблизително $0{,}03$. Тоест, въпреки първоначалното ви огромно доверие (съотношение $99:1$ в полза на $H_1$), след тези експериментални данни нещата се обръщат напълно: вероятността монетата да е честна става много по-малка от вероятността да е нечестна. Тъй като $H_1$ и $H_2$ образуват пълна група, лесно може да се пресметнат самите вероятности от $\mathbb{P}(H_1 \mid D) + \mathbb{P}(H_2 \mid D) = 1$.

\subsection{Пример 2: Тестове на летището (COVID или тероризъм)}

Да си представим, че на летището е въведена система за бързо тестване (аларма) за COVID инфектирани (или пък терористи). 
Нека събитието $H_1$ означава, че човекът е инфектиран, а $H_2$ - че е здрав.
В общата популация болните са много рядко срещани, да кажем $1\%$:
\[
\mathbb{P}(H_1) = 0{,}01, \quad \mathbb{P}(H_2) = 0{,}99
\]
Тестът (алармата $B$) работи отлично и хваща почти всички инфектирани - с вероятност 99\%:
\[
\mathbb{P}(B \mid H_1) = 0{,}99
\]
Ако човекът е здрав, тестът правилно не звъни с вероятност 80\%. Следователно тестът греши (дава фалшива аларма за здрав човек) с вероятност 20\%:
\[
\mathbb{P}(B \mid H_2) = 0{,}20
\]

Вие сте охрана и алармата $B$ току-що звънна за даден човек. Каква е вероятността той наистина да е инфектиран? Търсим $\mathbb{P}(H_1 \mid B)$.
По формулата на Бейс:
\begin{align*}
\mathbb{P}(H_1 \mid B) &= \frac{\mathbb{P}(B \mid H_1)\mathbb{P}(H_1)}{\mathbb{P}(B \mid H_1)\mathbb{P}(H_1) + \mathbb{P}(B \mid H_2)\mathbb{P}(H_2)} \\
&= \frac{0{,}99 \times 0{,}01}{0{,}99 \times 0{,}01 + 0{,}20 \times 0{,}99} \\
&= \frac{0{,}0099}{0{,}0099 + 0{,}1980} \\
&= \frac{0{,}0099}{0{,}2079} \\
&\approx \frac{1}{21}
\end{align*}

Резултатът е поразителен: въпреки че тестът засича правилно 99\% от истинските болни, вероятността човекът, за когото е звъннала алармата, реално да е болен, е едва 1 към 21 (под 5\%). Причината е, че търсеното събитие в генералната популация е много рядко (0,01), а фалшиво положителните резултати ($0{,}20$) сред масовата част от хората ($0{,}99$) създават огромен брой „фалшиви аларми“, които заглушават истинските. Това илюстрира защо прилагането на подобни системи в реалността изисква изключително висока специфичност на теста, когато търсим нещо много рядко срещано.


\chapter{Случайни величини. Индикатори и дискретни случайни величини}
\section{Случайни величини. Свойства}

\subsection{Въведение и интуиция}
Често при разглеждане на вероятностни модели, множеството от елементарни събития $\Omega$ може да съдържа обекти с много сложна структура — например хора, молекули, траектории на брауново движение или шесторки числа в тотото. Едно такова множество може да бъде характеризирано чрез всевъзможни характеристики (височина, тегло, доходи и др.).

На практика обаче, за даден експеримент или изследване най-често се интересуваме от конкретна характеристика, която по възможност е числова. Искаме на всяко елементарно събитие $\omega \in \Omega$ да съпоставим реално число. Това класическо изображение наричаме случайна величина. То ни позволява да забравим за сложната структура на елементарните събития и да работим изцяло с техните числови характеристики, което за повечето експерименти е абсолютно достатъчно.

\subsection{Формална дефиниция}
Нека $V = (\Omega, \mathcal{A}, \mathbb{P})$ е вероятностно пространство.
Разглеждаме изображение $X: \Omega \rightarrow \mathbb{R}$.

Ако имаме интервал (или произволно множество) $I \subset \mathbb{R}$, под пълния първообраз на $I$ чрез $X$ ще разбираме множеството от онези елементарни събития, чиято числова характеристика попада в $I$:
\[ X^{-1}(I) = \{ \omega \in \Omega : X(\omega) \in I \} \]
Например, ако $I = (a, b)$, множеството е:
\[ X^{-1}((a, b)) = \{ \omega \in \Omega : a < X(\omega) < b \} \]
което за краткост ще записваме просто като $\{ a < X < b \}$, като винаги подразбираме, че това е подмножество на $\Omega$.

За да можем да оперираме математически строго, ни е необходима специална структура. Трябва да можем да придаваме вероятности на множествата, които ни интересуват.

\begin{keydefn}
Изображението $X: \Omega \rightarrow \mathbb{R}$ се нарича \emph{случайна величина}, ако за всяко реално число $b \in \mathbb{R}$, множеството
\[ X^{-1}((-\infty, b)) = \{ \omega \in \Omega : X(\omega) < b \} \]
принадлежи на $\sigma$-алгебрата $\mathcal{A}$. Кратко записваме $\{ X < b \} \in \mathcal{A}$.
\end{keydefn}

Това условие означава, че можем да придаваме вероятност на такива множества — можем да изчисляваме вероятността $\mathbb{P}(X < b)$. Тоест, можем да отговорим на въпроси като „каква е вероятността доходът на човек да е по-малък от 1000 лева?“ или „каква е вероятността да тежи по-малко от 70 кг?“.

\begin{cor}
Ако $X$ е случайна величина, то за всеки две реални числа $a < b$, множествата $\{ a < X < b \}$, $\{ a \le X \le b \}$, $\{ a \le X < b \}$ и т.н. принадлежат на $\sigma$-алгебрата $\mathcal{A}$. Следователно ние можем да измерваме и да изчисляваме вероятностите случайната величина да попада в произволен интервал, бил той отворен, затворен, полуотворен или безкраен.
\end{cor}

\subsection{Операции със случайни величини}
Нека $V = (\Omega, \mathcal{A}, \mathbb{P})$ е вероятностно пространство и нека $X$ и $Y$ са две случайни величини, дефинирани върху него. Тук е изключително важно двете случайни величини да „вървят“ с едно и също вероятностно пространство, за да можем да дефинираме операциите събиране и умножение върху едни и същи елементарни събития. Тогава са изпълнени следните свойства:
\begin{enumerate}
    \item $c X$ е случайна величина за всяко $c \in \mathbb{R}$.
    \item $X + Y$ и $X - Y$ са случайни величини в $V$.
    \item $X Y$ е случайна величина.
    \item $\frac{X}{Y}$ е случайна величина, при условие че събитието $\{Y = 0\}$ е с вероятност 0. Тоест, изображението не връща стойност 0 (например никой няма 0 кг тегло) и можем да делим на него.
\end{enumerate}

\textbf{Скица на доказателство за $c X$:}
Ще разгледаме случая, когато константата $c < 0$. Нека $Z = c X$. Искаме да проверим, че за всяко реално число $b$, събитието $\{ Z < b \} \in \mathcal{A}$:
\[ \{ c X < b \} = \left\{ X > \frac{b}{c} \right\} \]
Тъй като $c < 0$, знакът на неравенството се обръща. Множеството $\left\{ X > \frac{b}{c} \right\}$ е точно $X^{-1}\left(\left(\frac{b}{c}, \infty\right)\right)$, което въз основа на предходното следствие принадлежи на $\sigma$-алгебрата $\mathcal{A}$. (Ако $c > 0$, знакът се запазва и доказателството следва директно от дефиницията).

\textbf{Идея на доказателството за $X + Y$:}
За да покажем, че сумата $X+Y$ е случайна величина, трябва да проверим дали множеството $\{ X + Y < b \}$ е в $\mathcal{A}$. Не е трудно (въпреки че не е съвсем тривиално) да се покаже, че това събитие може да се представи като изброимо обединение по всички рационални числа $q \in \mathbb{Q}$:
\[ \{ X + Y < b \} = \bigcup_{q \in \mathbb{Q}} ( \{ X < q \} \cap \{ Y < b - q \} ) \]
Ако $\omega$ принадлежи на това множество, то сумата $X(\omega) + Y(\omega) < b$. Виждаме, че $\{ X < q \} \in \mathcal{A}$ (по дефиницията за $X$) и $\{ Y < b - q \} \in \mathcal{A}$ (по дефиницията за $Y$). Тъй като $\mathcal{A}$ е $\sigma$-алгебра, сечението на две нейни множества отново е в $\mathcal{A}$. Тъй като множеството на рационалните числа $\mathbb{Q}$ е изброимо, изброимо обединение на елементи от $\mathcal{A}$ също принадлежи на $\mathcal{A}$. С това сложният въпрос се свежда до работа с отделните случайни величини.

\section{Индикаторни функции и техните свойства}

Преди да преминем към конкретната дефиниция на дискретни случайни величини, е полезно да въведем едно понятие, което не е директно свързано с вероятностите, а по-скоро с множества.

\begin{defn}
Нека $A$ е подмножество на $\Omega$. Индикаторната функция $\ind_A : \Omega \rightarrow \{0, 1\}$ се дефинира по следния начин:
\[
\ind_A(\omega) =
\begin{cases}
1, & \text{ако } \omega \in A \\
0, & \text{ако } \omega \notin A
\end{cases}
\]
\end{defn}

Това е най-простият прототип за бинарен експеримент (успех/неуспех). Цялото сложно пространство $\Omega$ се разцепва на две части: интересуваме се само дали събитието се е случило (стойност 1) или не (стойност 0). Използването на индикаторни функции е изключително удобно, защото ни позволява алгебрично да оперираме със сечения и обединения.

\textbf{Свойства на индикаторните функции:}
За произволни множества $A, B \subset \Omega$:
\begin{enumerate}
    \item \textbf{Сечение:} $\ind_{A \cap B} = \ind_A \cdot \ind_B$. \\
    Това равенство е вярно, защото $\ind_{A \cap B}(\omega) = 1$ тогава и само тогава, когато едновременно $\ind_A(\omega) = 1$ и $\ind_B(\omega) = 1$. Ако дори едното е 0, произведението пропада в 0.
    \item \textbf{Обединение:} $\ind_{A \cup B} = \ind_A + \ind_B - \ind_{A \cap B} = \ind_A + \ind_B - \ind_A \cdot \ind_B$. \\
    Тази формула директно отразява принципа на включване и изключване за две множества.
    \item \textbf{Допълнение:} $\ind_{A^c} = 1 - \ind_A$. \\
    Ако индикаторът за $A$ е 1, то $1 - 1 = 0$ (не е в допълнението), и обратно.
\end{enumerate}

С помощта на индикаторите, известни закони (като законите на Де Морган) се доказват алгебрично, без чертаене на диаграми на Вен:
\begin{align*}
\ind_{(A \cup B)^c} &= 1 - \ind_{A \cup B} = 1 - (\ind_A + \ind_B - \ind_A \ind_B) = 1 - \ind_A - \ind_B + \ind_A \ind_B \\
 &= (1 - \ind_A)(1 - \ind_B) = \ind_{A^c} \ind_{B^c} = \ind_{A^c \cap B^c}
\end{align*}
От равенството на индикаторните функции следва равенството на самите множества: $(A \cup B)^c = A^c \cap B^c$.

\begin{prop}
Ако множеството $A \in \mathcal{A}$, то индикаторната функция $\ind_A$ е случайна величина.
\end{prop}

\begin{proof}
Разглеждаме множеството $\{ \ind_A < b \}$ за произволно реално $b$:
\[
\{ \ind_A < b \} =
\begin{cases}
\emptyset, & \text{ако } b \le 0 \\
A^c, & \text{ако } 0 < b \le 1 \\
\Omega, & \text{ако } b > 1
\end{cases}
\]
Първият случай е така, защото индикаторът е винаги $\ge 0$; третият — защото е винаги $\le 1$; в средния случай единствената възможност е индикаторът да е $0$, т.е. $\omega \in A^c$.

Понеже $A \in \mathcal{A}$, неговото допълнение $A^c \in \mathcal{A}$. Празната съвкупност $\emptyset$ и цялото множество $\Omega$ също принадлежат на $\mathcal{A}$. Следователно и в трите възможни случая множеството попада в $\sigma$-алгебрата. Ето защо $\ind_A$ е случайна величина.
\end{proof}

Това е фундаментът („тухличката“), върху който се изгражда теорията на мярката и интеграла, работейки с линейни комбинации на такива индикатори.

\section{Дискретни случайни величини}\label{sec:discrete-random-variables}

\subsection{Дефиниция и разпределение}
Преминаваме към най-реалистичния клас от случайни величини — дискретните. В крайна сметка, светът и компютрите оперират дискретно и ние не можем да съхраняваме и оперираме директно с неизброима безкрайност.

Нека $V = (\Omega, \mathcal{A}, \mathbb{P})$ е вероятностно пространство. Нека разполагаме с редица от различни реални числа $x_j$ (която може да бъде крайна $x_1, \dots, x_n$ или изброимо безкрайна $x_1, x_2, \dots$) и съответна пълна група от събития $H_j \in \mathcal{A}$. (Пълна група означава, че тези събития са непресичащи се и обединението им изчерпва цялото $\Omega$).

\begin{defn}
Случайната величина $X$, дефинирана чрез
\[ X = \sum_{j} x_j \ind_{H_j} \]
се нарича \emph{дискретна случайна величина}.
\end{defn}

Тъй като събитията $H_j$ образуват пълна група, за всяко $\omega \in \Omega$ само един индикатор от тази сума не е нула. Следователно $X(\omega)$ винаги приема една конкретна стойност $x_j$. Самата величина $X$ е гарантирано случайна величина, тъй като (както видяхме) индикаторите са случайни величини, а тяхна линейна комбинация също е случайна величина.

\begin{defn}
Под \emph{разпределение} на дискретната случайна величина $X$ разбираме таблицата, съдържаща всички възможни стойности на $X$ и техните съответни вероятности $p_j$:
\[
\begin{array}{c|c|c|c|c}
X & x_1 & x_2 & \dots & x_n \\
\hline
\mathbb{P} & p_1 & p_2 & \dots & p_n
\end{array}
\]
Тук $p_j$ означава вероятността $X$ да приеме стойност $x_j$, тоест вероятността да настъпи събитието $H_j$ (записвано още като $\{X = x_j\}$). В сила са условията $p_j > 0$ и $\sum_j p_j = 1$.
\end{defn}

С тази таблица нашият модел напълно се абстрахира от сложното вероятностно пространство $\Omega$. Вече не се интересуваме дали зад модела стоят мъже и жени, температури или тегла. Важното е само какви стойности приема $X$ и с какви вероятности.

\subsection{Равенство по разпределение}
Много важно понятие в вероятностите е равенството по разпределение. Нека $X$ и $Y$ са дискретни случайни величини (които могат да идват от съвсем различни експерименти и вероятностни пространства).

\begin{defn}
Казваме, че $X$ е \emph{равна по разпределение} на $Y$ (записваме $X \stackrel{d}{=} Y$, от distribution), ако техните таблици на разпределение съвпадат. Тоест, за всяка възможна стойност $x_j$ е изпълнено:
\[ \mathbb{P}(X = x_j) = \mathbb{P}(Y = x_j) \]
\end{defn}

\begin{example}[Живот на процесор]\label{ex:04-1}
Нека $X$ измерва живота на даден процесор в дни. $X$ може да приема стойности $0, 1, 2, \dots$ (безкрайна редица) с вероятности $p_0, p_1, p_2, \dots$.
В друг модел, нека фирмата има правило да сменя процесорите най-много на 4000-ния ден. Тогава новата случайна величина $Y$ ще приема стойности от 0 до 4000. Нейните вероятности за $k < 4000$ съвпадат с тези на $X$, но за стойността 4000 вероятността акумулира всички останали: $\mathbb{P}(Y=4000) = \sum_{j=4000}^{\infty} p_j$. Следователно, въпреки че описват сходен процес, $X$ и $Y$ нямат еднакви таблици и не са равни по разпределение.

\end{example}

\begin{example}[Монета и зар]\label{ex:04-2}
Нека $X$ отразява хвърлянето на честна монета ($X=1$ за ези, $X=0$ за тура; и двете с вероятност $1/2$). Нека $Y$ отразява хвърлянето на честен зар, като отчитаме само дали числото е четно ($Y=1$) или нечетно ($Y=0$). Вероятностите за четно и нечетно при зар отново са $1/2$. Следователно таблиците им на разпределение са идентични и $X \stackrel{d}{=} Y$, въпреки че физическите генератори на случайност (монета и зар) са коренно различни.

\end{example}

\subsection{Хистограма}
За нагледно представяне на разпределението често се използва хистограма. Тя се построява като по хоризонталната ос се отбелязват стойностите $x_j$, а по вертикалната се издигат стълбчета с височина $p_j$.

\begin{example}\label{ex:04-3}
Да разгледаме хвърлянето на два честни зара, които се хвърлят независимо. Интересуваме се от сумата им. Възможните стойности са числата от 2 до 12. Вероятностите за всяка сума са съответно:
\[
\begin{array}{c|c|c|c|c|c|c|c|c|c|c|c}
\text{Стойност} & 2 & 3 & 4 & 5 & 6 & 7 & 8 & 9 & 10 & 11 & 12 \\
\hline
\text{Вероятност} & \frac{1}{36} & \frac{2}{36} & \frac{3}{36} & \frac{4}{36} & \frac{5}{36} & \frac{6}{36} & \frac{5}{36} & \frac{4}{36} & \frac{3}{36} & \frac{2}{36} & \frac{1}{36}
\end{array}
\]
Хистограмата за този експеримент ще изглежда като симетрична триъгълна пирамида, с най-висок връх (вероятност $6/36$) при стойността 7.

\end{example}

\section{Смяна на променливите на дискретни случайни величини}

Понякога се налага първоначалната случайна величина да бъде трансформирана в друга, по-удобна за нас. Например, ако разполагаме с резултат от зарче (1 до 6), но се интересуваме само дали е четен или не.

\subsection{Функция на една дискретна случайна величина}
Нека $X$ е дискретна случайна величина, а $g: \mathbb{R} \rightarrow \mathbb{R}$ е функция, която е добре дефинирана поне върху възможните стойности на $X$. Тогава $Y = g(X)$ също е дискретна случайна величина:
\[ Y = g(X) = \sum_{j} g(x_j) \ind_{H_j} \]

Ако функцията $g$ изобразява няколко различни стойности $x_i, x_j$ в едно и също число (както при четните числа от зара, които отиват в 1), ние можем да преномерираме и да обединим съответните колонки в таблицата на разпределението, като сумираме техните вероятности:
\[ \mathbb{P}(Y = y) = \sum_{j : g(x_j) = y} \mathbb{P}(X = x_j) \]

\begin{example}\label{ex:04-4}
Да се върнем към модела с процесора, чийто живот е $X \in \{0, 1, 2, \dots\}$. Искаме да знаем само дали процесорът е бил фабрично дефектен в нулевия ден. Дефинираме функцията $g(0) = 0$ (дефектен) и $g(x) = 1$ за $x \ge 1$ (работещ поне 1 ден). Тогава $Y = g(X)$ ще приема само две стойности:
\[ Y = 0 \cdot \ind_{\{X=0\}} + 1 \cdot \ind_{\{X \ge 1\}} \]
Табличката на разпределение на $Y$ става изключително проста: стойност $0$ с вероятност $p_0$ и стойност $1$ с вероятност $1 - p_0$. Така си спестяваме дългото и сложно разпределение на $X$, фокусирайки се само върху дефектите.

\end{example}

\subsection{Функция на две дискретни случайни величини}
Когато сменяме променливите на две случайни величини едновременно, е ключово те да са дефинирани върху \emph{едно и също} вероятностно пространство $V$.

Нека $X$ и $Y$ са две дискретни случайни величини във $V$, със съответните си пълни групи събития $H_j$ и $H_k^*$. Нека $g: \mathbb{R} \times \mathbb{R} \rightarrow \mathbb{R}$ е функция на две променливи (например $g(x,y) = x+y$). Тогава случайната величина $Z = g(X, Y)$ е:
\[ Z = g(X, Y) = \sum_{j, k} g(x_j, y_k) \ind_{H_j} \ind_{H_k^*} = \sum_{j, k} g(x_j, y_k) \ind_{H_j \cap H_k^*} \]
Тук използваме свойството на индикаторните функции, че произведението им е равно на индикатора на тяхното сечение. Това сечение $H_j \cap H_k^*$ отговаря на едновременното настъпване на събитието $X$ да приеме стойност $x_j$ и $Y$ да приеме стойност $y_k$.

\section{Независимост на дискретни случайни величини}

Концепцията за независимост е ключова. Ако имаме две дискретни случайни величини $X$ и $Y$ във $V$, интуитивно те са независими, ако стойността, която приема едната, не носи никаква информация за стойността на другата.

\begin{defn}[независимост на две случайни величини]
$X$ и $Y$ са \emph{независими} ($X \perp Y$), тогава и само тогава, когато за всеки две възможни техни стойности $x_i$ и $y_j$ е изпълнено:
\[ \mathbb{P}(X = x_i, Y = y_j) = \mathbb{P}(X = x_i) \mathbb{P}(Y = y_j) \]
Тоест, вероятността двете случайни величини едновременно да приемат конкретни стойности (настъпването на сечението на събитията) се разпада като произведение от техните индивидуални вероятности. Това е директен аналог на независимостта на събития: $\mathbb{P}(A_i \cap B_j) = \mathbb{P}(A_i)\mathbb{P}(B_j)$. За да направим това сечение, отново ни трябва $X$ и $Y$ да са в едно и също вероятностно пространство.
\end{defn}

\begin{defn}[независимост в съвкупност]
Нека имаме $n$ дискретни случайни величини $X_1, X_2, \dots, X_n$. Те се наричат независими в съвкупност, тогава и само тогава, когато за всяко подмножество от техни индекси $\{j_1, j_2, \dots, j_m\} \subset \{1, \dots, n\}$ (с размер $m \ge 2$) и за всички техни възможни стойности $x_{j_1}, \dots, x_{j_m}$ е вярно:
\[ \mathbb{P}(X_{j_1} = x_{j_1}, X_{j_2} = x_{j_2}, \dots, X_{j_m} = x_{j_m}) = \mathbb{P}(X_{j_1} = x_{j_1}) \mathbb{P}(X_{j_2} = x_{j_2}) \dots \mathbb{P}(X_{j_m} = x_{j_m}) \]
Записът може да изглежда сложен, но идеята му е проста: каквото и подмножество от случайните величини да вземете, съвместната вероятност за всяка тяхна конфигурация се представя като произведение от маргиналните (индивидуалните) вероятности на всяко едно от събитията.
\end{defn}

\section{Задачи}

\textbf{Домашна работа.} Нека $X$ и $Y$ са случайни величини върху едно и също вероятностно пространство. Довършете доказателството на равенството
\[
\{X+Y<b\}
=\bigcup_{q\in\mathbb{Q}}\bigl(\{X<q\}\cap\{Y<b-q\}\bigr),
\]
като докажете включването отляво надясно.


\chapter{Функция на разпределение. Математическо очакване и дисперсия}
\section{Функция на разпределение}

Функцията на разпределение е едно от основните понятия в теорията на вероятностите. На английски език терминът е \emph{Cumulative Distribution Function} (често съкращавано като CDF).

\begin{keydefn}[функция на разпределение]\label{def:cdf}
Нека $X$ е случайна величина, дефинирана върху вероятностно пространство
$(\Omega, \mathcal{A}, \mathbb{P})$. \emph{Функция на разпределение} на $X$ наричаме
\[ F_X(x) = \mathbb{P}(X < x), \qquad x \in \mathbb{R} , \]
тоест вероятността $X$ да приеме стойност, строго по-малка от $x$.
\end{keydefn}

\emph{Бележка относно терминологията:} В англоезичната литература (а и не само) функцията на разпределение по принцип се дефинира с нестрого неравенство, т.е. $F(x) = \mathbb{P}(X \le x)$. Въпреки това в част от българската традиция — и в този курс — се използва строго неравенство ($X < x$); срещат се и български учебници с нестрогия вариант. На практика това не води до съществени промени, особено при непрекъснати случайни величини, разгледани по-късно. Тази терминологична особеност следва да се има предвид.

\begin{prop}[Свойства на функцията на разпределение]\label{prop:cdf-props}
За всяка случайна величина $X$ функцията $F_X$ притежава следните свойства:
\begin{enumerate}
    \item[а)] $0 \le F_X(x) \le 1$ за всяко $x$;
    \item[б)] $F_X$ е монотонно ненамаляваща: ако $x_1 < x_2$, то $F_X(x_1) \le F_X(x_2)$;
    \item[в)] $\displaystyle \lim_{x \to -\infty} F_X(x) = 0$ и $\displaystyle \lim_{x \to +\infty} F_X(x) = 1$;
    \item[г)] $F_X$ е непрекъсната отляво (при конвенцията със строго неравенство).
\end{enumerate}
\end{prop}

\begin{proof}
а) следва от това, че $F_X(x)$ е вероятност. За б): ако $x_1 < x_2$, то
$\{X < x_1\} \subseteq \{X < x_2\}$, а вероятността е монотонна
(твърдение~\crosslecture{prop:monotone}{за монотонност от лекция 2}). За в): събитията
$\{X < x\}$ намаляват към $\emptyset$ при $x \to -\infty$ и растат към $\Omega$ при
$x \to +\infty$; остава да приложим непрекъснатостта на вероятностната мярка. Свойство г)
се проверява по същия начин чрез растяща редица $x_n \uparrow x$.
\end{proof}

Точките, в които $F_X$ има скок, са точно тези $x$, за които $\mathbb{P}(X = x) > 0$;
големината на скока е самата $\mathbb{P}(X = x)$. Това ще ни трябва по-нататък при
сходимостта по разпределение, където се изискват точки на непрекъснатост.

\subsection{Връзка с разпределението на дискретна случайна величина}

Дискретните случайни величини, разгледани в лекция 4, могат да бъдат описани чрез техните таблици на разпределение. Функцията на разпределение обаче е много по-общ и често по-удобен инструмент. Тя капсулира цялата информация за вероятностите, свързани със случайната величина $X$.

Нека $X$ е дискретна случайна величина, която приема стойности $x_1 < x_2 < x_3 < \dots$ със съответни вероятности $p_1, p_2, p_3, \dots$. (Допускаме, че стойностите са подредени в нарастващ ред).
Как бихме изчислили функцията на разпределение за тази случайна величина?
Тя има следния вид:
\[
F(x) =
\begin{cases}
0, & \text{ако } x \le x_1 \\
p_1, & \text{ако } x \in (x_1, x_2] \\
p_1 + p_2, & \text{ако } x \in (x_2, x_3] \\
p_1 + p_2 + p_3, & \text{ако } x \in (x_3, x_4] \\
\vdots & \vdots
\end{cases}
\]

Защо е така? Ако $x \le x_1$, тогава вероятността $X$ да приеме стойност, строго по-малка от $x_1$, е 0, тъй като $x_1$ е най-малката възможна стойност. Когато $x$ попадне в интервала $(x_1, x_2]$, единствената възможна стойност на $X$, която е строго по-малка от $x$, е $x_1$. Следователно вероятността е $p_1$. Този процес продължава стъпаловидно. Вижда се, че от тази функция лесно можем да възстановим оригиналната таблица с вероятностите, и обратно. Функцията на разпределение дава единен начин за кодиране на информацията на всяка една случайна величина, независимо дали е дискретна или непрекъсната.

\begin{example}[Индикаторна функция]\label{ex:05-1}
Нека $X$ е индикатор на дадено събитие $A$, тоест $X = \ind_A$. Това е най-простата дискретна случайна величина. Тя приема стойност 1 с вероятност $p_1 = \mathbb{P}(A)$ и стойност 0 с вероятност $p_0 = 1 - p_1$.
Графиката на функцията на разпределение $F(x)$ за този индикатор е стъпаловидна. Тя е 0 за всички стойности до 0 включително. В интервала $(0, 1]$ функцията скача на ниво $1 - p_1$, тъй като там $X$ може да бъде само 0. За стойности, по-големи от 1, функцията скача на ниво 1, защото вероятността $X$ да е по-малко от число, по-голямо от 1, е сигурно събитие (тъй като $X$ е или 0, или 1).
\[
F(x) =
\begin{cases}
0, & \text{ако } x \le 0 \\
1 - p_1, & \text{ако } x \in (0, 1] \\
1, & \text{ако } x > 1
\end{cases}
\]

\begin{lecfigure}[Функцията на разпределение на индикатор $\ind_A$ при
$p_1 = \mathbb{P}(A)$. Стъпаловидна е, с два скока — в $0$ и в $1$. Плътната точка
показва коя стойност се приема в самата точка на скока: при конвенцията
$F(x) = \mathbb{P}(X < x)$ функцията е непрекъсната отляво.\label{fig:cdf-indicator}]
\begin{tikzpicture}[x=2.1cm, y=2.3cm]
  \draw[axisline] (-0.75,0) -- (2.3,0) node[right, font=\small] {$x$};
  \draw[axisline] (0,-0.12) -- (0,1.32) node[above, font=\small] {$F(x)$};
  % level guides
  \draw[guide] (0,1) -- (2.1,1);
  \draw[guide] (0,0.62) -- (1,0.62);
  \node[left, font=\small] at (-0.03,1) {$1$};
  \node[left, font=\small] at (-0.03,0.62) {$1-p_1$};
  \node[below, font=\small] at (1,-0.04) {$1$};
  % the three pieces
  \draw[cdfstep] (-0.7,0) -- (0,0);
  \draw[cdfstep] (0,0.62) -- (1,0.62);
  \draw[cdfstep] (1,1) -- (2.1,1);
  % closed / open endpoints
  \fill[defframe] (0,0) circle (1.7pt);
  \draw[defframe, fill=white, line width=0.7pt] (0,0.62) circle (1.7pt);
  \fill[defframe] (1,0.62) circle (1.7pt);
  \draw[defframe, fill=white, line width=0.7pt] (1,1) circle (1.7pt);
\end{tikzpicture}
\end{lecfigure}

\end{example}

\section{Математическо очакване}

Следващата основна тема е свързана с две изключително важни числови характеристики на случайните величини: \textbf{математическо очакване} (на английски \emph{expectation}) и \textbf{дисперсия} (на английски \emph{variance} или \emph{variation}).

\subsection{Дефиниция на математическо очакване}

Първо ще разгледаме формалната дефиниция за дискретна случайна величина.

\begin{keydefn}
Нека $X$ е дискретна случайна величина, приемаща стойности $x_1, x_2, \dots$ със съответни вероятности $p_1, p_2, \dots$. Тогава под \emph{математическо очакване} на $X$, което бележим с $\E X$, разбираме сумата:
\[ \E X = \sum_{j} x_j p_j \]
ако редът от абсолютните стойности е сходящ, т.е.:
\[ \sum_{j} |x_j| p_j < \infty \]
\end{keydefn}

\textbf{Коментари към дефиницията:}
\begin{enumerate}
    \item Ако $X$ приема \emph{краен брой стойности} $\{x_1, \dots, x_n\}$, то сумата винаги е крайна и математическото очакване винаги съществува. То се задава просто като $\E X = \sum_{j=1}^n x_j p_j$.
    \item Ако $X$ приема \emph{безкраен брой стойности}, е възможно очакването да не съществува. Например, ако разгледаме случайна величина $X$, която приема стойности $j = 1, 2, 3, \dots$ с вероятности $p_j = \frac{6}{\pi^2 j^2}$, тогава очакването би било сума от вида $\sum_{j} j \cdot \frac{6}{\pi^2 j^2} = \frac{6}{\pi^2} \sum_{j} \frac{1}{j}$. Това е хармоничният ред, който е разходящ (сумата му е безкрайност), следователно очакването за тази случайна величина не съществува.
\end{enumerate}

\subsection{Интерпретация на математическото очакване}

\textbf{Физическа интерпретация:} Математическото очакване има много ясен физически смисъл. То играе ролята на \emph{център на тежестта}. Представете си, че имате единична маса, която е разпределена в различни точки (локации) $x_j$ по числовата ос, и всяка точка има маса $p_j$. Тогава очакването $\E X$ е точно координатата на центъра на тежестта на тази система. Както в механиката центърът на тежестта е точката, в която обектът се уравновесява, така и очакването е „центърът“ на разпределението на вероятностите.

\textbf{Математическа интерпретация (Минимизиране на грешката):} Представете си, че трябва да изберете само едно реално число $a$, което най-добре да характеризира дадена случайна величина $X$. Какво означава „най-добре“? Ако искаме да минимизираме \emph{средноквадратичната грешка}, тоест търсим числото $a$, което минимизира израза:
\[ \sum_{j} (x_j - a)^2 p_j \]
се оказва, че този минимум се достига точно когато $a = \E X$. В средноквадратичен смисъл, очакването най-добре описва „типичната“ стойност на случайната величина. Това свойство е фундаментално и ще бъде ключово по-нататък, когато се изучава условното математическо очакване.

\subsection{Примери за математическо очакване}

\begin{example}[Равномерно разпределение]\label{ex:05-2}
Нека $X$ е равномерно разпределена върху числата $\{1, 2, \dots, n\}$. Тоест, вероятността да приеме всяка една от тези стойности е $1/n$ ($p_j = 1/n$ за $j=1,\dots,n$).
Тогава очакването е:
\[ \E X = \sum_{j=1}^n x_j p_j = \sum_{j=1}^n j \frac{1}{n} = \frac{1}{n} \sum_{j=1}^n j \]
Това е просто средноаритметичното на числата.

\end{example}

\begin{example}[Две стратегии на рулетка]\label{ex:05-3}
Нека разгледаме игра на европейска рулетка, която има 37 числа (от 0 до 36).
\begin{itemize}
    \item \emph{Стратегия 1 (Залагане на цвят):} Залагаме на „черно“. Има 18 черни, 18 червени числа и една зелена нула. Залагаме 1 лев. Печелим 1 лев с вероятност $18/37$ и губим 1 лев (тоест печелим $-1$) с вероятност $19/37$. Нека $X$ е печалбата от тази стратегия.
    \[ \E X = (-1) \cdot \frac{19}{37} + 1 \cdot \frac{18}{37} = -\frac{1}{37} \]

    \item \emph{Стратегия 2 (Залагане на конкретно число):} Залагаме 1 лев на числото „1“. Печелим 35 лева чисто (връщат ни 36 лева) с вероятност $1/37$. Губим 1 лев с вероятност $36/37$. Нека $Y$ е печалбата от тази стратегия.
    \[ \E Y = (-1) \cdot \frac{36}{37} + 35 \cdot \frac{1}{37} = \frac{-36 + 35}{37} = -\frac{1}{37} \]
\end{itemize}

От гледна точка на математическото очакване, двете игри са напълно еквивалентни. Средно, дългосрочно, при всеки заложен лев вие губите $1/37$ от лева (около $2{,}7$ стотинки), независимо коя от двете стратегии ще изберете. Това показва, че очакването е мярка за средния резултат. Въпреки това, интуицията ни подсказва, че двете игри се усещат много различно. Тази разлика ще бъде обяснена и количествено измерена, когато въведем понятието \emph{дисперсия}.

\end{example}

\subsection{Свойства на математическото очакване}

Преди да продължим, ще формулираме и докажем основните свойства на очакването.

\begin{keylem}[Свойства на очакването]
Нека $X$ и $Y$ са две дискретни случайни величини, за които очакванията $\E X$ и $\E Y$ съществуват. Тогава са в сила следните свойства:
\begin{enumerate}
    \item[а)] Ако $X = C$ (константа), то $\E X = C$.
    \item[б)] Ако $c \in \mathbb{R}$ е константа, то $\E(cX) = c\E X$.
    \item[в)] $\E(X + Y) = \E X + \E Y$ (линейност на очакването).
    \item[г)] Ако в допълнение $X$ и $Y$ са \textbf{независими} ($X \perp\!\!\!\perp Y$), то $\E(XY) = \E X \cdot \E Y$.
\end{enumerate}
\end{keylem}

\begin{proof}
а) Ако $X$ приема само една единствена стойност $C$, то $\mathbb{P}(X=C) = 1$. Следователно, по дефиниция:
\[ \E X = C \cdot \mathbb{P}(X=C) = C \cdot 1 = C \]

б) Нека $Z = cX$. Можем да разглеждаме $Z$ като функция на $X$, т.е. $Z = g(X)$, където $g(x) = cx$. Очакването на функция от случайна величина се пресмята като:
\[ \E Z = \sum_{j} g(x_j) p_j = \sum_{j} c x_j p_j = c \sum_{j} x_j p_j = c\E X \]
Константата просто се изнася пред сумата.

в) Нека $Z = X + Y = g(X,Y)$, където функцията е $g(x,y) = x+y$. Използваме дефиницията за математическо очакване за функция от две случайни величини. Обхождаме всички възможни двойки стойности $x_i$ и $y_j$:
\[ \E(X+Y) = \sum_{i, j} (x_i + y_j) \mathbb{P}(X=x_i, Y=y_j) \]
Разкриваме скобите:
\[ = \sum_{i, j} x_i \mathbb{P}(X=x_i, Y=y_j) + \sum_{i, j} y_j \mathbb{P}(X=x_i, Y=y_j) \]
Сега ще прегрупираме сумите. Разглеждаме първата сума и първо сумираме по $j$:
\[ = \sum_{i} x_i \left( \sum_{j} \mathbb{P}(X=x_i, Y=y_j) \right) + \sum_{j} y_j \left( \sum_{i} \mathbb{P}(X=x_i, Y=y_j) \right) \]
Нека обърнем внимание на вътрешната сума $\sum_{j} \mathbb{P}(X=x_i, Y=y_j)$. Тъй като събитията $\{Y=y_j\}$ за всички възможни $j$ образуват пълна група от събития (те разделят достоверното събитие $\Omega$ на непресичащи се части), то обединението им е $\Omega$. Следователно събитието $\{X=x_i\}$ може да бъде записано като обединение на сеченията $\{X=x_i\} \cap \{Y=y_j\}$. Тогава:
\[ \sum_{j} \mathbb{P}(X=x_i, Y=y_j) = \mathbb{P}(X=x_i) \]
Прилагайки същото логическо разсъждение и за втората сума, получаваме:
\[ = \sum_{i} x_i \mathbb{P}(X=x_i) + \sum_{j} y_j \mathbb{P}(Y=y_j) = \E X + \E Y \]

г) За независимост. Нека разгледаме $Z = X \cdot Y$.
\[ \E(XY) = \sum_{i, j} x_i y_j \mathbb{P}(X=x_i, Y=y_j) \]
Тъй като по условие $X$ и $Y$ са независими, вероятността на сечението е равна на произведението от вероятностите:
\[ \mathbb{P}(X=x_i, Y=y_j) = \mathbb{P}(X=x_i) \mathbb{P}(Y=y_j) \]
Заместваме в сумата и групираме членовете:
\[ \E(XY) = \sum_{i, j} x_i y_j \mathbb{P}(X=x_i)\mathbb{P}(Y=y_j) = \left( \sum_{i} x_i \mathbb{P}(X=x_i) \right) \left( \sum_{j} y_j \mathbb{P}(Y=y_j) \right) = \E X \cdot \E Y \]

\end{proof}

\begin{prop}[Неотрицателност на очакването]\label{prop:nonneg-exp}
Ако $X$ е неотрицателна случайна величина (тоест вероятността $X \ge 0$ е единица), то $\E X \ge 0$.
Нещо повече, ако $X \ge 0$ и $\E X = 0$, то $\mathbb{P}(X = 0) = 1$.
\end{prop}

\begin{proof}
По дефиниция $\E X = \sum_{j} x_j p_j$. Тъй като всички стойности $x_j$ са неотрицателни и вероятностите $p_j$ са неотрицателни, то цялата сума също е неотрицателна.

За втората част: сумата $\sum_j x_j p_j$ е сума от неотрицателни събираеми и е равна на нула
само ако всяко събираемо е нула. Значи за всяко $j$ с $x_j > 0$ задължително $p_j = 0$;
следователно цялата вероятностна маса е съсредоточена в стойността $0$.
\end{proof}

\begin{supp}[Очакване чрез опашките]\label{supp:tail-sum}
Ако $X$ приема стойности в $\{0, 1, 2, \dots\}$, то
\[ \E X = \sum_{n \ge 0} \mathbb{P}(X > n) . \]
Наистина, $\mathbb{P}(X > n) = \sum_{k > n} \mathbb{P}(X = k)$, и като сменим реда на
сумиране,
\[
\sum_{n \ge 0} \sum_{k > n} \mathbb{P}(X = k)
 = \sum_{k \ge 1} \sum_{n=0}^{k-1} \mathbb{P}(X = k)
 = \sum_{k \ge 1} k\, \mathbb{P}(X = k) = \E X .
\]

\begin{lecfigure}[Смяната на реда на сумиране, нагледно. В стъпаловидната таблица е
разположена по една клетка за всяка двойка $(n,k)$ с $k > n$; в клетката стои
$p_k = \mathbb{P}(X=k)$. Съберем ли \emph{по редове}, ред $n$ дава опашката
$\mathbb{P}(X > n)$ — това е лявата страна. Съберем ли \emph{по стълбове}, стълб $k$ съдържа
точно $k$ клетки (за $n = 0, \dots, k-1$) и дава $k\,p_k$ — това е дефиницията на $\E X$.
Едни и същи числа, преброени по два начина.\label{fig:tail-sum}]
\begin{tikzpicture}[x=1.13cm, y=0.66cm]
  \def\NC{5}
  % cells: one for every pair (n,k) with k > n
  \foreach \n in {0,...,4} {
    \pgfmathtruncatemacro{\kstart}{\n+1}
    \foreach \k in {\kstart,...,\NC} {
      \fill[defback] ({\k-1.45},{-\n-0.4}) rectangle ({\k-0.55},{-\n+0.4});
      \draw[defframe, line width=0.4pt] ({\k-1.45},{-\n-0.4}) rectangle ({\k-0.55},{-\n+0.4});
      \node[font=\scriptsize] at ({\k-1},{-\n}) {$p_{\k}$};}
    % row label = the tail
    \node[left, font=\scriptsize] at (-0.62,{-\n}) {$\mathbb{P}(X > \n) =$};
    % the row continues
    \node[font=\scriptsize] at ({\NC+0.15},{-\n}) {$\cdots$};}
  \node[left, font=\scriptsize] at (-0.62,-5.05) {$\vdots$};
  \foreach \k in {1,...,\NC} { \node[font=\scriptsize] at ({\k-1},-5.05) {$\vdots$}; }
  % column totals
  \draw[exrule, line width=0.6pt] (-0.5,-5.5) -- ({\NC+0.45},-5.5);
  \foreach \k in {1,...,\NC} {
    \node[font=\scriptsize, below] at ({\k-1},-5.62) {$\k\,p_{\k}$}; }
  \node[font=\scriptsize, below] at ({\NC+0.15},-5.62) {$\cdots$};
  \node[left, font=\scriptsize] at (-0.62,-6.05) {стълбови суми:};
  % the two readings
  \node[right, font=\small, defframe, align=left] at ({\NC+0.75},-1.4)
       {по редове:\\[1pt] $\sum_{n \ge 0} \mathbb{P}(X > n)$};
  \node[right, font=\small, defframe, align=left] at ({\NC+0.75},-4.3)
       {по стълбове:\\[1pt] $\sum_{k \ge 1} k\,p_k = \E X$};
\end{tikzpicture}
\end{lecfigure}
Формулата често спестява работа. Например при $X \sim \Geo(p)$ имаме
$\mathbb{P}(X > n) = q^{\,n+1}$, откъдето веднага
$\E X = \sum_{n \ge 0} q^{\,n+1} = \frac{q}{1-q} = \frac{q}{p}$,
без да се налага да сумираме реда $\sum_k k q^k p$.
\end{supp}

\section{Пример от практиката: Групово тестване (Pool Testing)}

Този пример беше особено актуален по време на пандемията от COVID-19. Представете си, че искаме да тестваме популация, в която хората са здрави (незаразени) с вероятност $p = 0{,}95$ и са заразени с вероятност $1-p = 0{,}05$. Допускаме, че статусите на отделните хора са независими един от друг (което е разумно, освен ако не идват от едно и също домакинство).

Вместо да тестваме всеки човек индивидуално, можем да вземем проби от група от $n$ души и да ги смесим в един общ тест.
\begin{itemize}
    \item Ако общият тест е \emph{отрицателен}, това означава, че всички $n$ човека са здрави. Поради независимостта, вероятността за това събитие е $p^n$. В този случай сме изразходвали само 1 тест за всичките $n$ души.
    \item Ако общият тест е \emph{положителен}, което се случва с вероятност $1 - p^n$, това означава, че поне един човек от групата е заразен. За да разберем кой точно е заразен, се налага да тестваме всеки от тях поотделно. Така сме изразходвали 1 първоначален тест плюс още $n$ индивидуални теста, което прави общо $n+1$ теста.
\end{itemize}

Нека $X_n$ е случайната величина, отчитаща броя на изразходваните тестове за група от $n$ човека. Тя има следното разпределение:
\[
X_n =
\begin{cases}
1, & \text{с вероятност } p^n \\
n+1, & \text{с вероятност } 1-p^n
\end{cases}
\]

Интересува ни средният (очакваният) брой тестове, който се пада на един човек. Това се дава от отношението $\frac{\E X_n}{n}$:
\[ \frac{\E X_n}{n} = \frac{1}{n} \Big( 1 \cdot p^n + (n+1)(1-p^n) \Big) \]
Нека опростим този израз:
\[ \frac{1}{n} \Big( p^n + n + 1 - np^n - p^n \Big) = \frac{n + 1 - np^n}{n} = 1 + \frac{1}{n} - p^n \]

Следователно, пропорцията изразходвани тестове спрямо броя на хората е $1 + \frac{1}{n} - p^n$. Нашата цел е да \textbf{минимизираме} този израз като изберем оптималния размер на групата $n$ (при $n \ge 1$).
Ако разгледаме функцията $f(x) = 1 + \frac{1}{x} - p^x$ и я изследваме за минимум чрез нейната производна ($f'(x) = 0$), за стойност $p = 0{,}95$ ще установим, че минимумът се достига най-близо до цялото число $n = 5$.

Нека пресметнем средния брой тестове на човек при $n = 5$:
\[ \min_{n \ge 1} \frac{\E X_n}{n} = \frac{\E X_5}{5} \approx 0{,}43 \]
Това означава, че на един човек се падат средно $0{,}43$ теста. Оттук намираме очаквания брой тестове за група от петима:
\[ \E X_5 = 5 \cdot 0{,}43 = 2{,}15 \]\footnote{Числото $0{,}43$ е закръглено, затова и произведението е закръглено нагоре. Точната стойност е $\E X_5 = p^5 + 6(1-p^5) = 6 - 5 \cdot 0{,}95^5 \approx 2{,}131$. Изводът не се променя.}

Извод: Вместо да тестваме 5 човека с 5 теста, използвайки тази стратегия, ние очакваме да изразходваме средно само $2{,}15$ теста. Ако имаме 5000 души, ние бихме ги тествали със средно около 2150 теста (вместо 5000). Това е огромно спестяване на ресурси, напълно оправдано чисто математически.

\section{Дисперсия}

Както видяхме в пример~\ref{ex:05-3} с рулетката, математическото очакване само по себе си не винаги е достатъчно, за да опише пълноценно поведението на една случайна величина. То показва къде се намира „центърът на тежестта“, но не ни казва колко са „разпръснати“ възможните стойности около този център. За целта въвеждаме понятието \textbf{дисперсия}.

\subsection{Дефиниция на дисперсия}

\begin{keydefn}
Нека $X$ е дискретна случайна величина, която притежава крайно математическо очакване $\E X$. Тогава \emph{дисперсията} на $X$, която обозначаваме с $\Var X$, се дефинира като математическото очакване на квадрата на отклонението на $X$ от нейното очакване:
\[ \Var X := \E(X - \E X)^2 = \sum_{j} (x_j - \E X)^2 p_j \]
стига тази сума да е крайна, т.е. $\Var X < \infty$.
\end{keydefn}

Ако $X$ приема краен брой стойности, дисперсията винаги съществува. Дисперсията на практика представлява минимизираната средноквадратична грешка (за която говорихме по-рано), когато за характеризираща константа изберем точно математическото очакване $a = \E X$.

\begin{defn}[стандартно отклонение]\label{def:stddev}
Корен квадратен от дисперсията се нарича \emph{стандартно отклонение} (standard deviation):
\[ \sigma_X = \sqrt{\Var X} . \]
То има същата мерна единица като самата случайна величина.
\end{defn}

\subsection{Интерпретация на дисперсията}

Нека се върнем към пример~\ref{ex:05-3} с рулетката:
\begin{itemize}
    \item За първата стратегия ($X$, залагане на цвят): стойностите са $-1$ и $1$. Очакването е $-1/37$. Отклонението на възможните печалби/загуби от очакването е много малко. Следователно $\Var X$ ще е относително малка (около 1).
    \item За втората стратегия ($Y$, залагане на число): стойностите са $-1$ и $35$. Очакването отново е $-1/37$. Въпреки че средната загуба е същата, тук размахът на възможните стойности е огромен. Ако извадите очакването от 35, получавате голямо отклонение, което повдигнато на квадрат дава още по-голям принос към дисперсията. Следователно $\Var Y \gg \Var X$.
\end{itemize}
Дисперсията отчита \emph{волатилността}. В първата игра вие бавно и сигурно губите парите си с малки стъпки. Във втората игра имате шанса да реализирате голяма печалба, макар че дългосрочното математическо очакване остава същото. Разликата в поведението ще стане още по-ясна, когато се изучава Централната гранична теорема.

\begin{example}[Случайно блуждаене (Random Walk)]\label{ex:05-4}
Представете си център на едно въже и $n$ на брой човека, всеки от които може да дърпа въжето с единична сила. Всеки човек се позиционира случайно: отдясно на центъра с вероятност $1/2$ и отляво с вероятност $1/2$.
Нека $X_j$ е силата на $j$-тия човек. Ако е отдясно, $X_j = 1$; ако е отляво, $X_j = -1$.
Резултантната сила е сумата:
\[ F = \sum_{j=1}^n X_j \]
Математическото очакване на силата на един човек е $\E X_j = 1 \cdot \frac{1}{2} + (-1) \cdot \frac{1}{2} = 0$. Поради линейността на очакването, средната резултантна сила също е нула: $\E F = 0$.
Някои хора погрешно смятат, че щом средната сила е нула, то с нарастването на броя на хората $n$, вероятността силата да е точно нула клони към 1. Това е груба грешка. Напротив, дисперсията на силата е много голяма и расте линейно с $n$ ($\Var F = n \cdot \Var X_1 = n$). Следователно стандартното отклонение расте като $\sqrt{n}$. Вероятността частицата (въжето) да остане в покой всъщност е от порядъка на $1/\sqrt{n}$, което клони към $0$ при големи $n$. Частицата прави силно хаотично движение (блуждаене) напред-назад, без да стои на едно място.

\end{example}

\subsection{Свойства на дисперсията}

Ще докажем няколко важни свойства, които значително улесняват работата с дисперсии.

\begin{keythm}[Алтернативна формула за пресмятане на дисперсия]
Ако $X$ е дискретна случайна величина с крайна дисперсия (т.е. $\E X < \infty$ и $\Var X < \infty$), тогава:
\[ \Var X = \E X^2 - (\E X)^2 \]
\end{keythm}

\begin{proof}
Изхождаме от дефиницията и разкриваме скобите под очакването:
\[ \Var X = \E(X - \E X)^2 = \E\big( X^2 - 2X\E X + (\E X)^2 \big) \]
Използваме факта, че $\E X$ е просто константа, и прилагаме линейността на математическото очакване:
\[ = \E X^2 - \E(2X\E X) + \E\big( (\E X)^2 \big) \]
\[ = \E X^2 - 2\E X \cdot \E X + (\E X)^2 \]
\[ = \E X^2 - 2(\E X)^2 + (\E X)^2 = \E X^2 - (\E X)^2 \]
\end{proof}

Тази формула често е много по-удобна за практически пресмятания от самата дефиниция.

\begin{cor}
Нека $X$ е дискретна случайна величина и $\Var X < \infty$. Тогава:
\[ \E X^2 \ge (\E X)^2 \]
\end{cor}

\begin{proof}
Тъй като величината $Y = (X - \E X)^2$ е винаги неотрицателна, нейното очакване също е неотрицателно. Следователно $0 \le \Var X = \E Y = \E X^2 - (\E X)^2$. Като прехвърлим от другата страна, получаваме желаното неравенство.
\end{proof}

\begin{prop}[Свойства на дисперсията]\label{prop:var-props}
Нека $X$ и $Y$ са дискретни случайни величини с крайни дисперсии. Тогава са в сила следните свойства:
\begin{enumerate}
    \item[а)] $\Var X = 0$ тогава и само тогава, когато $X$ е константа с вероятност
    единица, т.е. $\mathbb{P}(X = \E X) = 1$.
    \item[б)] Ако $Z = cX$ за някаква константа $c \in \mathbb{R}$, то $\Var Z = c^2 \Var X$.
    \item[в)] Ако в допълнение $X$ и $Y$ са \textbf{независими} ($X \perp\!\!\!\perp Y$), то $\Var(X + Y) = \Var X + \Var Y$.
\end{enumerate}
\end{prop}

\begin{proof}
а) Ако $X = C$, знаем, че $\E X = C$. Замествайки в дефиницията за дисперсия, получаваме:
\[ \Var X = \E(C - C)^2 = \E(0) = 0 \]
Това е естествено — щом величината е константа, тя не се отклонява от средната си стойност, следователно разсейването (дисперсията) е нулево.
Обратно, ако $\Var X = 0$, то величината $Y = (X - \E X)^2$ е неотрицателна с нулево
очакване, откъдето по твърдение~\ref{prop:nonneg-exp} следва $\mathbb{P}(Y = 0) = 1$, тоест
$X = \E X$ с вероятност единица.

б) За $Z = cX$, намираме дисперсията:
\[ \Var(cX) = \E(cX)^2 - (\E(cX))^2 \]
\[ = \E(c^2 X^2) - (c\E X)^2 \]
\[ = c^2 \E X^2 - c^2 (\E X)^2 = c^2 \Big( \E X^2 - (\E X)^2 \Big) = c^2 \Var X \]
Важно е да се отбележи, че константата излиза от дисперсията повдигната на квадрат. Това е логично, тъй като мерната единица на дисперсията е на втора степен спрямо оригиналната величина (ако мерите в метри, дисперсията е в квадратни метри).

в) Да намерим дисперсията на сума от две независими случайни величини $X$ и $Y$:
\[ \Var(X + Y) = \E(X + Y)^2 - \big(\E(X + Y)\big)^2 \]
Разкриваме скобите:
\[ = \E(X^2 + 2XY + Y^2) - (\E X + \E Y)^2 \]
Използваме линейността на очакването:
\[ = \E X^2 + 2\E(XY) + \E Y^2 - \Big( (\E X)^2 + 2\E X\E Y + (\E Y)^2 \Big) \]
Групираме членовете:
\[ = \Big( \E X^2 - (\E X)^2 \Big) + \Big( \E Y^2 - (\E Y)^2 \Big) + 2\E(XY) - 2\E X\E Y \]
Първите две групи са точно $\Var X$ и $\Var Y$. Остава смесеният член: $2\big(\E(XY) - \E X\E Y\big)$.
Тъй като по условие $X$ и $Y$ са \textbf{независими}, знаем от свойствата на очакването, че $\E(XY) = \E X \cdot \E Y$. Следователно смесеният член се нулира. Така окончателно получаваме:
\[ \Var(X + Y) = \Var X + \Var Y \]
\end{proof}

Това свойство ни показва, че \emph{дисперсията е адитивна, но само за независими случайни величини}. Може да се направи геометрична аналогия: независимостта играе ролята на „ортогоналност“ на вектори. За ортогонални вектори квадратът от дължината на тяхната сума е равен на сумата от квадратите на техните дължини (подобно на Питагоровата теорема). Дисперсията е еквивалентът на този квадрат на дължината.

\section{Задачи}

В модела за групово тестване пренебрегнете биологичните ограничения и за $p=0{,}95$ намерете целия брой хора $n\geq 1$, който минимизира средния брой тестове на човек,
\[
1+\frac{1}{n}-p^n.
\]
Решете задачата числено, например чрез графика. Като алтернативен подход минимизирайте непрекъснатото продължение за $x\geq 1$ и сравнете стойностите при целите числа около получения минимум.


\chapter{Пораждащи функции и основни дискретни разпределения}
\section{Пораждащи функции}

Пораждащите функции представляват полезна трансформация на случайните величини. Подобно на архитектурна скица, която дава различен поглед върху една сграда в зависимост от ъгъла, пораждащата функция предоставя мощен аналитичен инструмент за изследване на свойствата на дискретните разпределения.

\begin{keydefn}
Нека $X : \Omega \rightarrow \mathbb{N} \cup \{0\}$ е целочислена случайна величина (т.е. възможните ѝ стойности са $0, 1, 2, \dots$). Тогава под \textbf{пораждаща функция} (probability generating function) на $X$ разбираме функцията $g_X(s)$, дефинирана като:
\[ g_X(s) := \E s^X = \sum_{n=0}^{\infty} s^n \mathbb{P}(X=n), \quad \text{за } s \in [-1, 1]. \]
\end{keydefn}

Този ред кодира цялото разпределение на $X$: коефициентът пред $s^n$ е точно $\mathbb{P}(X=n)$. Записът с безкраен ред е удобен и когато $X$ приема само краен брой стойности, защото тогава всички останали коефициенти просто са нула. Така на едно разпределение се съпоставя една аналитична функция, от която впоследствие могат да се извлекат отново отделните вероятности. Именно това е смисълът на аналогията със скиците на една сграда: различното представяне прави видими различни свойства на същия обект.

С помощта на пораждащата функция могат лесно да се пресмятат математическото очакване и дисперсията на случайната величина, както и да се възстановят нейните вероятности.

\begin{prop}
Нека $X$ е целочислена случайна величина с пораждаща функция $g_X(s)$. Тогава:
\begin{enumerate}
    \item[а)] $\E X = g_X'(1)$
    \item[б)] $\Var X = g_X''(1) + g_X'(1) - (g_X'(1))^2$
    \item[в)] $k! \mathbb{P}(X=k) = g_X^{(k)}(0)$, за $k \ge 0$.
\end{enumerate}
\end{prop}


\begin{proof}
а) По дефиниция $g_X(s) = \sum_{n=0}^{\infty} s^n \mathbb{P}(X=n)$. Диференцирайки по $s$, получаваме:
\[ g_X'(s) = \frac{d}{ds} \sum_{n=0}^{\infty} s^n \mathbb{P}(X=n) = \sum_{n=1}^{\infty} n s^{n-1} \mathbb{P}(X=n) = \E[X s^{X-1}]. \]
Замествайки $s=1$, намираме $g_X'(1) = \sum_{n=1}^{\infty} n \mathbb{P}(X=n) = \E X$.

б) Аналогично, втората производна е:
\[ g_X''(s) = \E[X(X-1)s^{X-2}]. \]
При $s=1$ получаваме $g_X''(1) = \E[X(X-1)] = \E X^2 - \E X$.
От формулата за дисперсия $\Var X = \E X^2 - (\E X)^2$, можем да изразим $\E X^2 = g_X''(1) + \E X = g_X''(1) + g_X'(1)$. Така:
\[ \Var X = g_X''(1) + g_X'(1) - (g_X'(1))^2. \]

в) За да намерим вероятността $X$ да приеме стойност $k$, диференцираме $g_X(s)$ $k$ пъти. Всички членове със степен по-малка от $k$ се анулират, а членът със степен $k$ става $k! \mathbb{P}(X=k)$. При заместване с $s=0$, всички останали членове (съдържащи $s$) се нулират, и остава точно $k! \mathbb{P}(X=k)$.
\end{proof}

\begin{cor}
Нека $X$ и $Y$ са целочислени случайни величини. Тогава $X \stackrel{d}{=} Y$ (равни по разпределение) тогава и само тогава, когато $g_X \equiv g_Y$.
\end{cor}

\begin{proof}
Ако $X\stackrel{d}{=}Y$, то $\mathbb{P}(X=n)=\mathbb{P}(Y=n)$ за всяко $n\ge0$, следователно съответните коефициенти в двата реда съвпадат и $g_X=g_Y$. Обратно, ако $g_X=g_Y$, то и производните им във всяка степен съвпадат. От вече доказаната формула
\[
k!\,\mathbb{P}(X=k)=g_X^{(k)}(0)=g_Y^{(k)}(0)=k!\,\mathbb{P}(Y=k)
\]
следва равенство на всички точкови вероятности, а значи и равенство по разпределение.
\end{proof}


Важно свойство на пораждащите функции се проявява при суми от независими случайни величини. 

\begin{defn}
Дискретните случайни величини $X_1, \dots, X_n$ (или $(X_j)_{j=1}^\infty$) имат \textbf{еднакво разпределение}, ако $X_i \stackrel{d}{=} X_j$ за всеки $i, j$.
\end{defn}


\begin{prop}
Нека $X_1, X_2, \dots, X_n$ са независими в съвкупност целочислени случайни величини. Ако $Y = \sum_{j=1}^n X_j$, то:
\[ g_Y(s) = \prod_{j=1}^n g_{X_j}(s). \]
Ако в допълнение $X_j$ са еднакво разпределени, то $g_Y(s) = (g_{X_1}(s))^n$.
\end{prop}

\begin{proof}
Ще го покажем за $n=2$. Нека $Y = X_1 + X_2$. Тогава:
\[ g_Y(s) = \E s^{X_1+X_2} = \sum_{j=0}^{\infty} \sum_{i=0}^{\infty} s^{i+j} \mathbb{P}(X_1=j, X_2=i). \]
Поради независимостта на $X_1$ и $X_2$, съвместната вероятност се разпада: $\mathbb{P}(X_1=j, X_2=i) = \mathbb{P}(X_1=j)\mathbb{P}(X_2=i)$.
\[ g_Y(s) = \sum_{j=0}^{\infty} s^j \mathbb{P}(X_1=j) \sum_{i=0}^{\infty} s^i \mathbb{P}(X_2=i) = g_{X_1}(s) g_{X_2}(s). \]
\end{proof}

Тук независимостта е съществена: тя позволява съвместната вероятност да се разложи на произведение. Ползата е, че прякото намиране на разпределението на една сума изисква да се обходят всички комбинации от стойности на събираемите, които дават един и същ резултат. При пораждащите функции тази комбинаторна задача се заменя с умножение на функции. Ако събираемите са и еднакво разпределени, произведението се свежда до една степен.

\begin{supp}[Пораждащата функция и функцията на разпределение]\label{supp:pgf-cdf}
Освен вероятностите и моментите, пораждащата функция кодира и функцията на разпределение:
за целочислена $X \ge 0$, при приетата тук строга конвенция $F_X(n) = \mathbb{P}(X < n)$,
\[ \sum_{n \ge 0} s^n F_X(n) = \frac{s\, g_X(s)}{1-s}, \qquad |s| < 1 . \]
Делението на $1-s$ съответства на натрупването: умножаването на реда по
$\frac{1}{1-s} = 1 + s + s^2 + \dots$ заменя всяка вероятност със сумата на всички преди
нея, което е точно преходът от $\mathbb{P}(X=n)$ към $F_X(n)$. Допълнителният множител
$s$ отчита строгото неравенство: той измества реда с една позиция, така че $\mathbb{P}(X=n)$
да не участва във $F_X(n)$. При широката конвенция $F_X(n) = \mathbb{P}(X \le n)$ множителят
отпада.
\end{supp}

\section{Основни дискретни разпределения}

Пораждащите функции ни дават единен апарат, с който да изведем очакването, дисперсията и вероятностите на изброените по-долу разпределения. Всички те се разглеждат върху схемата на Бернули.

\subsection{Схема на Бернули}

Схемата на Бернули е базов вероятностен модел, описващ поредица от независими експерименти с два възможни изхода (успех или неуспех) и постоянна вероятностна структура.
Двете изисквания трябва да се различават. Независимостта означава, че резултатът от един опит не променя вероятностите в останалите; еднаквото разпределение означава, че вероятността $p$ не се сменя от опит към опит. Например последователно хвърляне на различни монети може да даде независими резултати, но не и еднакво разпределени, ако монетите имат различни вероятности за ези.
Разглеждаме редица $(X_j)_{j=1}^{\infty}$, които са независими в съвкупност и еднакво разпределени:
\begin{align*}
\mathbb{P}(X_j=1) &= p \quad (\text{вероятност за успех}) \\
\mathbb{P}(X_j=0) &= 1-p = q \quad (\text{вероятност за неуспех})
\end{align*}

\medskip\noindent\textbf{Разпределение на Бернули.}
Случайната величина $X_1$ се нарича разпределена по Бернули, $X_1 \sim \Ber(p)$.
\[
\begin{array}{c|c|c}
X_1 & 0 & 1 \\ \hline
\mathbb{P} & q & p
\end{array}
\]
Пораждащата ѝ функция е:
\[ g_{X_1}(s) = q \cdot s^0 + p \cdot s^1 = q + ps. \]
Математическото очакване и дисперсията са:
\[ \E X_1 = p, \quad \Var X_1 = \E X_1^2 - (\E X_1)^2 = p - p^2 = pq. \]

\subsection{Биномно разпределение}

Биномното разпределение описва броя на успехите в първите $n$ експеримента от схема на Бернули. Записваме $X \sim \Bin(n, p)$.
Случайната величина $X$ може да се представи като сума от независими бернулиеви величини: $X = \sum_{j=1}^n X_j$.

\begin{prop}
Нека $X \sim \Bin(n, p)$. Тогава:
\begin{enumerate}
    \item[а)] $g_X(s) = (q+ps)^n$
    \item[б)] $\E X = np$
    \item[в)] $\Var X = npq$
    \item[г)] $\mathbb{P}(X=k) = \binom{n}{k} p^k q^{n-k}$, за $0 \le k \le n$.
\end{enumerate}
\end{prop}

\begin{lecfigure}[Биномното разпределение при $n=10$. Върхът стои около $\E X = np$:
при $p = 0{,}3$ това е $3$, при $p = 0{,}5$ — $5$. Симетрия има само при $p = \frac12$; за
$p < \frac12$ разпределението е изтеглено надясно, защото отляво го притиска границата
$k \ge 0$. Височините във всеки панел се сумират до
единица.\label{fig:binomial-pmf}]
\begin{tikzpicture}[x=0.52cm, y=9.5cm]
 \foreach \i/\lab/\dat in {%
   0/{$p = 0{,}3$}/{0/0.028,1/0.121,2/0.233,3/0.267,4/0.200,5/0.103,6/0.037,7/0.009,8/0.001,9/0.000,10/0.000},
   1/{$p = 0{,}5$}/{0/0.001,1/0.010,2/0.044,3/0.117,4/0.205,5/0.246,6/0.205,7/0.117,8/0.044,9/0.010,10/0.001}}
 {
  \begin{scope}[xshift=\i*7.2cm]
    \foreach \k/\h in \dat {
      \fill[defback] ({\k-0.36},0) rectangle ({\k+0.36},\h);
      \draw[defframe, line width=0.4pt] ({\k-0.36},0) rectangle ({\k+0.36},\h);}
    \draw[axisline] (-0.8,0) -- (10.9,0) node[below right, font=\small] {$k$};
    \foreach \k in {0,2,4,6,8,10} { \node[below, font=\scriptsize] at (\k,-0.004) {\k}; }
    \node[below, font=\small] at (5,-0.028) {\lab};
  \end{scope}}
\end{tikzpicture}
\end{lecfigure}

\begin{proof}
а) Тъй като $X = \sum_{j=1}^n X_j$ и $X_j$ са независими и еднакво разпределени, то $g_X(s) = (g_{X_1}(s))^n = (q+ps)^n$.

б) $\E X = g_X'(1) = np(q+ps)^{n-1} \big|_{s=1} = np(p+q)^{n-1} = np$.

в) Дисперсията се получава по същия начин чрез втората производна; сметката е оставена на читателя.

г) $\mathbb{P}(X=k)$ може да се намери чрез $k$-тата производна в нулата:
\[ k! \mathbb{P}(X=k) = g_X^{(k)}(0) = \left. \frac{d^k}{ds^k} (q+ps)^n \right|_{s=0} = n(n-1)\dots(n-k+1) p^k q^{n-k}. \]
Следователно $\mathbb{P}(X=k) = \frac{n(n-1)\dots(n-k+1)}{k!} p^k q^{n-k} = \binom{n}{k} p^k q^{n-k}$.
\end{proof}

Като приближен модел от лекцията беше разгледан потокът от хора, завръщащи се в страната в началото на пандемията. Ако за всеки пристигащ означим със стойност $1$ заразен човек и със стойност $0$ незаразен, приемем една и съща вероятност за заразяване $p$ и независимост между пристигащите, то всеки индикатор е бернулиев. Броят заразени сред първите $n$ пристигнали е сумата на тези индикатори и следователно има биномно разпределение. Лекторът изрично представи това като грубо приближение: еднаквото $p$ предполага сходен регионален риск, а независимостта изключва общ източник на заразяване.

\begin{supp}[Най-вероятна стойност (мода)]\label{supp:mode}
Естествен въпрос е къде разпределението достига максимума си. Отговорът се получава, като
сравним два съседни члена: разглеждаме отношението $\mathbb{P}(X=k)/\mathbb{P}(X=k-1)$ и
търсим къде то минава през $1$ — дотам вероятностите растат, след това намаляват.
За биномното разпределение
\[
\frac{\mathbb{P}(X=k)}{\mathbb{P}(X=k-1)} = \frac{n-k+1}{k}\cdot\frac{p}{q} \ \ge\ 1
\iff k \le p(n+1),
\]
откъдето модата е $k^* = \lfloor p(n+1) \rfloor$. По същия начин за Поасоновото
разпределение $k^* = \lfloor \lambda \rfloor$. Забележете, че модата и очакването са
близки, но по принцип различни числа.
\end{supp}

\subsection{Геометрично разпределение}

Геометричното разпределение моделира броя на неуспехите до настъпване на първия успех в схема на Бернули. Записваме $X \sim \Geo(p)$. Формално:
\[ X = \min\left\{n \ge 1 : \sum_{j=1}^n X_j = 1\right\} - 1. \]
В някои източници геометричното разпределение брои всички опити до първия успех и тогава не се изважда единица. В лекцията е избрана конвенцията за броя на неуспехите, така че възможните стойности започват от $0$. Например редицата $0,0,0,1$ дава $X=3$.

\begin{prop}
Нека $X \sim \Geo(p)$. Тогава:
\begin{enumerate}
    \item[а)] $\mathbb{P}(X=k) = q^k p$, за $k \ge 0$.
    \item[б)] $g_X(s) = \frac{p}{1-qs}$
    \item[в)] $\E X = \frac{q}{p}$, $\Var X = \frac{q}{p^2}$
\end{enumerate}
\end{prop}

\begin{proof}
а) Събитието $\{X=k\}$ означава, че първите $k$ експеримента са неуспешни, а $k+1$-вият е успешен. Поради независимостта:
\[ \mathbb{P}(X=k) = \mathbb{P}(X_1=0) \dots \mathbb{P}(X_k=0) \mathbb{P}(X_{k+1}=1) = q^k p. \]

б) За пораждащата функция имаме:
\[ g_X(s) = \sum_{n=0}^{\infty} s^n \mathbb{P}(X=n) = \sum_{n=0}^{\infty} s^n q^n p = p \sum_{n=0}^{\infty} (qs)^n = \frac{p}{1-qs}. \]

в) Чрез диференциране на $g_X(s)$ получаваме
\[
g_X'(s)=\frac{pq}{(1-qs)^2},
\qquad
g_X''(s)=\frac{2pq^2}{(1-qs)^3}.
\]
Следователно
\[
\E X=g_X'(1)=\frac{pq}{(1-q)^2}=\frac{q}{p}
\]
и, понеже $1-q=p$,
\begin{align*}
\Var X
  &=g_X''(1)+g_X'(1)-\bigl(g_X'(1)\bigr)^2 \\
  &=\frac{2q^2}{p^2}+\frac{q}{p}-\frac{q^2}{p^2}
    =\frac{q^2+pq}{p^2}
    =\frac{q}{p^2}.
\end{align*}
Така очакването и дисперсията се получават с пряка аналитична сметка от пораждащата функция, без да се пресмятат поотделно съответните степенни редове.
\end{proof}

Интерпретацията не зависи от това кой от двата изхода е наречен „успех“. При честна монета броят тури преди първото ези е геометричен с $p=1/2$ и има очакване $1$. В модела с вероятност за заразяване $p=0{,}05$ средният брой незаразени преди първия заразен е $0{,}95/0{,}05=19$. При баскетболист, който вкарва наказателен удар с вероятност $0{,}9$, броят вкарани кошове преди първия пропуск също е геометричен, но тук „успехът“, който прекратява броенето, е именно пропускът и неговата вероятност е $0{,}1$.

\medskip\noindent\textbf{Свойство безпаметност (Memoryless property).}
Геометричното разпределение е единственото дискретно разпределение, което „няма памет“. Тоест, фактът, че вече са настъпили определен брой неуспехи, не променя вероятността за бъдещите експерименти.

\begin{prop}[Безпаметност]
Ако $X \sim \Geo(p)$, то за $\forall m \ge 0$ и $\forall k \ge 0$:
\[ \mathbb{P}(X \ge m+k \given X \ge k) = \mathbb{P}(X \ge m). \]
\end{prop}

\begin{proof}
\[ \mathbb{P}(X \ge m+k \given X \ge k) = \frac{\mathbb{P}(X \ge m+k \text{ и } X \ge k)}{\mathbb{P}(X \ge k)} = \frac{\mathbb{P}(X \ge m+k)}{\mathbb{P}(X \ge k)}. \]
Тъй като $\mathbb{P}(X \ge k) = q^k$, получаваме $\frac{q^{m+k}}{q^k} = q^m = \mathbb{P}(X \ge m)$.
\end{proof}

Една интерпретация е животът на компонент: ако процесор вече е работил $k$ дни, при модел с безпаметност вероятността да работи още $m$ дни е същата като за току-що поставен процесор. Лекторът противопостави това на човешкия живот, при който възрастта очевидно променя оставащото време и геометричният модел не е подходящ.

Същото свойство обяснява защо историята на независими хазартни опити не помага за следващия залог. Ако три пъти поред се е паднало черно или тура, условната вероятност редицата да продължи още два пъти е същата, както без тази предходна информация. Записването на старите резултати не променя бъдещите вероятности; това важи и за миналите тиражи на „6 от 49“ при модела на независими тегления.

\subsection{Отрицателно биномно разпределение}

Отрицателното биномно разпределение обобщава геометричното. То брои броя на неуспехите до настъпването на $r$-тия успех в схема на Бернули. Означаваме го с $X \sim \NegBin(r, p)$, където $r \ge 1$ е цяло число.
Забележете, че $\NegBin(1, p) = \Geo(p)$.
Случайната величина може да се представи като сума от независими геометрични величини: $X = \sum_{j=1}^r Y_j$, където $Y_j \sim \Geo(p)$. В лекцията това е свързано с броя на пристигналите пътници преди да бъдат установени $r$ заразени: интересува ни не първият установен случай, а моментът, в който броят им достига капацитета, за който системата е подготвена.
Разлагането на сумата има конкретен смисъл: $Y_1$ брои неуспехите до първия успех, след него броенето започва отново и $Y_2$ брои неуспехите до втория успех, и така до $Y_r$. Поради независимостта на опитите тези блокове са независими и имат едно и също геометрично разпределение. Тази структура е причината очакването, дисперсията и пораждащата функция да се получават непосредствено от съответните величини за геометричното разпределение.

\begin{prop}
Нека $X \sim \NegBin(r, p)$. Тогава:
\begin{enumerate}
    \item[а)] $g_X(s) = \left( \frac{p}{1-qs} \right)^r$
    \item[б)] $\E X = r\frac{q}{p}$, $\Var X = r\frac{q}{p^2}$
    \item[в)] $\mathbb{P}(X=k) = \binom{r+k-1}{k} p^r q^k$, за $k \ge 0$.
\end{enumerate}
\end{prop}

\begin{proof}
От представянето като сума и независимостта на $Y_1,\dots,Y_r$ следва
\[
g_X(s)=\prod_{j=1}^r g_{Y_j}(s)
      =\left(\frac{p}{1-qs}\right)^r.
\]
Линейността на очакването и адитивността на дисперсията за независими величини дават
\[
\E X=\sum_{j=1}^r\E Y_j=r\frac{q}{p},
\qquad
\Var X=\sum_{j=1}^r\Var(Y_j)=r\frac{q}{p^2}.
\]

За вероятностите използваме $k$-тата производна на пораждащата функция:
\begin{align*}
k!\,\mathbb{P}(X=k)
  &=g_X^{(k)}(0) \\
  &=p^r q^k r(r+1)\dots(r+k-1).
\end{align*}
След деление на $k!$ получаваме
\[
\mathbb{P}(X=k)
  =\frac{r(r+1)\dots(r+k-1)}{k!}p^rq^k
  =\binom{r+k-1}{k}p^rq^k.
\]
Това е аналитичният извод от лекцията; пораждащата функция свежда комбинаторния коефициент до формално диференциране.
\end{proof}

\subsection{Поасоново разпределение}

Поасоновото разпределение често се използва за моделиране на събития, настъпващи рядко в непрекъснато време или пространство (напр. брой катастрофи, брой клиенти за единица време).
Като други начални модели бяха посочени броят голове в мач и броят видими звезди в участък от небето. Тук разпределението се задава направо чрез вероятностите си, а не като конкретен брой успехи в схема на Бернули. По-дълбокото му извеждане чрез процеси със стационарни и независими нараствания беше изрично оставено извън курса. Практическата връзка, която предстои, е друга: при подходящ режим Поасоновото разпределение приближава биномното и избягва тежките пресмятания с биномни коефициенти.

\begin{defn}
Случайната величина $X$ е разпределена по Поасон с параметър $\lambda > 0$, $X \sim \Poi(\lambda)$, ако вероятностите ѝ са дадени от:
\[ \mathbb{P}(X=k) = \frac{\lambda^k}{k!} e^{-\lambda}, \quad \text{за } k \ge 0. \]
\end{defn}

Коректността на това разпределение следва от развитието на експоненциалната функция в ред на Тейлър:
\[ \sum_{k=0}^{\infty} \mathbb{P}(X=k) = e^{-\lambda} \sum_{k=0}^{\infty} \frac{\lambda^k}{k!} = e^{-\lambda} e^{\lambda} = 1. \]

\begin{prop}
Нека $X \sim \Poi(\lambda)$. Тогава $g_X(s) = e^{\lambda s - \lambda}$, и $\E X = \Var X = \lambda$.
\end{prop}

\begin{proof}
За пораждащата функция имаме:
\[ g_X(s) = \sum_{n=0}^{\infty} s^n \frac{\lambda^n}{n!} e^{-\lambda} = e^{-\lambda} \sum_{n=0}^{\infty} \frac{(s\lambda)^n}{n!} = e^{-\lambda} e^{\lambda s} = e^{\lambda s - \lambda}. \]
Диференцирайки, получаваме $g_X'(s) = \lambda e^{\lambda s - \lambda}$, което при $s=1$ дава $\E X = \lambda$.
Втората производна е $g_X''(s) = \lambda^2 e^{\lambda s - \lambda}$, откъдето $\Var X = \lambda^2 + \lambda - (\lambda)^2 = \lambda$.
\end{proof}

\section{Задачи}

\begin{enumerate}
    \item Доказателството на формулата
    \[
    g_{X_1+\dots+X_n}(s)=\prod_{j=1}^n g_{X_j}(s)
    \]
    беше направено за две независими случайни величини. Обобщете го за произволен краен брой събираеми.
    \item За $X\sim\Bin(n,p)$ довършете оставеното пресмятане с втората производна на $g_X(s)=(q+ps)^n$ и изведете $\Var X=npq$.
    \item За геометрично разпределена величина пресметнете $\E X$ направо от дефиниционния ред, без да използвате пораждащата функция.
    \item За $X\sim\Poi(\lambda)$ намерете $\E X$ и $\Var X$ чрез първите две производни на $g_X(s)=e^{\lambda(s-1)}$.
\end{enumerate}


\chapter{Теорема на Поасон. Хипергеометрично разпределение, ковариация и корелация}
\section{Гранични теореми и приближения}

\subsection{Теорема на Поасон (Поасоново приближение)}

Досега разглеждахме Поасоновата случайна величина като модел, дефиниран чрез своите вероятности. За разлика от биномното разпределение, при което имаме ясен експериментален модел (хвърляне на монета или повтаряеми опити с два изхода), Поасоновото разпределение често възниква като граничен случай или като добро приближение на други разпределения.

Преди да формулираме основната теорема, нека припомним едно много полезно твърдение, свързано с пораждащите функции. Нека имаме редица от целочислени неотрицателни случайни величини $X_n$ и съответните им пораждащи функции $G_{X_n}(s)$. Нека $X$ е също целочислена неотрицателна случайна величина с пораждаща функция $G_X(s)$. Ако за всяко $s$ от дефиниционното множество имаме сходимост на пораждащите функции $\lim_{n \to \infty} G_{X_n}(s) = G_X(s)$, то за всяко цяло $k \ge 0$ е изпълнено и сходимост на самите вероятности:
\[
\lim_{n \to \infty} \mathbb{P}(X_n = k) = \mathbb{P}(X = k)
\]
Това твърдение ни позволява чрез граница на функции да идентифицираме граница на вероятности, което много често е технически значително по-лесно от директното доказване на сходимост на самите вероятности.

Ще използваме този факт, за да докажем една от важните ранни теореми в теорията на вероятностите, а именно Теоремата на Поасон. Тя илюстрира как един теоретичен резултат може да ни даде основание за полезни практически приближения.

\begin{keythm}[на Поасон]
Нека за всяко $n \ge 1$ е дадена случайна величина $X_n$, която има биномно разпределение $X_n \sim \Bin(n, p_n)$. Т.е. имаме редица от биномни случайни величини с нарастващ брой експерименти $n$ и променяща се вероятност за успех $p_n$. Нека освен това границата на произведението $n p_n$ е положително число:
\[
\lim_{n \to \infty} n p_n = \lambda > 0
\]
Тогава за всяко $k \ge 0$, границата на вероятностите $X_n$ да е равно на $k$ е равна на вероятността Поасонова случайна величина с параметър $\lambda$ да е равна на $k$:
\[
\lim_{n \to \infty} \mathbb{P}(X_n = k) = \frac{\lambda^k}{k!} e^{-\lambda}
\]
Грубо казано, ако броят експерименти по вероятността за успех се стабилизира до някакво число $\lambda$, сложните биномни вероятности се схождат към много по-лесно изчислимите Поасонови вероятности.
\end{keythm}

\begin{proof}[Доказателство (за частен случай)]
Няма да доказваме теоремата в пълната ѝ общност. За да предадем основната интуиция, ще разгледаме по-простия случай, при който не просто границата е $\lambda$, а самото произведение е точно равно на $\lambda$ за всяко $n$. Т.е. нека $\lambda = n p_n$, откъдето $p_n = \frac{\lambda}{n}$.

Пораждащата функция на биномната случайна величина $X_n$ е:
\[
G_{X_n}(s) = \left(1 - p_n + p_n s\right)^n
\]
Замествайки $p_n = \frac{\lambda}{n}$, получаваме:
\[
G_{X_n}(s) = \left(1 - \frac{\lambda}{n} + \frac{\lambda}{n} s\right)^n = \left(1 + \frac{\lambda (s-1)}{n}\right)^n
\]
Като използваме основното свойство на експоненциалната функция като граница, намираме:
\[
\lim_{n \to \infty} G_{X_n}(s) = \lim_{n \to \infty} \left(1 + \frac{\lambda (s-1)}{n}\right)^n = e^{\lambda (s-1)}
\]
Получената граница $e^{\lambda (s-1)}$ е точно пораждащата функция на Поасонова случайна величина $X \sim \Poi(\lambda)$. Благодарение на споменатото по-горе твърдение, това доказва, че $\lim_{n \to \infty} \mathbb{P}(X_n = k) = \frac{\lambda^k}{k!} e^{-\lambda}$.

\textbf{Практическо приложение (Приближение на Поасон):}
В практиката тази теорема не се използва толкова като строга граница (понеже на практика разполагаме с едно конкретно $n$), колкото като рецепта за приближение. Сложният израз за биномната вероятност $\binom{n}{k} p^k (1-p)^{n-k}$ изисква пресмятането на големи факториели и високи степени.

Ако разполагаме със случайна величина $Y \sim \Bin(n, p)$, където $n$ е достатъчно голямо, а $p$ е достатъчно малко, можем да използваме приближението:
\[
\mathbb{P}(Y = k) \approx \frac{\lambda^k}{k!} e^{-\lambda}, \quad \text{където } \lambda = n p
\]
\emph{Рецепта:} Приближението работи добре и може спокойно да се използва в случаите, когато $n \ge 100$ и $np \le 20$. Тогава, дори и да не знаем „границите“, твърдим, че биномните вероятности са изключително близки до Поасоновите с параметър $\lambda = np$.
\end{proof}


\section{Хипергеометрично разпределение}

Последното конкретно разпределение от фундаменталните дискретни случайни величини, което ще разгледаме, е хипергеометричното разпределение. За разлика от Поасоновото, то е директно обвързано с един много ясен физически и комбинаторен модел на извличане без връщане.

\begin{defn}[хипергеометричен модел]
Казваме, че $X$ има хипергеометрично разпределение с параметри $N$, $M$ и $n$ (записваме $X \sim \Hyp(N, M, n)$), ако $X$ описва следния процес:
Имаме общо $N$ на брой обекта (например продукти в партида). От тях точно $M$ са „маркирани“ (или „дефектни“), а останалите $N - M$ са немаркирани (или здрави). Ние изтегляме случайно $n$ на брой обекта от тези $N$ без връщане. Случайната величина $X$ представлява \emph{броят на изтеглените маркирани обекти} измежду тези $n$ обекта.
\end{defn}

Естествените параметрични ограничения са $N \ge M$ и $N \ge n$. Това разпределение намира огромно приложение в качествения контрол, където искаме да оценим каква пропорция от огромна партида стоки са дефектни, използвайки само малка извадка от тях.

\begin{prop}
Нека $X \sim \Hyp(N, M, n)$. Тогава:
\begin{enumerate}
    \item[а)] Вероятностното разпределение на $X$ се задава с:
    \[
    \mathbb{P}(X = k) = \frac{\binom{M}{k} \binom{N-M}{n-k}}{\binom{N}{n}}
    \]
    където $k$ удовлетворява естествените ограничения $0 \lor (n - N + M) \le k \le M \land n$.
    \item[б)] Математическото очакване е:
    \[
    \E X = n \frac{M}{N}
    \]
    \item[в)] Дисперсията е:
    \[
    \Var X = n \frac{M}{N} \frac{N-M}{N} \frac{N-n}{N-1}
    \]
\end{enumerate}
\end{prop}

\begin{proof}
\emph{За а):} Формулата следва от класическата вероятност. Всички възможни начини да изтеглим $n$ обекта от $N$ са $\binom{N}{n}$. За да изтеглим точно $k$ маркирани, трябва да изберем $k$ обекта от наличните $M$ маркирани (това става по $\binom{M}{k}$ начина) и останалите $n-k$ обекта да изберем от наличните $N-M$ немаркирани (това става по $\binom{N-M}{n-k}$ начина). Границите за $k$ гарантират, че не избираме повече маркирани, отколкото съществуват или отколкото теглим ($k \le \min(M, n)$), както и че няма как да изтеглим отрицателен брой обекти или повече немаркирани, отколкото съществуват.

\emph{За б):} Това свойство има скрита линейна структура, която изключително улеснява пресмятането на очакването. Дефинираме индикаторни случайни величини $X_j$ за всяко $j \in \{1, \dots, n\}$ по следния начин:
\[
X_j = \begin{cases}
1, & \text{ако } j\text{-тият изтеглен обект е маркиран} \\
0, & \text{ако не е маркиран}
\end{cases}
\]
Ясно е, че общият брой маркирани обекти $X$ е просто сумата на тези индикатори:
\[
X = \sum_{j=1}^n X_j
\]
Тъй като математическото очакване винаги е линейно (независимо дали събираемите са независими или не), можем да запишем:
\[
\E X = \E\left(\sum_{j=1}^n X_j\right) = \sum_{j=1}^n \E X_j
\]
Всяко $X_j$ е Бернулиева случайна величина $X_j \sim \Ber(p_j)$, където $p_j$ е вероятността $j$-тият изтеглен обект да е маркиран. Поради пълната симетрия на модела (няма значение дали теглите първи или $j$-ти, преди да сте видели резултатите от останалите тегления, вероятността е една и съща), вероятността на коя да е позиция да се падне маркиран обект е точно пропорцията на маркираните обекти в началото, а именно $p_j = \frac{M}{N}$.
Следователно, $\E X_j = p_j = \frac{M}{N}$. Замествайки обратно в сумата, получаваме:
\[
\E X = \sum_{j=1}^n \frac{M}{N} = n \frac{M}{N}
\]
\emph{За в):} Тук вече не можем да съберем дисперсиите наум, защото случайните величини $X_j$ \emph{не са независими} (изтеглянето на маркиран обект на първо теглене променя вероятностите за следващите тегления, тъй като тегленето е без връщане). Затова дисперсията на сумата не е просто сума от дисперсиите, а включва и ковариационните членове. Формулата се дава наготово.

\end{proof}

\begin{example}[Цветни маркери]\label{ex:07-1}
Имаме 8 маркера, от които 4 са цветни и 4 са черни. Т.е. $N = 8, M = 4$. Теглим $n = 3$ маркера на случаен принцип. Броят цветни маркери, които сме изтеглили сред тези три, се описва с хипергеометричната случайна величина $X \sim \Hyp(8, 4, 3)$.
\end{example}


\section{Съвместни разпределения}

Досега разглеждахме случайните величини изолирано. Често пъти обаче искаме да описваме сложни модели с повече от една случайна характеристика. Например, искаме да изследваме както разхода на един клиент в магазина, така и неговия пол. Затова трябва да въведем понятия за съвместна работа с две или повече случайни величини. За удобство на записа ще разглеждаме двойки $(X, Y)$, но концепциите се обобщават за произволен краен брой.

\subsection{Съвместни дискретни разпределения}

\begin{defn}
Нека $X$ и $Y$ са две дискретни случайни величини, дефинирани върху едно и също вероятностно пространство. Под \textbf{съвместно разпределение} на $X$ и $Y$ разбираме съвкупността от всички вероятности:
\[
p_{i,j} = \mathbb{P}(X = x_i, Y = y_j)
\]
за всички възможни стойности $x_i$ на $X$ и $y_j$ на $Y$. Обикновено това се представя във формата на таблица, където редовете отговарят на възможните стойности на $X$, а колоните — на възможните стойности на $Y$. Всяка клетка $(i, j)$ съдържа вероятността двете събития да настъпят едновременно. Сумата на всички числа (клетки) в тази таблица е равна на 1.
\end{defn}

\textbf{Маргинални разпределения:}
Ако съберем всички вероятности по даден ред (фиксирайки $X = x_i$ и сумирайки по всички възможни стойности на $Y$), по формулата за пълната вероятност получаваме \emph{маргиналното разпределение} на $X$:
\[
\mathbb{P}(X = x_i) = \sum_{j} \mathbb{P}(X = x_i, Y = y_j) = \sum_{j} p_{i,j}
\]
Това е просто индивидуалното разпределение на $X$, такова каквото би било, ако въобще не се интересувахме от $Y$. Аналогично, сумирайки по колони, получаваме маргиналното разпределение на $Y$:
\[
\mathbb{P}(Y = y_j) = \sum_{i} \mathbb{P}(X = x_i, Y = y_j) = \sum_{i} p_{i,j}
\]

\begin{example}\label{ex:07-2}
Хвърляме два стандартни зара (номерирани от 1 до 6). Дефинираме две случайни величини:
\begin{itemize}
    \item $X$: броят на хвърлените шестици (възможни стойности: 0, 1, 2);
    \item $Y$: броят на хвърлените единици (възможни стойности: 0, 1, 2).
\end{itemize}

Нека попълним таблицата на съвместното им разпределение. Общият брой възможни изходи при хвърляне на два зара е $6 \times 6 = 36$.
\begin{itemize}
    \item $\mathbb{P}(X=0, Y=0)$: Търсим вероятността да няма нито една шестица и нито една единица. Значи и двата зара трябва да покажат числа между 2 и 5 (общо 4 възможности за всеки зар). Вероятността е $\frac{4 \times 4}{36} = \frac{16}{36}$.
    \item $\mathbb{P}(X=1, Y=0)$: Търсим една шестица и нула единици. За първия зар имаме шестица, а за втория число от 2 до 5 (4 възможности). Или обратното. Общо $2 \times 4 = 8$ възможности. Вероятността е $\frac{8}{36}$.
    \item $\mathbb{P}(X=0, Y=1)$: По симетрия (една единица и нула шестици), вероятността също е $\frac{8}{36}$.
    \item $\mathbb{P}(X=1, Y=1)$: Една шестица и една единица. Вариантите са (6,1) и (1,6). Вероятността е $\frac{2}{36}$.
    \item $\mathbb{P}(X=2, Y=0)$: Две шестици, вероятност $\frac{1}{36}$.
    \item $\mathbb{P}(X=0, Y=2)$: Две единици, вероятност $\frac{1}{36}$.
    \item Всички останали комбинации (напр. $X=2, Y=1$) са невъзможни (не може да имаме общо повече от 2 зара), така че вероятността им е 0.
\end{itemize}

Сумата от всички вероятности в таблицата е $\frac{16+8+8+2+1+1}{36} = \frac{36}{36} = 1$.

Ако изчислим маргиналното разпределение на $X$ чрез сумиране по редовете, получаваме:
\[
\mathbb{P}(X=0) = \frac{16+8+1}{36} = \frac{25}{36}
\]
\[
\mathbb{P}(X=1) = \frac{8+2}{36} = \frac{10}{36}
\]
\[
\mathbb{P}(X=2) = \frac{1}{36}
\]
Което е точно очакваното биномно разпределение за броя на успехите (падане на шестица, $p = 1/6$) при 2 опита.


\end{example}

\subsection{Съвместна функция на разпределение (Joint CDF)}

\begin{defn}
Нека $X$ и $Y$ са две произволни случайни величини (дискретни или непрекъснати). Тогава \textbf{съвместна функция на разпределение} (или кумулативна функция) на $X$ и $Y$ е функцията $F_{X,Y} : \mathbb{R}^2 \to [0, 1]$, зададена като:
\[
F_{X,Y}(x,y) := \mathbb{P}(X < x, Y < y) \quad \forall (x,y) \in \mathbb{R}^2
\]
Геометрично, ако разглеждаме стойностите на $X$ и $Y$ като координати на двумерен вектор в равнината, $F_{X,Y}(x,y)$ е вероятността този случаен вектор да попадне в безкрайния правоъгълник на югозапад от точката $(x, y)$ (строго по-малко от $x$ по хоризонтала и строго по-малко от $y$ по вертикала).
\end{defn}

\begin{prop}[Маргинални функции от съвместната]\label{prop:marginal-from-joint}
От съвместната функция на разпределение се възстановяват маргиналните:
\[
F_{X,Y}(x, \infty) = F_X(x), \qquad F_{X,Y}(\infty, y) = F_Y(y), \qquad F_{X,Y}(\infty, \infty) = 1 .
\]
\end{prop}

Възстановяването става, като заместим съответно другата координата с безкрайност. Понеже събитието $\{Y < \infty\}$ е достоверно (вероятност 1), то сечението на $\{X < x\}$ с достоверно събитие не променя вероятността:
\[
F_{X,Y}(x, \infty) = \mathbb{P}(X < x, Y < \infty) = \mathbb{P}(X < x) = F_X(x)
\]
Аналогично:
\[
F_{X,Y}(\infty, y) = F_Y(y)
\]
Очевидно е също, че $F_{X,Y}(\infty, \infty) = \mathbb{P}(X < \infty, Y < \infty) = 1$.

\begin{keythm}[Критерий за независимост]\label{thm:indep-criterion}
Две случайни величини $X$ и $Y$ са независими тогава и само тогава, когато съвместната им функция на разпределение се разпада на произведение от маргиналните им функции на разпределение за всяка точка $(x,y)$:
\[
F_{X,Y}(x, y) = F_X(x) F_Y(y) \quad \forall (x,y) \in \mathbb{R}^2 .
\]
\end{keythm}

Това е най-общият критерий за независимост. За дискретни разпределения обикновено се ползва еквивалентният и по-интуитивен критерий с вероятности в конкретни точки: $\mathbb{P}(X = x, Y = y) = \mathbb{P}(X = x) \mathbb{P}(Y = y)$.

\begin{remark}[Сума на независими величини и зависимост]\label{rem:sum-dependence}
Ако случайните величини $X$ и $Y$ са независими и стандартно нормално разпределени ($N(0,1)$), а $Z$ зависи от $Y$, следва ли, че $X+Y$ е независимо от $Z$?

Не — и контрапримерът е елементарен. Нека $Y = Z$. Тогава сумата $X+Y$ е всъщност $X+Z$, което очевидно зависи директно от $Z$. Може да има специфични случаи на независимост, но в общия случай $X+Y$ и $Z$ са зависими.
\end{remark}


\section{Ковариация и корелация}

След като дефинирахме какво означава две случайни величини да са независими във вероятностен смисъл, преминаваме към изследването на конкретни видове зависимости. Понякога знаем, че две величини са зависими, но искаме да измерим силата и посоката на тяхната \emph{линейна} зависимост. Търсим да отговорим на въпроса доколко $Y$ може да се представи като линейна функция на $X$, т.е. $Y \approx aX + b$.

\subsection{Ковариация}

Първата стъпка към дефиниране на такава мярка е ковариацията. За да изолираме отместванията на разпределенията по реалната права, ние първо ги центрираме, като изваждаме техните математически очаквания.
Ако дефинираме центрирани случайни величини $\widetilde{X} = X - \E X$ и $\widetilde{Y} = Y - \E Y$, то е очевидно, че тяхното очакване е 0 ($\E\widetilde{X} = 0$).

\begin{defn}
Ковариация на $X$ и $Y$ се нарича математическото очакване на произведението на техните центрирани версии (ако това очакване съществува):
\[
\Cov(X, Y) = \E[(X - \E X)(Y - \E Y)] = \E[\widetilde{X}\widetilde{Y}]
\]
Ковариацията може да се изчислява и по алтернативна формула, която се извежда директно чрез разкриване на скобите и използване на линейността на математическото очакване:
\[
\Cov(X, Y) = \E[XY] - \E X \E Y
\]
\end{defn}
\begin{cor}
Ако $X$ и $Y$ са независими (най-силното свойство на липса на зависимост), то математическото очакване на произведението им се разпада на произведение от очакванията им ($\E[XY] = \E X\E Y$). Следователно, при независими случайни величини ковариацията е точно 0. Важно е да се отбележи, че обратното не е вярно — нулева ковариация означава само липса на линейна зависимост, но величините може да са свързани чрез по-сложна, нелинейна функция.
\end{cor}

\begin{prop}[Дисперсия на сума]\label{prop:var-sum}
За произволни случайни величини с крайни дисперсии
\[
\Var(aX + bY) = a^2 \Var X + b^2 \Var Y + 2ab\,\Cov(X,Y),
\]
и по-общо
\[
\Var\Big( \sum_{j=1}^{n} X_j \Big)
 = \sum_{j=1}^{n} \Var X_j + 2 \sum_{i<j} \Cov(X_i, X_j) .
\]
\end{prop}

\begin{proof}
Достатъчно е да разкрием квадрата под очакването. Полагаме $\tilde X = X - \E X$ и
$\tilde Y = Y - \E Y$; тогава $a X + b Y$ има центрирана част $a\tilde X + b\tilde Y$ и
\[
\Var(aX+bY) = \E\big[(a\tilde X + b\tilde Y)^2\big]
 = a^2 \E \tilde X^2 + b^2 \E \tilde Y^2 + 2ab\,\E[\tilde X \tilde Y],
\]
което е точно търсеното. Общият случай се получава по същия начин: при повдигане на
$\sum_j \tilde X_j$ на квадрат диагоналните членове дават дисперсиите, а всяка
недиагонална двойка се появява два пъти и дава $2\Cov(X_i, X_j)$.
\end{proof}

Ако величините са независими в съвкупност, всички ковариации се нулират и се връщаме към
адитивността на дисперсията от лекция 5. Точно тези „ковариационни членове“ липсват при
хипергеометричното разпределение, където изтеглянията не са независими.

\emph{Недостатък на ковариацията:} Ако умножим случайните величини по константа, например 10 пъти, тяхната ковариация ще се увеличи 100 пъти: $\Cov(10X, 10Y) = 100 \Cov(X, Y)$. Това математическо мащабиране променя стойността на мярката, въпреки че линейната структура на зависимостта между величините реално не се е променила.

\subsection{Коефициент на корелация}

За да компенсираме ефекта от мащабирането, ние не само центрираме случайните величини, но ги и \emph{нормираме}, разделяйки ги на корен квадратен от техните дисперсии (т.е. на стандартните им отклонения).

\begin{defn}
Коефициент на корелация $\rho(X, Y)$ между две случайни величини с крайни дисперсии $\Var X$ и $\Var Y$ се дефинира като:
\[
\rho(X, Y) = \frac{\Cov(X, Y)}{\sqrt{\Var X} \sqrt{\Var Y}}
\]
Ако дефинираме центрираните и нормирани случайни величини като:
\[
\overline{X} = \frac{X - \E X}{\sqrt{\Var X}}, \quad \overline{Y} = \frac{Y - \E Y}{\sqrt{\Var Y}}
\]
може лесно да се провери, че $\E\overline{X} = 0$ и $\Var\overline{X} = 1$. В тези термини, коефициентът на корелация е просто очакването на тяхното произведение:
\[
\rho(X, Y) = \E[\overline{X}\overline{Y}]
\]
Коефициентът на корелация е безразмерна величина. Ако пресметнем корелацията на $(10X, 10Y)$, множителят $100$ се появява както в числителя (ковариацията), така и в знаменателя (от корените на дисперсиите), и се съкращава. Следователно мащабирането не променя корелацията.
\end{defn}

\begin{keythm}[Граници на корелацията]\label{thm:corr-bounds}
За произволни случайни величини с крайни дисперсии е вярно, че:
\[
|\rho(X, Y)| \le 1 .
\]
\end{keythm}

Това е вероятностната форма на \emph{неравенството на Коши\,--\,Буняковски} (в
литературата и като неравенство на Коши\,--\,Шварц). Връзката с геометрията е
пряка: ковариацията играе ролята на скаларно произведение, стандартните
отклонения — на дължини, а коефициентът на корелация — на косинус от ъгъла между
двата вектора, откъдето и границите $\pm 1$.
\begin{proof}
Ще използваме очевидния факт, че квадратът на всяка случайна величина е неотрицателен, следователно и очакването му е неотрицателно:
\[
0 \le \E[(\overline{X} \pm \overline{Y})^2]
\]
Разкриваме скобите, ползвайки линейността на очакването:
\[
0 \le \E[\overline{X}^2 \pm 2\overline{X}\overline{Y} + \overline{Y}^2] = \E\overline{X}^2 \pm 2\E[\overline{X}\overline{Y}] + \E\overline{Y}^2
\]
Тъй като $\overline{X}$ и $\overline{Y}$ имат нулеви очаквания, техните дисперсии съвпадат с очакването на квадратите им: $\E\overline{X}^2 = \Var\overline{X} = 1$ и $\E\overline{Y}^2 = \Var\overline{Y} = 1$. Заместваме това заедно с дефиницията за $\rho(X, Y) = \E[\overline{X}\overline{Y}]$ и получаваме:
\[
0 \le 1 \pm 2\rho(X, Y) + 1 = 2 \pm 2\rho(X, Y)
\]
Следователно $1 \pm \rho(X, Y) \ge 0$, откъдето правим извода, че $|\rho(X, Y)| \le 1$.
\end{proof}

Следващата теорема обяснява точно какво ни казва коефициентът на корелация и кога той достига своите екстремални стойности $\pm 1$.

\begin{keythm}[Критерий за линейна зависимост]
Нека $X$ и $Y$ са случайни величини с крайни дисперсии. Тогава между тях има строго линейна зависимост от вида $Y = aX + b$ (за някакви константи $a \neq 0$ и $b$) тогава и само тогава, когато $|\rho(X, Y)| = 1$.
\end{keythm}

\begin{proof}[Доказателство ($\implies$)]
Нека е дадено, че $Y = aX + b$. Ще изразим връзката между нормираните величини $\overline{Y}$ и $\overline{X}$.
\[
\overline{Y} = \frac{Y - \E Y}{\sqrt{\Var Y}} = \frac{aX + b - (a\E X + b)}{\sqrt{\Var Y}} = \frac{a}{\sqrt{\Var Y}}(X - \E X)
\]
Тъй като $\sqrt{\Var X} = \frac{\sqrt{\Var Y}}{|a|}$, следва че:
\[
\overline{Y} = \frac{a}{|a|} \frac{X - \E X}{\sqrt{\Var X}} = v \overline{X}
\]
където $v = \text{sgn}(a) = \pm 1$. Така получихме връзката $\overline{Y} = \pm \overline{X}$.
Сега пресмятаме коефициента на корелация:
\[
\rho(X, Y) = \E[\overline{X}\overline{Y}] = \E[\overline{X}(\pm \overline{X})] = \pm \E\overline{X}^2 = \pm 1
\]
Тъй като $\E\overline{X}^2 = \Var\overline{X} = 1$. Това показва, че $|\rho(X, Y)| = 1$.
\end{proof}

\begin{proof}[Доказателство ($\Longleftarrow$)]
Нека сега е дадено, че $|\rho(X, Y)| = 1$. Ще разгледаме случая $\rho(X, Y) = 1$ (доказателството за $\rho = -1$ е аналогично).
От предишното доказателство знаем, че:
\[
\E[(\overline{X} - \overline{Y})^2] = 2 - 2\rho(X, Y)
\]
При $\rho(X, Y) = 1$, получаваме $\E[(\overline{X} - \overline{Y})^2] = 0$.
Тъй като интегрираме (търсим математическо очакване) на строго неотрицателна случайна величина и получаваме 0, това означава, че самата случайна величина е тъждествено равна на нула (с вероятност 1):
\[
(\overline{X} - \overline{Y})^2 = 0 \implies \overline{X} = \overline{Y}
\]
Заместваме дефинициите за центрирани и нормирани величини:
\[
\frac{X - \E X}{\sqrt{\Var X}} = \frac{Y - \E Y}{\sqrt{\Var Y}}
\]
От това уравнение лесно можем да изразим $Y$ като линейна функция на $X$:
\[
Y = \frac{\sqrt{\Var Y}}{\sqrt{\Var X}} X - \frac{\sqrt{\Var Y}}{\sqrt{\Var X}} \E X + \E Y
\]
Което е точно от вида $Y = aX + b$, където $a = \frac{\sqrt{\Var Y}}{\sqrt{\Var X}}$ и $b = \E Y - a\E X$. С това теоремата е доказана.
\end{proof}

В заключение, коефициент на корелация близо до 1 (или -1) индикира силна права (или обратна) линейна зависимост между величините. Коефициент близо до 0 означава липса на линейна зависимост.

\section{Задачи}

\begin{enumerate}
    \item За фиксирано $k$ докажете директно границата от теоремата на Поасон:
    \[
    \lim_{n\to\infty}\binom{n}{k}\left(\frac{\lambda}{n}\right)^k
    \left(1-\frac{\lambda}{n}\right)^{n-k}
    =e^{-\lambda}\frac{\lambda^k}{k!}.
    \]
    \item За $X\sim\Hyp(N,M,n)$ докажете оставената като техническо упражнение формула
    \[
    \Var X=n\frac{M}{N}\frac{N-M}{N}\frac{N-n}{N-1}.
    \]
    \item В примера с осемте маркера, четири от които са цветни, се теглят три маркера без връщане. Пресметнете вероятностите да бъдат изтеглени съответно $0$, $1$, $2$ и $3$ цветни маркера.
    \item От дефиницията на съвместната функция на разпределение проверете
    \[
    F_{X,Y}(x,\infty)=F_X(x),\qquad
    F_{X,Y}(\infty,y)=F_Y(y).
    \]
    \item Разкрийте скобите в дефиницията на ковариацията и докажете
    \[
    \Cov(X,Y)=\E[XY]-\E X\,\E Y.
    \]
    \item[Домашно 2.] Нека $X$ и $Y$ са независими стандартно нормално разпределени случайни величини, а $Z$ зависи от $Y$. Следва ли, че $X+Y$ е независимо от $Z$? Разгледайте като контрапример случая $Y=Z$ и разпишете зависимостта. Решението е обсъдено в забележката към раздела за независимост.
\end{enumerate}


\chapter{Условно математическо очакване. Непрекъснати случайни величини}
\section{Условно математическо очакване}

\subsection{Мотивация и средноквадратично приближение}
В много практически ситуации разполагаме със случайна величина $X$, която описва дадена числова характеристика (например приходи, разходи на клиент, време на работа на процесор и т.н.). Вече знаем, че обикновеното математическо очакване $\E X$ е числото, което най-добре характеризира $X$ в средноквадратичен смисъл. Това означава, че дисперсията $\Var X = \E(X - \E X)^2$ е минималното средноквадратично отклонение от произволна константа $a$:
\[
\Var X = \min_{a \in \mathbb{R}} \E(X - a)^2
\]
Следователно, очакването е най-добрата константа за оценка на $X$, ако не разполагаме с никаква допълнителна информация.

Какво обаче се случва, ако разполагаме с допълнителна, наблюдаема случайна величина $Y$, която носи информация за $X$? Например, ако $X$ е разходът на клиент, а $Y$ е неговият пол (мъж или жена), е логично да очакваме, че оценката ни за $X$ може да се подобри, ако се съобразим с наблюдаваната стойност на $Y$. В такъв случай вече не търсим просто едно число $a$, което да минимизира средноквадратичната грешка, а функция $G(Y)$, зависеща от наблюдаваната стойност на $Y$. Целта ни е да намерим функция $G$, за която се достига следният минимум:
\[
\min_{G} \E(X - G(Y))^2
\]

Нека разгледаме прост модел, в който $Y$ е бинарна случайна величина, която разбива пространството от елементарни събития $\Omega$ на две части — събитието $A$ (например клиентът е мъж) и неговото допълнение $\overline{A}$ или $A^c$ (клиентът е жена). Можем да запишем $Y$ чрез индикаторната функция $\ind_A$:
\[
Y = \ind_A = \begin{cases}
1, & \text{ако се сбъдне } A \\
0, & \text{ако се сбъдне } A^c
\end{cases}
\]
Нека вероятностите са съответно $\mathbb{P}(Y=1) = \mathbb{P}(A) = p$ и $\mathbb{P}(Y=0) = \mathbb{P}(A^c) = 1-p = q$. Тъй като $Y$ приема само две стойности, всяка възможна функция $G(Y)$ от нея е изключително проста и може да бъде записана като линейна комбинация от индикаторите на $A$ и $A^c$ с неизвестни константи $a, b \in \mathbb{R}$:
\[
G(Y) = a \ind_A + b \ind_{A^c}
\]
Тогава оптимизационната ни задача се свежда до минимизиране на функцията на две променливи $f(a, b)$:
\[
\min_{a, b \in \mathbb{R}} \underbrace{\E(X - a \ind_A - b \ind_{A^c})^2}_{f(a,b)} = \min_{a, b \in \mathbb{R}} f(a,b)
\]
Разкриваме скобите, като използваме свойствата на индикаторите, че $\ind_A^2 = \ind_A$ и че тяхното произведение $\ind_A \ind_{A^c} = 0$:
\[
f(a,b) = \E[X^2 + a^2 \ind_A + b^2 \ind_{A^c} - 2aX\ind_A - 2bX\ind_{A^c}]
\]
Вземайки предвид, че $\E\ind_A = \mathbb{P}(A)$, диференцираме получената функция $f(a,b)$ спрямо променливите $a$ и $b$, и приравняваме частните производни на нула, за да намерим минимума:
\[
0 = \frac{\partial f}{\partial a} = 2a\mathbb{P}(A) - 2\E[X\ind_A] \implies a = \frac{\E[X\ind_A]}{\mathbb{P}(A)} = \frac{\E[XY]}{\mathbb{P}(Y=1)}
\]
\[
0 = \frac{\partial f}{\partial b} = 2b\mathbb{P}(A^c) - 2\E[X\ind_{A^c}] \implies b = \frac{\E[X\ind_{A^c}]}{\mathbb{P}(A^c)} = \frac{\E[X(1-Y)]}{\mathbb{P}(Y=0)}
\]
Полученият резултат е минимум (тъй като изразът представлява квадратична форма с положителен хесиан). Функцията, за която се достига минимумът, отбелязваме с $G^*(Y)$:
\[
G^*(Y) = \frac{\E[X\ind_A]}{\mathbb{P}(A)} \ind_A + \frac{\E[X\ind_{A^c}]}{\mathbb{P}(A^c)} \ind_{A^c} =: \E(X\given Y)
\]
Тази функция дефинира \emph{условното математическо очакване} на $X$ при условие $Y$. Вместо да използваме едно единствено число, вече използваме две различни стойности в зависимост от това каква стойност сме наблюдавали за $Y$. Ако клиентът е мъж, използваме оценката $a$, ако е жена — оценката $b$.

За да илюстрираме как условната вероятност естествено възниква в този контекст, нека приемем, че $X$ също е бинарна случайна величина, например $X = \ind_B$. Тогава $\E[X\ind_A] = \E[\ind_B \ind_A] = \mathbb{P}(A \cap B)$. Замествайки във формулата, получаваме:
\[
G^*(Y) = \frac{\mathbb{P}(A \cap B)}{\mathbb{P}(A)} \ind_A + \frac{\mathbb{P}(A^c \cap B)}{\mathbb{P}(A^c)} \ind_{A^c} = \mathbb{P}(B\given A)\ind_A + \mathbb{P}(B\given A^c)\ind_{A^c}
\]
Виждаме, че коефициентите пред индикаторите са точно условните вероятности $\mathbb{P}(B\given A)$ и $\mathbb{P}(B\given A^c)$.

\subsection{Дефиниция за дискретни случайни величини}

Въз основа на изложената геометрична мотивация (като проекция в средноквадратичен смисъл), можем да дефинираме понятието формално за произволни дискретни случайни величини.

\begin{keydefn}
Нека $X$ и $Y$ са дискретни случайни величини. Условното математическо очакване на $X$ при условие, че $Y$ е приело конкретната стойност $y_j$, се дефинира като сума по всички възможни стойности на $X$ с техните условни вероятности:
\[
\E[X \given Y=y_j] = \sum_i x_i \mathbb{P}(B_i \given A_j)
\]
където $B_i = \{X=x_i\}$ са възможните стойности на $X$, а $A_j = \{Y=y_j\}$.
\end{keydefn}

Това е просто число (за дадено фиксирано събитие $Y=y_j$). Ако обаче разглеждаме $Y$ като случайна величина, можем да дефинираме \emph{условното математическо очакване на $X$ спрямо $Y$} (което бележим с $\E[X\given Y]$) като нова дискретна случайна величина, чиято стойност зависи изцяло от това каква стойност ще приеме $Y$:
\[
\E[X \given Y] = \sum_j \E[X \given Y=y_j] \ind_{A_j}
\]

\begin{prop}
Ако $X$ и $Y$ са дискретни случайни величини, то $\E[X\given Y]$ е също дискретна случайна величина, която има следната таблица на разпределение:
\[
\begin{array}{c|c|c|c|c}
\E[X \given Y] & \E[X \given Y=y_1] & \dots & \E[X \given Y=y_j] & \dots \\
\hline
\mathbb{P} & p_1 = \mathbb{P}(A_1) & \dots & p_j = \mathbb{P}(A_j) & \dots
\end{array}
\]
Тази случайна величина изцяло се предопределя от стойностите на $Y$ и приема съответната условна очаквана стойност със същата вероятност, с която $Y$ приема стойността $y_j$.
\end{prop}

\subsection{Свойства на условното математическо очакване}

Условното очакване притежава серия от важни свойства, аналогични на обикновеното математическо очакване. Те са валидни както за дискретни, така и за непрекъснати величини, но тук ще бъдат обосновани интуитивно през призмата на дискретните величини.

\begin{prop}[Свойства на условното математическо очакване]\label{prop:cond-exp-props}
\begin{enumerate}
    \item[а)] \textbf{Линейност:}
    $\E[aX + bZ \given Y] = a\E[X \given Y] + b\E[Z \given Y]$.
    \item[б)] \textbf{Независимост:} ако $X \perp\!\!\!\perp Y$, то $\E[X \given Y] = \E X$.
    \item[в)] \textbf{Детерминираност:} ако $X = f(Y)$, то $\E[X \given Y] = X$.
    \item[г)] \textbf{Повторно очакване (Закон за пълното математическо очакване):}
    $\E[\E[X \given Y]] = \E X$.
    \item[д)] \textbf{Изнасяне на функция от условието:}
    $\E[f(X,Y) \given Y=y_j] = \E[f(X, y_j) \given Y=y_j]$.
\end{enumerate}
\end{prop}

\begin{proof}
а) Използвайки дефиницията за индикаторно представяне:
\begin{align*}
\E[aX + bZ \given Y]
  &= \sum_j \frac{\E[(aX+bZ)\ind_{A_j}]}{\mathbb{P}(A_j)}\ind_{A_j} \\
  &= a \sum_j \frac{\E[X\ind_{A_j}]}{\mathbb{P}(A_j)}\ind_{A_j}
   + b \sum_j \frac{\E[Z\ind_{A_j}]}{\mathbb{P}(A_j)}\ind_{A_j} ,
\end{align*}
което е точно $a\E[X \given Y] + b\E[Z \given Y]$.

б) Тъй като $X$ и $\ind_{A_j}$ са независими, то $\E[X\ind_{A_j}] = \E X \cdot \E\ind_{A_j} = \E X \cdot \mathbb{P}(A_j)$. Тогава в дефиниционната сума вероятностите $\mathbb{P}(A_j)$ се съкращават и остава $\E X \sum_j \ind_{A_j} = \E X \cdot 1 = \E X$ (защото индикаторите на пълна група от събития се сумират до 1).

г) Прилагаме оператора $\E$ върху дефиниционната сума на $\E[X\given Y]$:
\begin{align*}
\E[\E[X\given Y]]
  &= \E\left[\sum_j \frac{\E[X\ind_{A_j}]}{\mathbb{P}(A_j)} \ind_{A_j}\right]
   = \sum_j \frac{\E[X\ind_{A_j}]}{\mathbb{P}(A_j)} \E\ind_{A_j} \\
  &= \sum_j \E[X\ind_{A_j}]
   = \E\left[\sum_j X\ind_{A_j}\right] = \E X
\end{align*}
\end{proof}

Свойства в) и д) са непосредствени. За в): ако познавайки $Y$ ние вече знаем с абсолютна точност стойността на $X$, то случайността отпада и най-добрата оценка е самото $X$. За д): тъй като в условието знаем, че $Y=y_j$, можем да заместим $Y$ с тази константа в самата функция.

Свойство б) има ясен смисъл: когато $Y$ не носи информация за $X$, условното очакване съвпада с безусловното, защото най-доброто число за приближение на $X$ (в средноквадратичен смисъл) остава просто нейното обикновено математическо очакване.

\begin{example}[Смес от разпределения]\label{ex:08-1}

Условното математическо очакване е изключително полезно при т.нар. смеси (mixture distributions), където параметрите на дадено разпределение самите те са случайни величини. Нека разгледаме Бернулиева случайна величина $X \sim \Ber(Y)$, чийто параметър на успеха е случайната величина $Y$:
\[
Y = \begin{cases}
\frac{2}{3}, & \text{с вероятност } p \\
\frac{1}{3}, & \text{с вероятност } 1-p
\end{cases}
\]
Това означава, че първо теглим параметъра $Y$. Ако се падне $2/3$, тогава $X$ е Бернулиева с вероятност за успех $2/3$. Ако $Y$ е $1/3$, тогава $X$ е Бернулиева с вероятност за успех $1/3$. Търсим пълната вероятност $\mathbb{P}(X=1)$, която за индикаторна променлива напълно съвпада с нейното математическо очакване $\E X$. Използвайки свойството за повторно очакване $\E X = \E[\E[X \given Y]]$, имаме:
\begin{align*}
\E[X \given Y] &= \E\left[X \given Y=\frac{2}{3}\right] \ind_{\{Y=2/3\}} + \E\left[X \given Y=\frac{1}{3}\right] \ind_{\{Y=1/3\}} \\
&= \frac{2}{3} \ind_{\{Y=2/3\}} + \frac{1}{3} \ind_{\{Y=1/3\}}
\end{align*}
Вземайки безусловното математическо очакване от двете страни, получаваме крайната вероятност за успех:
\[
\mathbb{P}(X=1) = \E[\E[X \given Y]] = \frac{2}{3} \mathbb{P}\left(Y=\frac{2}{3}\right) + \frac{1}{3} \mathbb{P}\left(Y=\frac{1}{3}\right) = \frac{2}{3}p + \frac{1}{3}(1-p)
\]

\end{example}

\section{Условни разпределения}

Аналогично на условното очакване, можем да дефинираме цялостното \emph{условно разпределение} на една случайна величина спрямо друга.

\begin{defn}
Нека $X$ и $Y$ са дискретни случайни величини. Условното разпределение на $X$ при условие $Y=y_j$ (за всяка възможна стойност на $Y$ с $\mathbb{P}(Y=y_j)>0$) се задава чрез условните вероятности за всяка възможна стойност $x_i$:
\[
\begin{array}{c|c|c|c|c|c}
X \given Y=y_j & x_1 & x_2 & \dots & x_i & \dots \\
\hline
\mathbb{P} & \mathbb{P}(X=x_1 \given Y=y_j) & \dots & \dots & \mathbb{P}(X=x_i \given Y=y_j) & \dots
\end{array}
\]
Това на практика са изменените вероятности за $X$, които произтичат от допълнителната информация, че е настъпило събитието $Y=y_j$.
\end{defn}

\begin{example}\label{ex:08-2}
Да разгледаме хвърляне на два зара. Нека $Y$ е броят на падналите се шестици, а $X$ е броят на падналите се единици. И двете случайни величини приемат стойности в множеството $\{0, 1, 2\}$. Съвместното им разпределение е зададено със следната таблица (известна от предходна лекция):
\[
\begin{array}{c|c|c|c}
Y \setminus X & 0 & 1 & 2 \\
\hline
0 & 16/36 & 8/36 & 1/36 \\
1 & 8/36 & 2/36 & 0 \\
2 & 1/36 & 0 & 0
\end{array}
\]
Искаме да намерим условното разпределение на $X$ при условие, че $Y=0$. Маргиналната вероятност за $Y=0$ е $\mathbb{P}(Y=0) = \frac{16}{36} + \frac{8}{36} + \frac{1}{36} = \frac{25}{36}$. За да намерим условните вероятности $\mathbb{P}(X=x_i \given Y=0)$, просто разделяме съвместната вероятност от таблицата на маргиналната вероятност:
\begin{align*}
\mathbb{P}(X=0 \given Y=0) &= \frac{16/36}{25/36} = \frac{16}{25} \\
\mathbb{P}(X=1 \given Y=0) &= \frac{8/36}{25/36} = \frac{8}{25} \\
\mathbb{P}(X=2 \given Y=0) &= \frac{1/36}{25/36} = \frac{1}{25}
\end{align*}
Този резултат има напълно ясна интуиция. Ако знаем със сигурност, че $Y=0$ (не се е паднала нито една шестица), това е еквивалентно на това да хвърляме зарове само с по 5 стени (защото сме елиминирали напълно шестицата). Възможните изходи от хвърлянето на два зара намаляват до $5^2 = 25$. Вероятността да не се падне единица е $\left(\frac{4}{5}\right)^2 = \frac{16}{25}$, да се падне точно една единица е $2 \cdot \frac{1}{5} \cdot \frac{4}{5} = \frac{8}{25}$, а две единици - $\left(\frac{1}{5}\right)^2 = \frac{1}{25}$.

\end{example}

\section{Полиномно разпределение}

Полиномното разпределение е естествено многомерно обобщение на биномното разпределение. Докато при биномното разпределение имаме само два възможни изхода (успех и неуспех), при полиномното разполагаме с $r \ge 1$ възможни изхода.
Нека провеждаме $n$ независими експеримента. Всеки експеримент може да завърши с един от $r$ изхода. Съответните вероятности за тези изходи са $p_1, p_2, \dots, p_r$, като сумата им естествено е $\sum_{j=1}^{r} p_j = 1$.

Интересува ни съвместното разпределение на вектора от случайни величини $(X_1, X_2, \dots, X_r)$, където $X_j$ е броят на настъпванията на изход $j$. Зададени са цели неотрицателни числа $k_1, k_2, \dots, k_r$, такива че сборът им дава общия брой експерименти: $k_1 + k_2 + \dots + k_r = n$.
Вероятността да се получи точно тази конфигурация от резултати се пресмята по следната формула:
\[
\mathbb{P}(X_1=k_1, X_2=k_2, \dots, X_r=k_r) = \frac{n!}{k_1! \dots k_r!} p_1^{k_1} p_2^{k_2} \dots p_r^{k_r}
\]
Тази формула се извежда аналогично на биномното разпределение. Вероятността за една конкретна наредба на тези изходи (т.е. конкретна последователност от $n$ събития) е произведението $p_1^{k_1} \dots p_r^{k_r}$. Тъй като нас не ни интересува конкретният ред на настъпване на събитията, а само общият им брой, трябва да умножим по броя на всички възможни конфигурации, които водят до този краен резултат.
Броят на начините да изберем $k_1$ места за първия изход е $\binom{n}{k_1}$. За останалите $n-k_1$ свободни места избираме $k_2$ места за втория изход, което може да стане по $\binom{n-k_1}{k_2}$ начина, и така нататък. Произведението от тези биномни коефициенти се опростява алгебрично до мултиномния коефициент:
\[
\binom{n}{k_1} \binom{n-k_1}{k_2} \dots \binom{n-k_1-\dots-k_{r-1}}{k_r} = \frac{n!}{k_1! \dots k_r!}
\]
Ако $r=2$, получаваме класическото биномно разпределение, където $k_2 = n - k_1$ и $p_2 = 1 - p_1$.

\section{Непрекъснати случайни величини}

\subsection{Въведение и мотивация}
Досега разглеждахме дискретни случайни величини, които могат да приемат най-много изброимо множество от стойности. Има обаче много случайни характеристики от реалния живот (например температура, тегло, време на живот на процесор, приходи), които чисто теоретично могат да приемат неизброимо много стойности — всяка стойност от някакъв реален интервал.

От гледна точка на компютърното моделиране, неизброимостта е фикция — ние винаги съхраняваме данните с крайна точност (чрез тактовете на процесора). В такъв случай защо въобще са ни нужни непрекъснати случайни величини?
Отговорът се крие в ефективността. Ако искаме да моделираме равномерна случайна величина в интервала $(0,1)$ с много голяма точност (например 100 милиона възможни стойности), трябва да пазим масив със 100 милиона вероятности, въпреки че всички те са напълно еднакви. Много по-изчислително ефективно е да използваме една проста функция (плътност), която компактно да описва поведението на модела.
Освен това, на практика изключително рядко се интересуваме от точната вероятност температурата да бъде съвършено равняваща се на дадено число с безкрайна точност (която вероятност клони към нула). Много по-често ни интересува вероятността стойността да попадне в даден диапазон.

\subsection{Дефиниция и плътност на разпределение}

\begin{defn}[Абсолютно непрекъсната случайна величина]
Случайната величина $X$ се нарича непрекъсната, ако съществува неотрицателна функция $f_X: \mathbb{R} \to [0, \infty)$, наричана \emph{плътност на разпределение} (probability density function), такава че:
\begin{enumerate}
    \item $f_X(x) \ge 0 \quad \forall x \in \mathbb{R}$
    \item Лицето под цялата крива е единица: $\int_{-\infty}^{\infty} f_X(x) dx = 1$
    \item За всеки две числа $a < b$ (включително $\pm\infty$) е изпълнено:
    \[
    \mathbb{P}(X \in (a, b)) = \int_{a}^{b} f_X(x) dx
    \]
\end{enumerate}
\end{defn}

Графиката на функцията $f_X(x)$ ни показва къде случайните величини са по-склонни да приемат стойности. Лицето под графиката между точките $a$ и $b$ е точно равно на вероятността $X$ да попадне в интервала $(a,b)$.

От свойствата на интеграла веднага следва, че вероятността $X$ да приеме точно една конкретна стойност $c$ е нула:
\[
\mathbb{P}(X = c) = \int_{c}^{c} f_X(x) dx = 0
\]
Поради това за непрекъснати случайни величини границите на интервала нямат никакво значение при пресмятане на вероятност:
\[
\mathbb{P}(a < X < b) = \mathbb{P}(a \le X < b) = \mathbb{P}(a < X \le b) = \mathbb{P}(a \le X \le b)
\]

\begin{example}[Медицински застраховки]\label{ex:08-3}
Нека разходите за медицински застраховки в една фирма са зададени с модела $M = 100\,000 \cdot X$, където $X$ е непрекъсната случайна величина с плътност:
\[
f_X(x) = \begin{cases}
5(1-x)^4, & x \in (0,1) \\
0, & \text{иначе}
\end{cases}
\]
Тази плътност е нула за отрицателни стойности и за стойности по-големи от 1. В интервала $(0,1)$ представлява намаляваща крива (започваща от стойност 5 при $x=0$). Търсим вероятността разходите на фирмата да надвишат 10 000 единици:
\[
\mathbb{P}(M > 10\,000) = \mathbb{P}\left(100\,000 \cdot X > 10\,000\right) = \mathbb{P}\left(X > \frac{1}{10}\right)
\]
Тъй като $X$ приема стойности само до 1, това се свежда до определен интеграл:
\[
\mathbb{P}\left(\frac{1}{10} < X < 1\right) = \int_{1/10}^{1} 5(1-x)^4 dx
\]
Този интеграл се решава изключително лесно и е равен на $\left(\frac{9}{10}\right)^5 \approx 0{,}59$.
Вместо да пази огромен масив с дискретни вероятности за всяка възможна стотинка, фирмата използва компактна функция за оценка на риска.

\end{example}

\subsection{Функция на разпределение}

Функцията на разпределение $F_X(x)$ за непрекъсната случайна величина $X$ се дефинира по обичайния начин като вероятността $X$ да бъде по-малко от $x$:
\[
F_X(x) = \mathbb{P}(X < x) = \int_{-\infty}^{x} f_X(y) dy \quad \forall x \in \mathbb{R}
\]
Това представлява натрупване (акумулиране) на площта под плътността от минус безкрайност до текущата точка $x$. Диференциалът $dy$ в този контекст трябва да се разбира като мярка.
Ако плътността $f_X$ е непрекъсната функция в дадена точка $x_0$, то от основната теорема на интегралното смятане следва, че производната на функцията на разпределение е равна на плътността:
\[
\left. \frac{dF_X}{dx} \right|_{x=x_0} = f_X(x_0) \implies F_X' = f_X
\]
Следователно, ако знаем плътността, чрез интегриране възстановяваме функцията на разпределение, а ако знаем функцията на разпределение, чрез диференциране възстановяваме плътността.

\section{Смяна на променливите при непрекъснати случайни величини}

Често имаме даден модел със случайна величина $X$, но се интересуваме от някаква нейна трансформация $Y = g(X)$ (например ако $X$ е приходът на клиент, $Y$ може да бъде неговият разход за конкретни мероприятия, който зависи функционално от прихода). Искаме да намерим плътността на новата случайна величина $Y$.

\begin{thm}
Нека $X$ е непрекъсната случайна величина, а $g: \mathbb{R} \to \mathbb{R}$ е строго монотонно растяща или строго монотонно намаляваща функция (поне в областта, където $X$ приема стойности с положителна плътност). Тогава $Y = g(X)$ е също непрекъсната случайна величина, чиято плътност се задава от формулата:
\[
f_Y(y) = f_X(g^{-1}(y)) \left| (g^{-1}(y))' \right|
\]
\end{thm}

\begin{proof}
Търсим функция $f_Y$, която при интегриране над произволен интервал $(a, b)$ да дава вероятността $\mathbb{P}(Y \in (a, b))$.
Разглеждаме първо случая, когато $g$ е строго монотонно растяща ($g^{\uparrow}$).
\[
\mathbb{P}(Y \in (a, b)) = \mathbb{P}(g(X) \in (a, b)) = \mathbb{P}(X \in (g^{-1}(a), g^{-1}(b)))
\]
Тъй като познаваме плътността на $X$, тази вероятност е просто интеграл:
\[
\mathbb{P}(Y \in (a, b)) = \int_{g^{-1}(a)}^{g^{-1}(b)} f_X(w) dw
\]
За да превърнем границите на интегриране отново в $a$ и $b$, правим субституцията $w = g^{-1}(v)$. Тогава диференциалът е $dw = (g^{-1}(v))' dv$. Границите на интегриране се променят от $g^{-1}(a)$ до $g^{-1}(b)$ съответно на $a$ до $b$:
\[
\int_{a}^{b} f_X(g^{-1}(v)) (g^{-1}(v))' dv
\]
Тъй като $g$ е монотонно растяща функция, нейната обратна функция $g^{-1}$ също е монотонно растяща, което означава, че производната $(g^{-1}(v))'$ е положителна, т.е. съвпада със своя модул.

Нека сега разгледаме случая, когато $g$ е строго монотонно намаляваща ($g^{\downarrow}$). Когато приложим $g^{-1}$ към неравенството $a < g(X) < b$, посоките на неравенствата се обръщат, защото функцията обръща наредбата:
\[
\mathbb{P}(g(X) \in (a, b)) = \mathbb{P}(X \in (g^{-1}(b), g^{-1}(a))) = \int_{g^{-1}(b)}^{g^{-1}(a)} f_X(w) dw
\]
При същата субституция $w = g^{-1}(v)$ границите на интегриране стават от $b$ до $a$. За да обърнем нормално интеграла да бъде от $a$ до $b$, трябва да добавим знак минус:
\[
-\int_{a}^{b} f_X(g^{-1}(v)) (g^{-1}(v))' dv
\]
Но тъй като $g$ е монотонно намаляваща, $(g^{-1}(v))'$ е отрицателна функция, следователно $-(g^{-1}(v))' = \left| (g^{-1}(v))' \right|$. И в двата случая стигаме до един и същ обединен израз:
\[
\int_{a}^{b} f_X(g^{-1}(v)) \left| (g^{-1}(v))' \right| dv
\]
Тъй като вероятността $\mathbb{P}(Y \in (a, b))$ се пресмята чрез интеграл над интервала $(a,b)$ от тази конкретна функция, то по дефиниция подинтегралната функция е именно търсената плътност на $Y$.
\end{proof}

\section{Числени характеристики на непрекъснатите величини}

\subsection{Математическо очакване и дисперсия}

По аналогичен начин на дискретните случайни величини, където математическото очакване се получаваше като сума от произведенията на стойностите с техните вероятности $\E X = \sum x_i p_i$, при непрекъснатите величини сумата логично се заменя с интеграл (тъй като интегралът е граница на сума), а вероятностите преминават в плътност.

\begin{defn}
Нека $X$ е непрекъсната случайна величина. Ако интегралът $\int_{-\infty}^{\infty} |x| f_X(x) dx < \infty$ е краен, то \emph{математическото очакване} на $X$ се дефинира като:
\[
\E X = \int_{-\infty}^{\infty} x f_X(x) dx
\]
От гледна точка на механиката, ако разглеждаме плътността $f_X(x)$ като разпределение на масата върху пръчка, очакването $\E X$ е точно центърът на тежестта на тази пръчка.
\end{defn}

\begin{defn}[дисперсия на непрекъсната случайна величина]\label{def:var-continuous}
\emph{Дисперсията} $\Var X$ продължава да се дефинира като средноквадратичното отклонение от очакването:
\[
\Var X = \E(X - \E X)^2 = \int_{-\infty}^{\infty} (x - \E X)^2 f_X(x) \, dx .
\]
\end{defn}

Ако искаме да пресметнем математическото очакване на функция от случайната величина $g(X)$, не е необходимо първо да намираме плътността на $Y=g(X)$. Използваме директната формула (която се доказва строго чрез теория на мярката, но следва интуитивно от току-що доказаната теорема за смяна на променливите):
\[
\E[g(X)] = \int_{-\infty}^{\infty} g(x) f_X(x) dx
\]

\textbf{Свойства на очакването и дисперсията:}
Следните свойства остават изцяло в сила както за дискретни, така и за непрекъснати случайни величини (нека $a, b, c$ са константи):
\begin{itemize}
    \item Очакване на константа: $\E c = c$
    \item Дисперсия на константа: $\Var c = 0$
    \item Линейност на очакването: $\E(aX + b) = a\E X + b$ и $\E(aX + bY) = a\E X + b\E Y$
    \item Изместване на дисперсията: $\Var(X + c) = \Var X$
    \item Мащабиране на дисперсията: $\Var(aX) = a^2 \Var X$
    \item Независимост: Ако $X \perp\!\!\!\perp Y$, то $\E[XY] = \E X \cdot \E Y$ и $\Var(X + Y) = \Var X + \Var Y$
\end{itemize}

\subsection{Равномерно разпределение}

Равномерното разпределение (Uniform distribution) е най-простият клас непрекъснати случайни величини. То се използва изключително често като априорно разпределение, когато нямаме никаква допълнителна информация или експертни познания за системата и искаме да избегнем субективизъм (т.е. стартираме на чисто).

Казваме, че $X$ е равномерно разпределена в интервала $[a, b]$, където $a < b$ са крайни числа, ако плътността ѝ е постоянна в този интервал и нула извън него. За да бъде лицето под графиката на плътността равно на 1 (както изисква дефиницията), плътността трябва да е:
\[
f_X(x) = \begin{cases}
\frac{1}{b-a}, & x \in [a, b] \\
0, & \text{иначе}
\end{cases}
\]
Лесно се проверява, че $\int_{a}^{b} \frac{1}{b-a} dx = \frac{1}{b-a} \cdot (b-a) = 1$.
Нарича се равномерно, защото за подинтервали с еднаква дължина вероятността случайната величина да попадне в тях е напълно еднаква, без значение къде се намира подинтервалът в рамките на $[a, b]$.

Функцията на разпределение за равномерното разпределение се получава чрез директно интегриране:
\[
F_X(x) = \begin{cases}
0, & x \le a \\
\frac{x-a}{b-a}, & x \in (a, b) \\
1, & x \ge b
\end{cases}
\]
Ако $x \in (a, b)$, то $F_X(x) = \int_{-\infty}^{x} f_X(y) dy = \int_{a}^{x} \frac{1}{b-a} dy = \frac{x-a}{b-a}$.
За $x \le a$ интегралът е от 0, следователно $0$. За $x \ge b$ плътността след $b$ отново е 0, така че целият интеграл се свежда до площта над целия интервал $[a, b]$, която е 1.

\section{Задачи}

\begin{enumerate}
    \item Нека $Y=aX+b$, където $a\ne0$. Използвайте теоремата за смяна на променливите, за да намерите плътността на $Y$ чрез плътността на $X$.
    \item Проверете от дефиницията, че функцията
    \[
    f(x)=\frac{1}{b-a}\ind_{[a,b]}(x),\qquad a<b,
    \]
    е плътност: тя е неотрицателна и интегралът ѝ върху реалната права е единица.
\end{enumerate}


\chapter{Непрекъснати разпределения. Гама и хи-квадрат разпределение}
\section{Класове непрекъснати случайни величини}

В тази част от лекцията ще разгледаме три основни класа непрекъснати случайни величини (равномерно, нормално и експоненциално разпределение), които имат важна роля в теорията на вероятностите, и ще покажем как структурно могат да бъдат свързани със своите стандартизирани версии чрез подходящи трансформации. Тази техника значително улеснява пресмятанията на техните моменти.

\subsection{Равномерно разпределение}

Казваме, че случайната величина $X$ има равномерно разпределение в интервала $[a, b]$, и записваме $X \sim U(a, b)$, където $-\infty < a < b < \infty$, ако нейната плътност се задава с:
\[
f_X(x) = 
\begin{cases} 
\frac{1}{b-a}, & \text{за } x \in [a, b] \\ 
0, & \text{иначе} 
\end{cases}
\]

Една от важните роли на равномерното разпределение е, че когато знаем, че възможните стойности на даден експеримент лежат в интервала $[a, b]$, но нямаме допълнителна информация, то избирайки това разпределение ние внасяме най-малко субективизъм (принцип на максималната ентропия). Всеки подинтервал с една и съща дължина има една и съща вероятност.

Функцията на разпределение $F_X(x)$ е:
\[
F_X(x) = 
\begin{cases} 
0, & \text{за } x \le a \\ 
\frac{x-a}{b-a}, & \text{за } a < x < b \\ 
1, & \text{за } x \ge b 
\end{cases}
\]

Очакваната стойност се пресмята стандартно като център на тежестта:
\[
\mathbb{E}[X] = \int_{a}^{b} x \frac{1}{b-a} dx = \frac{a+b}{2}
\]

\textbf{Структурна връзка със стандартно равномерно разпределение:} 
Вместо да пресмятаме сложни интеграли за всяко конкретно $a$ и $b$, можем да сведем $X$ до стандартното равномерно разпределение $Y \sim U(0, 1)$. Ако разгледаме трансформацията $y = g(x) = \frac{x-a}{b-a}$, която е строго монотонно растяща, то обратната ѝ функция е $x = g^{-1}(y) = a + (b-a)y$. 

Ако дефинираме $Z = g(X) = \frac{X-a}{b-a}$, по теоремата за смяна на променливите плътността на $Z$ се дава от:
\[
f_Z(y) = f_X(g^{-1}(y)) \left| \frac{d}{dy} g^{-1}(y) \right| = f_X(a + (b-a)y) \cdot (b-a)
\]
Понеже $y \in [0, 1] \implies a + (b-a)y \in [a, b]$, в този интервал $f_X$ е $\frac{1}{b-a}$. Следователно:
\[
f_Z(y) = \frac{1}{b-a} \cdot (b-a) = 1 \quad \text{за } y \in [0, 1]
\]
Това доказва, че $Z \sim U(0, 1)$. 

Чрез обратната връзка $X = a + (b-a)Z$ можем лесно да изчислим моментите на $X$. Очакването е:
\[
\mathbb{E}[X] = \mathbb{E}[a + (b-a)Z] = a + (b-a)\mathbb{E}[Z] = a + (b-a)\frac{1}{2} = \frac{a+b}{2}
\]
Дисперсията е:
\[
\text{Var}(X) = \text{Var}(a + (b-a)Z) = (b-a)^2 \text{Var}(Z)
\]
За $Z \sim U(0,1)$, вторият момент е $\mathbb{E}[Z^2] = \int_0^1 z^2 \cdot 1 dz = \frac{1}{3}$. Тогава $\text{Var}(Z) = \frac{1}{3} - \left(\frac{1}{2}\right)^2 = \frac{1}{12}$. Оттук получаваме:
\[
\text{Var}(X) = \frac{(b-a)^2}{12}
\]

\subsection{Нормално разпределение (Гаусово)}

Случайната величина $X$ има нормално разпределение с параметри $\mu$ (средно) и $\sigma^2$ (дисперсия, $\sigma > 0$), което се записва като $X \sim N(\mu, \sigma^2)$, ако нейната плътност е:
\[
f_X(x) = \frac{1}{\sqrt{2\pi}\sigma} e^{-\frac{(x-\mu)^2}{2\sigma^2}}, \quad \forall x \in \mathbb{R}
\]
Тази крива, известна като камбанка (bell-shaped curve), има огромно значение в природата и се появява естествено благодарение на Централната гранична теорема (ЦГТ). Интегралът от тази плътност над цялата реална ос е равен на 1.

\begin{defn}[стандартно нормално разпределение]\label{def:std-normal}
Специалният случай, когато $\mu = 0$ и $\sigma^2 = 1$, се нарича \emph{стандартно нормално разпределение} и се бележи със $Z \sim N(0, 1)$. Неговата плътност е
\[
f_Z(x) = \frac{1}{\sqrt{2\pi}} e^{-\frac{x^2}{2}} .
\]
\end{defn}
Функцията на разпределение на стандартното нормално разпределение няма затворена форма чрез елементарни функции и традиционно се означава с главната буква $\Phi(x)$:
\[
\Phi(x) = \int_{-\infty}^{x} \frac{1}{\sqrt{2\pi}} e^{-\frac{y^2}{2}} dy
\]

\textbf{Стандартизиране на нормална случайна величина:} 
Нека $X \sim N(\mu, \sigma^2)$. Въвеждаме трансформацията $Y = g(X) = \frac{X-\mu}{\sigma}$. Обратната функция е $x = g^{-1}(y) = \sigma y + \mu$ с производна $\sigma$.
Плътността на $Y$ е:
\[
f_Y(y) = f_X(\sigma y + \mu) \cdot \sigma = \frac{1}{\sqrt{2\pi}\sigma} e^{-\frac{(\sigma y + \mu - \mu)^2}{2\sigma^2}} \cdot \sigma = \frac{1}{\sqrt{2\pi}} e^{-\frac{y^2}{2}}
\]
Това съвпада с плътността на $Z$, следователно $\frac{X-\mu}{\sigma} \sim N(0, 1)$. Обратната връзка е $X = \sigma Z + \mu$.

Тази структурна връзка позволява лесно изчисляване на моментите. Очакването е:
\[
\mathbb{E}[X] = \sigma\mathbb{E}[Z] + \mu
\]
Тъй като подинтегралната функция $y e^{-y^2/2}$ за очакването на $Z$ е нечетна и се интегрира в симетрични граници $(-\infty, \infty)$, следва че $\mathbb{E}[Z] = 0$. Следователно $\mathbb{E}[X] = \mu$.

За дисперсията имаме $\text{Var}(X) = \sigma^2 \text{Var}(Z) = \sigma^2 \mathbb{E}[Z^2]$. За да пресметнем втория момент $\mathbb{E}[Z^2]$, ще използваме елегантен метод на Ричард Файнман — диференциране под знака на интеграла по параметър. 
Знаем, че интегралът от общата нормална плътност е 1, т.е.:
\[
\sqrt{2\pi}\sigma = \int_{-\infty}^{\infty} e^{-\frac{y^2}{2\sigma^2}} dy
\]
Диференцираме двете страни по параметъра $\sigma$:
\[
\sqrt{2\pi} = \int_{-\infty}^{\infty} \frac{\partial}{\partial \sigma} \left( e^{-\frac{y^2}{2\sigma^2}} \right) dy = \int_{-\infty}^{\infty} e^{-\frac{y^2}{2\sigma^2}} \left(-\frac{y^2}{2}\right) \left(-\frac{2}{\sigma^3}\right) dy
\]
Оттук получаваме равенството:
\[
\sqrt{2\pi}\sigma^3 = \int_{-\infty}^{\infty} y^2 e^{-\frac{y^2}{2\sigma^2}} dy
\]
Като заместим $\sigma = 1$, намираме втория момент на стандартното нормално разпределение:
\[
\sqrt{2\pi} = \int_{-\infty}^{\infty} y^2 e^{-\frac{y^2}{2}} dy \implies \mathbb{E}[Z^2] = \frac{1}{\sqrt{2\pi}} \int_{-\infty}^{\infty} y^2 e^{-\frac{y^2}{2}} dy = \frac{\sqrt{2\pi}}{\sqrt{2\pi}} = 1
\]
Следователно $\text{Var}(Z) = 1$, а дисперсията на първоначалната величина е $\text{Var}(X) = \sigma^2$.

За пресмятане на вероятности с общо нормално разпределение, свеждаме до $\Phi(x)$:
\[
\mathbb{P}(X < x) = \mathbb{P}\left(\frac{X-\mu}{\sigma} < \frac{x-\mu}{\sigma}\right) = \Phi\left(\frac{x-\mu}{\sigma}\right)
\]
Например, $\Phi(0) = \frac{1}{2}$ поради симетрията. Също така класическа стойност е $\Phi(1{,}96) \approx 0{,}975$, което означава, че $95\%$ от вероятностната маса е в интервала $[-1{,}96, 1{,}96]$.

\subsection{Експоненциално разпределение}

Казваме, че $X$ е експоненциално разпределена с параметър $\lambda > 0$, записвано $X \sim \text{Exp}(\lambda)$, ако нейната плътност е:
\[
f_X(x) = 
\begin{cases} 
\lambda e^{-\lambda x}, & \text{за } x \ge 0 \\ 
0, & \text{иначе} 
\end{cases}
\]

Функцията на разпределение е:
\[
F_X(x) = \int_{0}^{x} \lambda e^{-\lambda y} dy = 1 - e^{-\lambda x}
\]
Често е по-удобно да се работи с опашката на разпределението (tail probability), която се бележи с $\bar{F}_X(x)$:
\[
\bar{F}_X(x) = \mathbb{P}(X \ge x) = 1 - F_X(x) = e^{-\lambda x}
\]

Този клас разпределения е фундаментален заради \textbf{свойството си на безпаметност (memorylessness)}. То отразява феномена, при който ако един процес/компонент е работил $x$ време, вероятността да продължи да работи още $y$ време е същата, сякаш е започнал да работи отначало (чисто нов). Формално:
\begin{align*}
\mathbb{P}(X \ge x+y \mid X \ge x)
  &= \frac{\mathbb{P}(X \ge x+y \text{ и } X \ge x)}{\mathbb{P}(X \ge x)}
   = \frac{\mathbb{P}(X \ge x+y)}{\mathbb{P}(X \ge x)} \\
  &= \frac{e^{-\lambda(x+y)}}{e^{-\lambda x}} = e^{-\lambda y} = \mathbb{P}(X \ge y)
\end{align*}

Обратно, ако търсим функция $\bar{F}$, която да удовлетворява функционалното уравнение за безпаметност $\bar{F}(x+y) = \bar{F}(x)\bar{F}(y)$, и изискваме отговарящото ѝ разпределение да притежава плътност, единственото възможно решение е експоненциалната функция $\bar{F}(x) = e^{-\lambda x}$.

Структурната връзка тук се изразява чрез факта, че ако $X \sim \text{Exp}(\lambda)$, то мащабираната величина $Y = \lambda X \sim \text{Exp}(1)$. Проверката е директна:
\[
\mathbb{P}(Y \le x) = \mathbb{P}(\lambda X \le x) = \mathbb{P}\left(X \le \frac{x}{\lambda}\right) = 1 - e^{-\lambda \left(\frac{x}{\lambda}\right)} = 1 - e^{-x} = F_Y(x)
\]
Тъй като $Y$ е с параметър $1$, нейните моменти са $\mathbb{E}[Y] = 1$ и $\text{Var}(Y) = 1$. Чрез линейните свойства на очакването намираме за $X$:
\[
\mathbb{E}[X] = \mathbb{E}\left[\frac{Y}{\lambda}\right] = \frac{1}{\lambda}\mathbb{E}[Y] = \frac{1}{\lambda}
\]
\[
\text{Var}(X) = \text{Var}\left(\frac{Y}{\lambda}\right) = \frac{1}{\lambda^2}\text{Var}(Y) = \frac{1}{\lambda^2}
\]

\section{Непрекъснати съвместни разпределения}

В практиката най-често се работи със съвкупност от случайни величини, които може да имат сложни зависимости помежду си. Ето защо въвеждаме случайни вектори.

\begin{defn}
Векторът $X = (X_1, \dots, X_n)$ е непрекъснат вектор от случайни величини, ако съществува функция $f_X: \mathbb{R}^n \to [0, \infty)$, наречена \textbf{съвместна плътност} (joint probability density function), такава че:
1. $f_X(x) \ge 0 \quad \forall x \in \mathbb{R}^n$
2. $\int_{-\infty}^{\infty} \dots \int_{-\infty}^{\infty} f_X(x_1, \dots, x_n) dx_1 \dots dx_n = 1$
3. За всяко подмножество $D \subseteq \mathbb{R}^n$: 
   \[
   \mathbb{P}(X \in D) = \int \dots \int_D f_X(x) dx
   \]
\end{defn}

Множеството, върху което съвместната плътност е строго положителна, се нарича дефиниционно множество:
\[
\mathcal{D}(f_X) = \{ x \in \mathbb{R}^n : f_X(x) > 0 \}
\]

Ако знаем съвместната плътност на целия вектор, можем да възстановим плътността на всеки индивидуален компонент $X_j$.

\begin{defn}[маргинална плътност]\label{def:marginal}
\emph{Маргинална плътност} на компонента $X_j$ наричаме
\[
f_{X_j}(x_j) = \int_{-\infty}^{\infty} \dots \int_{-\infty}^{\infty} f_X(x_1, \dots, x_n) \, dx_1 \dots dx_{j-1} dx_{j+1} \dots dx_n ,
\]
тоест резултатът от „дезинтегриране“ по всички останали променливи.
\end{defn}

\textbf{Съвместна функция на разпределение:}
Тя се дефинира общо (не само за непрекъснати величини) като:
\[
F_X(x) = \mathbb{P}(X_1 \le x_1, \dots, X_n \le x_n) = \int_{-\infty}^{x_1} \dots \int_{-\infty}^{x_n} f_X(y_1, \dots, y_n) dy_1 \dots dy_n
\]
От функцията на разпределение може да се възстанови плътността чрез $n$-кратна смесена производна:
\[
f_X(x_1, \dots, x_n) = \frac{\partial^n}{\partial x_1 \dots \partial x_n} F_X(x_1, \dots, x_n)
\]

\textbf{Независимост:}
Две случайни величини $X_1$ и $X_2$ са независими ($X_1 \perp \!\!\! \perp X_2$), тогава и само тогава, когато съвместната им функция на разпределение се разпада на произведение от индивидуалните:
\[
F_X(x_1, x_2) = F_{X_1}(x_1) F_{X_2}(x_2)
\]
В случая на непрекъснат вектор, това е еквивалентно на разпадане на съвместната плътност:
\[
f_X(x_1, x_2) = f_{X_1}(x_1) f_{X_2}(x_2)
\]

Случайните величини $X_1, \dots, X_n$ са \textbf{независими в съвкупност}, ако за всяко $k$ ($2 \le k \le n$) и за всеки избор на $k$ различни индекса $j_1, \dots, j_k$, съвместната плътност на съответния подвектор е произведение от маргиналните плътности:
\[
f_Y(y_{j_1}, \dots, y_{j_k}) = \prod_{i=1}^{k} f_{X_{j_i}}(y_{j_i})
\]
В частност, плътността на целия вектор е $f_X(x_1, \dots, x_n) = \prod_{i=1}^n f_{X_i}(x_i)$.

\textbf{Свойства на очакването:}
Ако $X$ е непрекъснат вектор и $g$ е функция, то математическото очакване на $g(X)$ се задава като:
\[
\mathbb{E}[g(X)] = \int_{-\infty}^{\infty} \dots \int_{-\infty}^{\infty} g(x) f_X(x) dx
\]
Очакването е винаги линеен оператор. Ако $X = (X_1, X_2)$, то за функцията $g(X_1, X_2) = X_1 + X_2$:
\[
\mathbb{E}[X_1 + X_2] = \int_{-\infty}^{\infty} \int_{-\infty}^{\infty} (x_1 + x_2) f_X(x_1, x_2) dx_1 dx_2
\]
Можем да разделим интеграла на две части. Като интегрираме първата част по $x_2$, получаваме маргиналната плътност на $X_1$:
\[
\int_{-\infty}^{\infty} x_1 \left( \int_{-\infty}^{\infty} f_X(x_1, x_2) dx_2 \right) dx_1 = \int_{-\infty}^{\infty} x_1 f_{X_1}(x_1) dx_1 = \mathbb{E}[X_1]
\]
Аналогично, втората част е $\mathbb{E}[X_2]$. Следователно $\mathbb{E}[X_1 + X_2] = \mathbb{E}[X_1] + \mathbb{E}[X_2]$. 

Ако допълнително $X_1$ и $X_2$ са независими, то тогава:
\[
\mathbb{E}[X_1 X_2] = \mathbb{E}[X_1] \mathbb{E}[X_2]
\]
\[
\text{Var}(X_1 + X_2) = \text{Var}(X_1) + \text{Var}(X_2)
\]

\section{Смяна на променливите при случайни вектори}

Често се налага да намерим разпределението на нов случаен вектор $Y$, който е функция от познат вектор $X$.

\begin{thm}[Смяна на променливите]
Нека $X = (X_1, X_2)$ е непрекъснат вектор със съвместна плътност $f_X$ и дефиниционна област $\mathcal{D}(f_X)$. 
Нека $g: \mathcal{D}(f_X) \to \mathbb{R}^2$ е взаимно еднозначно изображение, т.е. $Y = g(X) = (g_1(X), g_2(X))$.
Обратната функция бележим с $X = g^{-1}(Y) = h(Y) = (h_1(Y), h_2(Y))$. 
Нека $g$ и $h$ са непрекъснати и диференцируеми в съответните си области $\mathcal{D}(f_X)$ и $g(\mathcal{D}(f_X))$.
Тогава съвместната плътност на $Y$ съществува и е равна на:
\[
f_Y(y) = f_X(h(y)) |J(y)|, \quad \forall y \in g(\mathcal{D}(f_X))
\]
където $J(y)$ е якобианата (детерминантата на матрицата от производни):
\[
J(y) = \det \begin{pmatrix} 
\frac{\partial h_1}{\partial y_1} & \frac{\partial h_1}{\partial y_2} \\ 
\frac{\partial h_2}{\partial y_1} & \frac{\partial h_2}{\partial y_2} 
\end{pmatrix}
\]
\end{thm}

Доказателството стъпва директно върху теоремата за смяна на променливите при двукратни интеграли от математическия анализ, като се запише вероятността $\mathbb{P}(Y \in A) = \mathbb{P}(X \in h(A))$.

\textbf{Обща схема за сума на независими случайни величини:}
Нека $X_1$ и $X_2$ са независими, откъдето $f_X(x_1, x_2) = f_{X_1}(x_1) f_{X_2}(x_2)$. Интересуваме се от плътността на сумата $Y_2 = X_1 + X_2$.
За да приложим теоремата, допълваме с тривиална променлива $Y_1 = X_1$:
\[
Y = \begin{pmatrix} Y_1 \\ Y_2 \end{pmatrix} = \begin{pmatrix} X_1 \\ X_1 + X_2 \end{pmatrix} = g(X)
\]
Обратната трансформация е лесна: $X_1 = Y_1$ и $X_2 = Y_2 - Y_1$. Следователно:
\[
X = h(Y) = \begin{pmatrix} Y_1 \\ Y_2 - Y_1 \end{pmatrix}
\]
Якобианата е $J(y) = \det \begin{pmatrix} 1 & 0 \\ -1 & 1 \end{pmatrix} = 1$.
Съвместната плътност на $Y$ става:
\[
f_Y(y_1, y_2) = f_X(h(y_1, y_2)) \cdot 1 = f_{X_1}(y_1) f_{X_2}(y_2 - y_1)
\]
За да намерим търсената маргинална плътност на сумата $Y_2$, интегрираме по $y_1$:

\begin{defn}[конволюция]\label{def:convolution}
Плътността на сумата на две независими непрекъснати случайни величини се дава от \emph{конволюцията}
\[
f_{Y_2}(y_2) = \int f_{X_1}(y_1) f_{X_2}(y_2 - y_1) \, dy_1 .
\]
\end{defn}

Тънкостта при прилагането ѝ е правилното определяне на границите на интеграла, тъй като те зависят от дефиниционните области на изходните величини (виж следващия пример).

\section{Гама разпределение и Хи-квадрат разпределение}

\subsection{Гама разпределение}

Казваме, че $X$ има Гама разпределение с параметри $\alpha > 0$ и $\beta > 0$, записвано $X \sim \Gamma(\alpha, \beta)$, ако плътността ѝ е:
\[
f_X(x) = 
\begin{cases}
\frac{\beta^\alpha}{\Gamma(\alpha)} x^{\alpha - 1} e^{-\beta x}, & \text{за } x \ge 0 \\
0, & \text{иначе}
\end{cases}
\]
където $\Gamma(\alpha) = \int_0^\infty x^{\alpha-1} e^{-x} dx$ е добре познатата Гама функция. Важни нейни стойности са $\Gamma(1) = 1$ и $\Gamma(n+1) = n!$.

Забележете връзката с експоненциалното разпределение: ако фиксираме $\alpha = 1$, плътността става $f_X(x) = \frac{\beta^1}{1} x^{0} e^{-\beta x} = \beta e^{-\beta x}$. Това означава, че $\Gamma(1, \beta) \equiv \text{Exp}(\beta)$.

\begin{prop}[Сума на Гама величини]
Нека $X_1, \dots, X_n$ са независими в съвкупност случайни величини, като $X_i \sim \Gamma(\alpha_i, \beta)$ (имат общ мащабен параметър $\beta$). Тогава тяхната сума също има Гама разпределение:
\[
Y = \sum_{j=1}^n X_j \sim \Gamma\left(\sum_{i=1}^n \alpha_i, \beta\right)
\]
\end{prop}

\begin{proof}[Доказателство (за $n=2$)]
Използваме конволюционната схема от предишния раздел. Търсим плътността на $Y_2 = X_1 + X_2$. Тъй като $X_1 \ge 0$ и $X_2 \ge 0$, то $y_1 \ge 0$ и $y_2 - y_1 \ge 0 \implies 0 \le y_1 \le y_2$. Границите на интеграла са от $0$ до $y_2$:
\[
f_{Y_2}(y_2) = \int_{0}^{y_2} f_{X_1}(y_1) f_{X_2}(y_2 - y_1) dy_1
\]
Заместваме плътностите на двете Гама разпределения:
\[
f_{Y_2}(y_2) = \int_{0}^{y_2} \left( \frac{\beta^{\alpha_1}}{\Gamma(\alpha_1)} y_1^{\alpha_1 - 1} e^{-\beta y_1} \right) \left( \frac{\beta^{\alpha_2}}{\Gamma(\alpha_2)} (y_2 - y_1)^{\alpha_2 - 1} e^{-\beta (y_2 - y_1)} \right) dy_1
\]
Експоненциалните членове се опростяват: $e^{-\beta y_1} e^{-\beta y_2 + \beta y_1} = e^{-\beta y_2}$, което не зависи от $y_1$ и излиза пред интеграла. Остава да се реши:
\[
f_{Y_2}(y_2) = \frac{\beta^{\alpha_1+\alpha_2}}{\Gamma(\alpha_1)\Gamma(\alpha_2)} e^{-\beta y_2} \int_{0}^{y_2} y_1^{\alpha_1 - 1} (y_2 - y_1)^{\alpha_2 - 1} dy_1
\]
С помощта на субституцията $y_1 = u y_2$, интегралът се свежда до Бета функцията $B(\alpha_1, \alpha_2) = \frac{\Gamma(\alpha_1)\Gamma(\alpha_2)}{\Gamma(\alpha_1+\alpha_2)}$ и крайният резултат е точно плътността на $\Gamma(\alpha_1+\alpha_2, \beta)$:
\[
f_{Y_2}(y_2) = \frac{\beta^{\alpha_1+\alpha_2}}{\Gamma(\alpha_1+\alpha_2)} y_2^{\alpha_1+\alpha_2-1} e^{-\beta y_2}
\]
От това следва, че Гама разпределението е \textbf{безгранично делимо}. Например $\Gamma(3{,}5, \beta)$ може да се представи като сума от 7 независими $\Gamma(0{,}5, \beta)$ величини.
\end{proof}

\subsection{Хи-квадрат разпределение ($\chi^2$)}

\begin{prop}
Нека $Z_1, \dots, Z_n$ са независими в съвкупност стандартно нормални случайни величини ($Z_i \sim N(0, 1)$). Тогава сумата от техните квадрати има $\chi^2$ разпределение с $n$ степени на свобода, което е еквивалентно на Гама разпределение:
\[
Y = \sum_{j=1}^n Z_j^2 \sim \chi^2(n) \equiv \Gamma\left(\frac{n}{2}, \frac{1}{2}\right)
\]
\end{prop}

\begin{proof}
Ще покажем, че квадратът на една стандартно нормална величина $T_1 = Z_1^2$ има разпределение $\chi^2(1) \equiv \Gamma(1/2, 1/2)$. 
Започваме с функцията на разпределение за $t > 0$:
\[
F_{T_1}(t) = \mathbb{P}(T_1 < t) = \mathbb{P}(Z_1^2 < t) = \mathbb{P}(-\sqrt{t} < Z_1 < \sqrt{t}) = \frac{1}{\sqrt{2\pi}} \int_{-\sqrt{t}}^{\sqrt{t}} e^{-\frac{y^2}{2}} dy
\]
Поради четността на функцията, това е:
\[
F_{T_1}(t) = 2 \cdot \frac{1}{\sqrt{2\pi}} \int_{0}^{\sqrt{t}} e^{-\frac{y^2}{2}} dy
\]
Диференцираме по $t$, за да намерим плътността (използвайки правилото за диференциране на интеграл с променлива горна граница $\sqrt{t}$, чиято производна е $\frac{1}{2\sqrt{t}}$):
\[
f_{T_1}(t) = \frac{d}{dt} F_{T_1}(t) = \frac{2}{\sqrt{2\pi}} e^{-\frac{(\sqrt{t})^2}{2}} \cdot \frac{1}{2\sqrt{t}} = \frac{1}{\sqrt{2\pi}} t^{-1/2} e^{-t/2}
\]

Сега ще сравним този резултат с плътността на Гама разпределението при параметри $\alpha = 1/2$ и $\beta = 1/2$:
\[
f_\Gamma(t) = \frac{(1/2)^{1/2}}{\Gamma(1/2)} t^{(1/2) - 1} e^{-t/2} = \frac{1}{\sqrt{2}\Gamma(1/2)} t^{-1/2} e^{-t/2}
\]
Тъй като и двата израза представляват плътност (и следователно интегралите им са 1), константите отпред трябва да са равни. От равенството $\frac{1}{\sqrt{2\pi}} = \frac{1}{\sqrt{2}\Gamma(1/2)}$ можем пътем да изчислим красивия факт, че $\Gamma(1/2) = \sqrt{\pi}$.
Най-важното е, че двете плътности съвпадат. Следователно $Z_1^2 \sim \Gamma(1/2, 1/2)$.

Тъй като $Y = \sum_{j=1}^n Z_j^2$ е сума на $n$ независими Гама величини с еднакъв мащабен параметър $\beta = 1/2$ и формов параметър $\alpha_j = 1/2$, от предходното твърдение следва директно, че:
\[
Y \sim \Gamma\left(\sum_{j=1}^n \frac{1}{2}, \frac{1}{2}\right) = \Gamma\left(\frac{n}{2}, \frac{1}{2}\right) \equiv \chi^2(n)
\]
С това доказателството е завършено.
\end{proof}


\chapter{Видове сходимост. Неравенство на Чебишов и закони за големите числа}
\section{Въведение. Видове сходимости на случайни величини}

Законите за големите числа са едни от най-фундаменталните резултати в теорията на вероятностите. Често пъти, когато използваме различни методи, компютърни пакети или статистически анализи (като например анализ на данни), обосновката за валидността на тези методи произтича именно от този фундамент. За тяхното формулиране и разбиране първо е необходимо понятието за сходимост в теорията на вероятностите.

Ще разгледаме три основни типа сходимост на случайни величини. Съществуват и други видове, но тези три са най-важните за нашите цели. За първите два типа сходимост ще допуснем, че имаме редица от случайни величини $X_n$ (за $n \ge 1$) и една гранична случайна величина $X$. Важно е да отбележим, че за тези първи две дефиниции всички случайни величини живеят в едно и също вероятностно пространство $(\Omega, \mathcal{A}, \mathbb{P})$. Тоест, имаме едно общо множество от елементарни събития $\Omega$, една сигма-алгебра $\mathcal{A}$ и една вероятностна мярка $\mathbb{P}$.

\subsection{Сходимост почти сигурно (п.с.)}

За удобство, нека означим с $L_X$ множеството от всички елементарни събития $\omega \in \Omega$, за които числовата редица $X_n(\omega)$ се схожда към числото $X(\omega)$ в класическия смисъл на математическия анализ:
\[
L_X = \left\{ \omega \in \Omega \mathrel{\bigg|} \lim_{n \to \infty} X_n(\omega) = X(\omega) \right\}
\]
За всяко конкретно $\omega$, $X_n(\omega)$ е просто една редица от реални числа. Ако тази числова редица клони към $X(\omega)$, то събитието $\omega$ принадлежи на $L_X$.

\begin{defn}[Сходимост почти сигурно]
Казваме, че редицата $X_n$ се схожда почти сигурно към $X$, което се записва като $X_n \xrightarrow{\text{п.с.}} X$, ако вероятността на събитието $L_X$ е равна на единица:
\[
\mathbb{P}(L_X) = \mathbb{P}\left( \lim_{n \to \infty} X_n = X \right) = 1
\]
\end{defn}

Интуитивно това означава, че ако изберете случайно едно елементарно събитие (т.е. проведете експеримент), вероятността да изберете такова $\omega$, че редицата да се сходи, е $1$. Това не означава задължително, че сходимостта е налице за абсолютно всяко възможно $\omega$ — множеството от изключения (където няма сходимост) може да не е празно, но то задължително има вероятност (мярка) 0.

\subsection{Сходимост по вероятност}

За втория тип сходимост нека въведем събитието $A_{n, \varepsilon}$, което се състои от всички елементарни събития, за които разликата между $X_n$ и $X$ надвишава дадено число $\varepsilon > 0$:
\[
A_{n, \varepsilon} = \{ \omega \in \Omega \mathrel{\big|} |X_n(\omega) - X(\omega)| > \varepsilon \}
\]
Това са онези изходи от експеримента, при които стойността на $X_n$ стои на разстояние по-голямо от $\varepsilon$ спрямо стойността на $X$.

\begin{defn}[Сходимост по вероятност]
Казваме, че редицата $X_n$ се схожда към $X$ по вероятност (означава се с $X_n \xrightarrow{\mathbb{P}} X$), ако за всяко $\varepsilon > 0$ е изпълнено:
\[
\lim_{n \to \infty} \mathbb{P}(|X_n - X| > \varepsilon) = 0
\]
\end{defn}

Това означава, че при достатъчно голямо $n$, вероятността разстоянието между $X_n$ и $X$ да бъде съществено (повече от $\varepsilon$) става пренебрежимо малка. Вероятността не е задължително да намалява монотонно, но в граница тя клони към нула.

\emph{Забележка:} В горната дефиниция е напълно достатъчно да се разглеждат само $\varepsilon$ от вида $\frac{1}{r}$ за всяко естествено число $r \ge 1$. Ако $\varepsilon$ е произволно реално число, винаги можем да намерим естествено $r$, такова че $\frac{1}{r+1} < \varepsilon \le \frac{1}{r}$. Тъй като условието $|X_n - X| > \varepsilon$ влече $|X_n - X| > \frac{1}{r+1}$, събитията се влагат едно в друго (т.е. $A_{n, \varepsilon} \subseteq A_{n, \frac{1}{r+1}}$). Работата с изброимо количество стойности на $\varepsilon$ е много полезна в теоретичните доказателства.

\subsection{Сходимост по разпределение}

При предишните два вида сходимост беше ключово всички случайни величини да са дефинирани в едно и също вероятностно пространство $\Omega$. Сходимостта по разпределение обаче е най-слабият, но може би най-полезният тип сходимост в практиката, тъй като не зависи от конкретното пространство. Тя се интересува единствено от вероятностните свойства (разпределенията) на величините.

Нека $F_X(x)$ е функцията на разпределение на случайната величина $X$. Ще означим с $C_{F_X}$ множеството от всички реални числа $x \in \mathbb{R}$, в които функцията $F_X$ е непрекъсната. За непрекъснати случайни величини (с плътност), $C_{F_X} = \mathbb{R}$. Ако обаче $X$ има скокове (дискретна случайна величина, напр. на Бернули), $C_{F_X}$ е цялата реална права без точките на скок. Точките на скок са точно тези $x$, за които $\mathbb{P}(X = x) > 0$.

\begin{defn}[Сходимост по разпределение]
Нека $X_n$ е случайна величина в пространство $(\Omega_n, \mathcal{A}_n, \mathbb{P}_n)$ за $n \ge 1$, а $X$ е случайна величина в пространство $(\Omega, \mathcal{A}, \mathbb{P})$. Казваме, че $X_n$ се схожда по разпределение към $X$ (записваме $X_n \xrightarrow{d} X$), ако за всяко $x \in C_{F_X}$ е вярно, че:
\[
\lim_{n \to \infty} F_{X_n}(x) = F_X(x)
\]
\end{defn}

Този вид сходимост означава, че въпреки че експериментите може да са коренно различни (едното изследване може да е върху хора, друго върху животни, трето хвърляне на монети), вероятностната структура на случайните величини става все по-близка. Това е сходимост на самите функции на разпределение, а не сходимост на реализациите на величините.

\section{Връзка между видовете сходимост}

Нека сега градираме видовете сходимост по тяхната „сила“.

\begin{thm}
\begin{enumerate}
    \item[а)] Ако $X_n \xrightarrow{\text{п.с.}} X$, то $X_n \xrightarrow{\mathbb{P}} X$.
    \item[б)] Ако $X_n \xrightarrow{\mathbb{P}} X$, то $X_n \xrightarrow{d} X$.
\end{enumerate}
\end{thm}

Най-силната сходимост е „почти сигурно“, следвана от „по вероятност“, и най-слабата е „по разпределение“. Поради това сходимостта по разпределение често се постига най-лесно.

Обратните твърдения по принцип не са верни, и двата контрапримера си заслужават да се видят.

\begin{example}[По вероятност, но не почти сигурно]\label{ex:10-typewriter}
Работим върху $\Omega = [0,1]$ с $\mathbb{P}$ равна на дължината. Разглеждаме индикаторите на
последователни интервали, които обикалят отрезъка:
\[
\xi_{n,k} = \ind_{\left[\frac{k-1}{n},\, \frac{k}{n}\right]},
\qquad k = 1, \dots, n, \quad n = 1, 2, \dots
\]
и ги подреждаме в една редица: $\xi_{1,1};\ \xi_{2,1}, \xi_{2,2};\ \xi_{3,1}, \xi_{3,2}, \xi_{3,3};\ \dots$

Тази редица клони към $0$ \emph{по вероятност}: за фиксирано $\varepsilon \in (0,1)$
събитието $\{\xi_{n,k} > \varepsilon\}$ е точно интервалът с дължина $\frac1n$, тоест
неговата вероятност е $\frac1n \to 0$.

Но никъде не клони \emph{почти сигурно}: за всяко фиксирано $\omega \in [0,1]$ и за всяко
$n$ съществува точно едно $k$, за което $\omega$ попада в съответния интервал, тоест
$\xi_{n,k}(\omega) = 1$. Значи редицата приема стойност $1$ безкрайно често и стойност $0$
безкрайно често — тя не се схожда за нито едно $\omega$.
\end{example}

\begin{example}[По разпределение, но не по вероятност]\label{ex:10-samedist}
Нека $\xi, \xi_1, \xi_2, \dots$ са независими и еднакво разпределени, всички $N(0,1)$.
Понеже всички имат една и съща функция на разпределение, тривиално
$\xi_n \xrightarrow{d} \xi$.

Сходимост по вероятност обаче няма. Разликата $\xi_n - \xi$ е сума на две независими
стандартно нормални величини (с обратен знак), тоест $\xi_n - \xi \sim N(0,2)$ за всяко $n$.
Тогава $\mathbb{P}(|\xi_n - \xi| > \varepsilon)$ е една и съща положителна константа за всяко
$n$ и не клони към нула.

Този пример показва и защо сходимостта по разпределение не се интересува от
вероятностното пространство: тя гледа само разпределенията, не и самите реализации.
\end{example}

\begin{proof}[Скица на доказателството на част а)]
Знаем, че $\mathbb{P}(L_X) = 1$, където $L_X = \left\{ \lim_{n\to\infty} X_n = X \right\}$. По дефиниция за сходимост на числова редица, множеството $L_X$ може да бъде представено чрез сечения и обединения:
\[
L_X = \bigcap_{r=1}^{\infty} \bigcup_{n=1}^{\infty} \bigcap_{k=n}^{\infty} A_{k, \frac{1}{r}}^{c}
\]
където $A_{k, \frac{1}{r}}^{c} = \left\{ |X_k - X| \le \frac{1}{r} \right\}$. Изразът отразява факта, че „за всяко $r$ (аналог на $\varepsilon$), съществува $n$, такова че за всяко $k \ge n$, отклонението $|X_k - X|$ е не повече от $\frac{1}{r}$“.

Взимайки допълнението по правилата на Де Морган, получаваме множеството, в което няма сходимост:
\[
L_X^c = \bigcup_{r=1}^{\infty} \underbrace{\bigcap_{n=1}^{\infty} \bigcup_{k=n}^{\infty} A_{k, \frac{1}{r}}}_{B_r}
\]
Тъй като $\mathbb{P}(L_X) = 1$, то $\mathbb{P}(L_X^c) = 0$. От това следва, че:
\[
0 = \mathbb{P}\left(\bigcup_{r=1}^{\infty} B_r\right) \ge \mathbb{P}(B_\ell) \quad \forall \ell \ge 1
\]
Следователно $\mathbb{P}(B_r) = 0$ за всяко $r \ge 1$.

Нека дефинираме $C_{n, r} = \bigcup_{k=n}^{\infty} A_{k, \frac{1}{r}}$. Забелязваме, че с нарастването на $n$, обединението включва все по-малко множества, така че редицата е намаляваща: $C_{n, r} \supseteq C_{n+1, r} \supseteq \dots$
Тъй като $B_r = \bigcap_{n=1}^{\infty} C_{n, r}$ и редицата е монотонно намаляваща, от непрекъснатостта на вероятностната мярка следва:
\[
\mathbb{P}(B_r) = 0 \implies \lim_{n\to\infty} \mathbb{P}(C_{n, r}) = 0
\]
Но $C_{n, r}$ очевидно съдържа първия си елемент $A_{n, \frac{1}{r}} = \left\{ |X_n - X| > \frac{1}{r} \right\}$. Щом вероятността на по-голямото множество $C_{n, r}$ клони към 0, тогава и:
\[
\lim_{n\to\infty} \mathbb{P}\left(|X_n - X| > \frac{1}{r}\right) = 0
\]
Това е вярно за всяко $r$, което доказва, че $X_n$ клони към $X$ по вероятност.

\end{proof}

\begin{proof}[Скица на доказателството на част б)]
Тук искаме да свържем събитието $\{X_n < x\}$, което участва във функцията на разпределение
на $X_n$, със събитието $\{X < x\}$. Директна връзка няма — трябва да минем през малка
пертурбация $\varepsilon$. Означаваме
\[
A_{n,\varepsilon} = \{ |X_n - X| > \varepsilon \},
\qquad\text{като}\qquad
\mathbb{P}(A_{n,\varepsilon}) \xrightarrow[n \to \infty]{} 0
\]
за всяко $\varepsilon > 0$ по условие (това е точно сходимостта по вероятност).

Разбиваме $\{X_n < x\}$ по $A_{n,\varepsilon}$ и допълнението му, което не променя нищо:
\[
\{X_n < x\} = \big(\{X_n < x\} \cap A_{n,\varepsilon}^c\big) \cup \big(\{X_n < x\} \cap A_{n,\varepsilon}\big).
\]
Върху $A_{n,\varepsilon}^c$ имаме $|X_n - X| \le \varepsilon$, тоест $X$ се различава от $X_n$ с
най-много $\varepsilon$. Затова
\[
\{X_n < x\} \subseteq \{X < x + \varepsilon\} \cup A_{n,\varepsilon},
\qquad
\{X < x - \varepsilon\} \cap A_{n,\varepsilon}^c \subseteq \{X_n < x\}.
\]
Прилагайки вероятност към двете влагания, получаваме „сандвич“:
\[
F_X(x-\varepsilon) - \mathbb{P}(A_{n,\varepsilon})
\ \le\ F_{X_n}(x) \ \le\
F_X(x+\varepsilon) + \mathbb{P}(A_{n,\varepsilon}).
\]
При $n \to \infty$ членовете $\mathbb{P}(A_{n,\varepsilon})$ отпадат и остава
\[
F_X(x-\varepsilon) \ \le\ \liminf_n F_{X_n}(x) \ \le\ \limsup_n F_{X_n}(x) \ \le\ F_X(x+\varepsilon).
\]
Накрая пускаме $\varepsilon \to 0$. Ако $x$ е точка на непрекъснатост на $F_X$, двата края
клонят към $F_X(x)$ и следователно $F_{X_n}(x) \to F_X(x)$ — точно това е сходимост по
разпределение.

\end{proof}

\subsection{Обратна връзка в частни случаи}
В общия случай от сходимост по разпределение не следва сходимост по вероятност. Съществува обаче едно важно изключение: когато граничната случайна величина е константа.

\begin{prop}
Ако $X_n \xrightarrow{d} c$, където $c$ е константа, то $X_n \xrightarrow{\mathbb{P}} c$ (разбира се, допускайки, че случайните величини са дефинирани в едно и също вероятностно пространство).
\end{prop}

\emph{Идея на доказателството:} При използваната тук конвенция функцията на разпределение на константата $c$ е $F_c(x) = 0$ за $x \le c$ и $F_c(x) = 1$ за $x > c$. Единствената точка на прекъсване е $x=c$. Тъй като $X_n \xrightarrow{d} c$, за всяко $x \neq c$ имаме сходимост на функциите на разпределение. От това може да се покаже, че вероятността $X_n$ да се отклони от $c$ с повече от $\varepsilon$ се състои от сумата на вероятностите $\mathbb{P}(X_n > c + \varepsilon)$ и $\mathbb{P}(X_n < c - \varepsilon)$. И двете клонят към 0, тъй като в точките $c+\varepsilon$ и $c-\varepsilon$ функцията на разпределение $F_c$ е съответно 1 и 0, към които се схождат $F_{X_n}$.

\section{Неравенство на Чебишов}

Когато въведохме дисперсията, обсъдихме, че тя играе ролята на средноквадратична грешка. Неравенството на Чебишов отговаря на въпроса: „Ако знаем нещо за дисперсията, каква е горната граница на вероятността разликата между $X$ и очакването на $X$ да е по-голяма от някакво число $a$?“ Това ни позволява лесно да оценим вероятността за големи отклонения.

\begin{prop}[Неравенство на Марков]\label{prop:markov}
Ако $Y \ge 0$ е неотрицателна случайна величина с крайно очакване, то
\[ \mathbb{P}(Y > a) \le \frac{\E Y}{a} \quad \forall a > 0 . \]
\end{prop}

\begin{proof}
Понеже $Y \ge 0$, за всяко $a>0$ е в сила поточковото неравенство
$Y \ge a\,\ind_{\{Y > a\}}$: върху събитието $\{Y>a\}$ лявата страна е по-голяма от $a$, а
извън него дясната страна е нула. Взимаме очакване и използваме монотонността му заедно с
$\E \ind_{\{Y>a\}} = \mathbb{P}(Y>a)$:
\[ \E Y \ge a\,\mathbb{P}(Y > a) . \]
\end{proof}

Неравенството на Чебишов се получава от него с едно заместване: прилагаме
твърдение~\ref{prop:markov} към $Y = (X - \E X)^2 \ge 0$ и към прага $a^2$. По-долу е
дадено и прякото доказателство, както беше направено на лекцията.

\begin{keythm}[Неравенство на Чебишов]
Нека $X$ е случайна величина с дисперсия $\Var X$. Тогава:
\[
\mathbb{P}(|X - \E X| > a) \le \frac{\Var X}{a^2} \quad \forall a > 0
\]
\end{keythm}

\begin{proof}
Нека положим $Y = X - \E X$. Търсим да оценим $\mathbb{P}(|Y| > a)$, което може да се запише като очакване на индикаторната функция: $\E[\ind_{\{|Y| > a\}}]$.
От свойствата на дисперсията имаме $\Var X = \Var Y = \E[Y^2]$ (тъй като $\E Y = 0$).

Представяме единицата като сума от два индикатора: $1 = \ind_{\{|Y| \le a\}} + \ind_{\{|Y| > a\}}$. Тогава:
\[
\E[Y^2] = \E[Y^2 \ind_{\{|Y| \le a\}}] + \E[Y^2 \ind_{\{|Y| > a\}}]
\]
Тъй като и двете събираеми са неотрицателни, можем да премахнем първото и да получим неравенство:
\[
\E[Y^2] \ge \E[Y^2 \ind_{\{|Y| > a\}}]
\]
Върху събитието, за което индикаторът $\ind_{\{|Y| > a\}}$ не е нула (тоест когато $|Y| > a$), със сигурност е изпълнено $Y^2 > a^2$. Затова можем да заменим $Y^2$ с $a^2$:
\[
\E[Y^2 \ind_{\{|Y| > a\}}] \ge \E[a^2 \ind_{\{|Y| > a\}}] = a^2 \E[\ind_{\{|Y| > a\}}] = a^2 \mathbb{P}(|Y| > a)
\]
Окончателно получаваме $\Var X \ge a^2 \mathbb{P}(|Y| > a)$, или:
\[
\mathbb{P}(|X - \E X| > a) \le \frac{\Var X}{a^2}
\]
\end{proof}

\subsection{Следствия и пример}
Аналогично неравенство за произволни моменти $m \ge 1$ се получава, като приложим твърдение~\ref{prop:markov} към $Y = |X - \E X|^m$ и прага $a^m$ (названието Марков или Чебишов тук не се фиксира, тъй като няма съществена разлика в типа доказателство):
\[
\mathbb{P}(|X - \E X| > a) \le \frac{\E[|X - \E X|^m]}{a^m}
\]
Когато $m=2$, получаваме класическото неравенство на Чебишов.

\begin{example}\label{ex:10-1}
Каква е вероятността една случайна величина $X$ да се отклони от очакването си с повече от $10$ пъти своето стандартно отклонение?
Нека за удобство $\E X = 0$. Търсим $\mathbb{P}(|X| > 10\sqrt{\Var X})$. В неравенството на Чебишов избираме $a = 10\sqrt{\Var X}$. Получаваме:
\[
\mathbb{P}(|X| > 10\sqrt{\Var X}) \le \frac{\Var X}{(10\sqrt{\Var X})^2} = \frac{\Var X}{100 \Var X} = \frac{1}{100}
\]
Това е изключително силен универсален резултат — за \emph{всяка} случайна величина вероятността да надвиши 10 пъти стандартното си отклонение никога не надхвърля $1\%$.

\end{example}

\section{Закони за големите числа (ЗГЧ)}

Законите за големите числа формализират интуитивната идея, че при извършване на голям брой независими експерименти, средноаритметичното от резултатите се приближава до теоретичното математическо очакване.

\subsection{Слаб закон за големите числа}

\begin{keythm}[Слаб ЗГЧ]
Нека $(X_i)_{i=1}^\infty$ е редица от независими и еднакво разпределени случайни величини, такива че $\E|X_1| < \infty$. Тогава за редицата е в сила законът за големите числа, тоест:
\[
\frac{\sum_{i=1}^n X_i}{n} \xrightarrow[n \to \infty]{\mathbb{P}} \E X_1
\]
\emph{Бележка:} Тъй като $X_i$ са еднакво разпределени, всички те имат едно и също математическо очакване, $\E X_i = \E X_1$.
\end{keythm}

\begin{proof}[Доказателство (в случая, когато дисперсията съществува)]
Ще докажем теоремата при допълнителното по-силно условие, че $\Var X_1 < \infty$, тъй като тогава можем директно да използваме неравенството на Чебишов.
Трябва да покажем, че за всяко $\varepsilon > 0$:
\[
\lim_{n \to \infty} \mathbb{P} \left( \left| \frac{\sum_{i=1}^n X_i}{n} - \E X_1 \right| > \varepsilon \right) = 0
\]
Тъй като $\E X_i = \E X_1$ за всяко $i$, можем да внесем очакването вътре в сумата:
\[
\mathbb{P} \left( \left| \frac{\sum_{i=1}^n X_i}{n} - \frac{\sum_{i=1}^n \E X_i}{n} \right| > \varepsilon \right) = \mathbb{P} \left( \frac{\left| \sum_{i=1}^n (X_i - \E X_i) \right|}{n} > \varepsilon \right)
\]
Като умножим двете страни на неравенството в скобите по $n$, получаваме:
\[
\mathbb{P} \left( \left| \sum_{i=1}^n (X_i - \E X_i) \right| > \varepsilon n \right)
\]
Сега прилагаме неравенството на Чебишов за случайната величина $Y = \sum_{i=1}^n (X_i - \E X_i)$, чието очакване е 0. Като заместим $a = \varepsilon n$, имаме:
\[
\mathbb{P}(|Y| > \varepsilon n) \le \frac{\Var \left( \sum_{i=1}^n (X_i - \E X_i) \right)}{(\varepsilon n)^2}
\]
Тук ключово използваме \emph{независимостта} на случайните величини $X_i$. Поради нея, дисперсията на сумата е равна на сумата от дисперсиите:
\[
\frac{\sum_{i=1}^n \Var(X_i)}{\varepsilon^2 n^2}
\]
Но тъй като величините са еднакво разпределени, те имат еднаква дисперсия $\Var X_1$. Сумата се състои от $n$ такива събираеми:
\[
= \frac{n \Var X_1}{n^2 \varepsilon^2} = \frac{\Var X_1}{n \varepsilon^2}
\]
Тъй като дисперсията $\Var X_1$ и $\varepsilon^2$ са фиксирани константи, при $n \to \infty$ този израз клони към $0$. С това доказахме сходимостта по вероятност.
\end{proof}

\subsection{Усилен закон за големите числа (УЗГЧ)}

Слабият закон гарантира сходимост само по вероятност. Усиленият закон утвърждава много по-силен резултат — сходимост почти сигурно.

\begin{keythm}[Усилен закон за големите числа]\label{thm:slln}
Нека $(X_i)_{i=1}^\infty$ е редица от независими и еднакво разпределени случайни величини с крайни очаквания $\E|X_1| < \infty$. Тогава средното аритметично се схожда почти сигурно към теоретичното очакване:
\[
\frac{\sum_{i=1}^n X_i}{n} \xrightarrow[n \to \infty]{\text{п.с.}} \E X_1
\]
Това е най-полезният и често срещан вид на закона, тъй като в повечето практически задачи (като например събиране на емпирични данни) работим с независими наблюдения от един и същи феномен.
\end{keythm}

\subsection{Примери и приложения на ЗГЧ}

\begin{example}[Схема на Бернули: гласуване или ваксини]\label{ex:10-2}
Нека $X_i \sim \Ber(p)$ с $p \in (0,1)$. Например, $X_i = 1$, ако $i$-тият човек гласува за партия А (или ако ваксината проработи при него) и $0$ в противен случай. Не знаем истинската вероятност $p$, но чрез УЗГЧ знаем, че:
\[
\frac{\sum_{i=1}^n X_i}{n} \xrightarrow[n \to \infty]{\text{п.с.}} p = \E X_1
\]
Тоест, ако попитаме достатъчно голям брой хора, пропорцията от гласове (средноаритметичното на $X_i$) ще клони почти сигурно към истинската вероятност $p$.
\end{example}

\begin{example}[Проверка на грамаж]\label{ex:10-3}
Ако на етикета пише, че в един сандвич има 50 грама шунка, можем да измерим грамажа на голям брой $n$ сандвичи. Нека $X_i$ е количеството шунка в $i$-тия сандвич. Средното аритметично $\frac{1}{n}\sum X_i$ трябва да клони към теоретичното очакване $\E X_1$. Ако измерим средно 46 грама при достатъчно голямо $n$, по силата на ЗГЧ имаме сериозни основания да се усъмним, че истинското очакване е 50 грама.
\end{example}

\begin{example}[Игра на рулетка]\label{ex:10-4}
Връщаме се към \crosslecture{ex:05-3}{примера с рулетката от лекция 5}. Печелите 1 лев с вероятност $\frac{18}{37}$ и губите 1 лев с вероятност $\frac{19}{37}$. Очакваната ви печалба от една игра е:
\[
\E X_1 = 1 \cdot \frac{18}{37} - 1 \cdot \frac{19}{37} = -\frac{1}{37}
\]
Колкото повече играете, толкова средната ви печалба на игра клони към:
\[
\frac{\sum_{i=1}^n X_i}{n} \xrightarrow[n \to \infty]{\text{п.с.}} -\frac{1}{37}
\]
което обяснява защо дългосрочно винаги сте на загуба.
\end{example}

\begin{example}[Дърпане на въже: случайно блуждаене]\label{ex:10-5}\label{ex:tug-of-war}
Връщаме се към \crosslecture{ex:05-4}{модела със случайно дърпане на въже от лекция 5}. Имаме $n$ играчи (или молекули), които дърпат въжето наляво или надясно с вероятност $\frac{1}{2}$. Силата $X_i = 1$ за надясно и $-1$ за наляво. Очакваната сила е $\E X_i = 0$. Резултантната сила е $S_n = \sum_{i=1}^n X_i$. Според ЗГЧ, средната сила $\frac{S_n}{n}$ клони към 0. Важно е обаче да не се бърка сходимостта на средната сила с вероятността дясната страна да победи ($\mathbb{P}(S_n > 0)$). Пропорцията $\frac{S_n}{n}$ клони към 0, но вероятността резултантната сила да е положителна клони към $\frac{1}{2}$, а не към 0! Това е често допускана грешка от неразбиране на разликата между самите пропорции и вероятностите за техните стойности.
\end{example}

\begin{example}[Симулации Монте Карло]\label{ex:10-6}\label{ex:monte-carlo}
Продължаваме \crosslecture{ex:monte-carlo-volume}{геометричната интерпретация на Монте Карло от лекция 2}. Искаме да намерим обема $|A|$ на дадено множество $A$, вложено в единичния куб. Алгоритъмът Монте Карло „хвърля“ случайно и равномерно точки $X_i$ в единичния куб. Дефинираме $Y_i = 1$, ако $X_i \in A$, и $Y_i = 0$, ако $X_i \notin A$. Величината $Y_i$ е Бернулиева с параметър $|A|$ (тъй като обемът на единичния куб е 1). Според УЗГЧ:
\[
\frac{\sum_{i=1}^n Y_i}{n} \xrightarrow[n \to \infty]{\text{п.с.}} |A|
\]
Тоест, просто като броим точките, попаднали в $A$, и ги разделим на общия брой хвърлени точки, ние числено апроксимираме търсения обем $|A|$. Колко точно точки трябва да хвърлим, за да постигнем определена точност, е въпрос, на който отговаря Централната Гранична Теорема, която ще разгледаме в следващите лекции.
\end{example}

\section{Задачи}

\begin{enumerate}
    \item Докажете, че в дефиницията за сходимост по вероятност е достатъчно условието да се проверява за $\varepsilon=1/r$, $r\in\mathbb{N}$.
    \item Довършете доказателството, че от $X_n\xrightarrow{\mathbb{P}}X$ следва $X_n\xrightarrow{d}X$.
    \item Довършете доказателството на частния обратен резултат: ако $X_n\xrightarrow{d}c$ за константа $c$, то $X_n\xrightarrow{\mathbb{P}}c$.
    \item Нека $\E X=0$ и $\Var X<\infty$. Използвайте неравенството на Чебишов, за да докажете универсалната оценка
    \[
    \mathbb{P}\left(|X|>10\sqrt{\Var X}\right)\le\frac{1}{100}.
    \]
\end{enumerate}


\chapter{Централна гранична теорема и функции на моментите}
\section{От усиления закон към скоростта на сходимост}

За да поставим новия въпрос за скоростта на сходимост, първо припомняме означенията от усиления закон. Разглеждаме редица от независими и еднакво разпределени (н.е.р.) случайни величини $(X_i)_{i=1}^\infty$. Фактът, че са взаимно независими и еднакво разпределени, означава, че всяка една от тях има същата функция на разпределение като всяка друга, тоест функцията на разпределение $F(x)$ е една и съща за всяко $X_i$.

Нека $\mu = \E[X_1]$ бъде математическото очакване на първата случайна величина. Тъй като величините са еднакво разпределени, $\mu$ е математическото очакване на всички останали величини в редицата. Дефинираме парциалната сума $S_n$:
\[ S_n = \sum_{j=1}^n X_j \]
$S_n$ представлява нова случайна величина за всяко $n$. Редицата от случайни величини $S_1, S_2, \dots$ разбира се не е нито от независими, нито от еднакво разпределени величини, тъй като те натрупват стойностите на първите $n$ елемента.

Съгласно \crosslecture{thm:slln}{теоремата за усиления закон за големите числа от предходната лекция}, средното аритметично $\frac{S_n}{n}$ се схожда почти сигурно (п.с.) към математическото очакване $\mu$:
\[ \frac{S_n}{n} \xrightarrow[n \to \infty]{\text{п.с.}} \mu \]
Това означава, че ако извършим серия от експерименти, средното аритметично от резултатите ще се доближава неограничено до очакването на първата случайна величина с вероятност единица. Ако означим разликата (или грешката) с $E_n = \frac{S_n}{n} - \mu$, то е ясно, че $E_n$ клони към нула почти сигурно при $n \to \infty$.

Въпросът обаче е: \emph{колко бързо} се осъществява тази сходимост? Можем ли да кажем нещо допълнително за поведението на грешката $E_n = \frac{S_n}{n} - \mu$? Знаем, че клони към нула, но това е единствената информация дотук. За да отговорим на този въпрос и да изследваме поведението на тази грешка, се нуждаем от Централната гранична теорема (ЦГТ).

\section{Централна гранична теорема (ЦГТ)}

Нека формулираме теоремата, която описва прецизно поведението на тази разлика при голямо $n$.

\begin{keythm}[Централна гранична теорема]\label{thm:clt}
Нека $(X_i)_{i=1}^\infty$ са независими и еднакво разпределени (н.е.р.) случайни величини със средно $\mu = \E[X_1]$ и дисперсия $\Var[X_1] = \sigma^2 < \infty$. Тогава за
\[ Z_n := \frac{\sqrt{n}}{\sigma} E_n = \frac{\sqrt{n}}{\sigma} \left( \frac{S_n}{n} - \mu \right) = \frac{S_n - n\mu}{\sigma\sqrt{n}} \]
е в сила сходимост по разпределение
\[ Z_n \xrightarrow[n \to \infty]{d} Z \sim N(0,1), \]
където $Z$ е стандартно нормално разпределена случайна величина (със средно $0$ и дисперсия $1$).
\end{keythm}

Забележете разликата със закона за големите числа: тук изискваме съществуването на втория момент (дисперсията), което моментално влече съществуването на очакването.

\subsection{Коментари и интуиция зад ЦГТ}

Какво ни казва най-грубо този резултат? От алгебрична гледна точка можем да изразим разликата като:
\[ \frac{S_n}{n} - \mu \approx \frac{\sigma}{\sqrt{n}} Z \]
Това означава, че грешката между емпиричното средно $S_n/n$ и истинското средно $\mu$ е от порядъка на $\frac{\sigma}{\sqrt{n}}$, умножено по някаква случайна пертурбация $Z$, концентрирана около нулата (тъй като разпределението на $Z$ е нормално с връх в нулата). Скоростта на сходимост към нулата се предопределя изцяло от детерминистичното количество $\frac{\sigma}{\sqrt{n}}$, докато флуктуациите идват от $Z$.

\textbf{Универсалност:} Защо този резултат е толкова важен? Защото ние никъде не специфицираме какво е разпределението на изходните случайни величини $X_i$. Те могат да бъдат дискретни, непрекъснати, с най-разнообразни плътности или вероятностни маси. Единственото, което зависи, е да познаваме техните две основни характеристики: средно $\mu$ и дисперсия $\sigma^2$. Веднъж фиксирани ли са те, кумулативният ефект при събиране винаги асимптотично клони към нормално разпределение.

\textbf{Сходимост по разпределение:} Какво означава $Z_n \xrightarrow{d} Z$? Това означава сходимост на функциите на разпределение. Тъй като функцията на разпределение на стандартното нормално разпределение $\Phi(z)$ е непрекъсната за всяко $z \in \mathbb{R}$, имаме:
\[ \mathbb{P}(Z_n < b) \xrightarrow[n \to \infty]{} \mathbb{P}(Z < b) = \Phi(b). \]
По същия начин, за всеки интервал $(a, b)$:
\[ \mathbb{P}(a \le Z_n < b) \xrightarrow[n \to \infty]{} \mathbb{P}(a \le Z < b) = \Phi(b) - \Phi(a). \]
Поради непрекъснатостта на $Z$, няма значение дали интервалите са отворени или затворени.

Ако преработим неравенството $a \le \frac{S_n - n\mu}{\sigma\sqrt{n}} < b$, получаваме:
\[ \mathbb{P}\left( n\mu + a\sigma\sqrt{n} \le S_n < n\mu + b\sigma\sqrt{n} \right) \xrightarrow[n \to \infty]{} \Phi(b) - \Phi(a). \]
Ако вземем много широки граници, например $a = -100$ и $b = 100$, интегралът на нормалната плътност в тези граници е на практика единица, следователно вероятността парциалната сума $S_n$ да попадне в този обхват също е приблизително 1.

\subsection{Примери за приложение на теорема~\ref{thm:clt}}

\begin{example}[Игра на дърпане на въже]\label{ex:11-1}
Връщаме се към \crosslecture{ex:tug-of-war}{модела с дърпане на въже от предходната лекция}. Около въжето се подреждат хора отляво и отдясно. Всеки човек има вероятност $1/2$ да застане отдясно и $1/2$ да застане отляво, и всички дърпат с еднаква сила. Нека $X_j$ бъде приносът на $j$-тия човек, като $X_j = 1$ с вероятност $1/2$ (дърпа надясно) и $X_j = -1$ с вероятност $1/2$ (дърпа наляво). Резултантната сила е:
\[ S_n = \Delta_n - \Lambda_n = \sum_{j=1}^n X_j \]
където $\Delta_n$ и $\Lambda_n$ са съответно броят хора отдясно и отляво. Ясно е, че очакването е $\E[X_1] = 0$. Дисперсията е:
\[ \Var[X_1] = \E[X_1^2] = (1)^2 \frac{1}{2} + (-1)^2 \frac{1}{2} = 1. \]
Следователно $\mu=0$ и $\sigma=1$. По ЦГТ имаме:
\[ \frac{S_n}{\sqrt{n}} \approx Z \sim N(0,1). \]
Каква е вероятността дясната страна да победи (т.е. $S_n > 0$)?
\[ \mathbb{P}(S_n > 0) = \mathbb{P}\left(\frac{S_n}{\sqrt{n}} > 0\right) \approx \mathbb{P}(Z > 0) = \frac{1}{2}. \]
Въпреки че съществува вероятност за равенство ($S_n = 0$), тя клони към нула при $n \to \infty$, тъй като разпределението се апроксимира от непрекъснато такова.

\end{example}

\begin{example}[Непрекъсната случайна величина]\label{ex:11-2}
Нека сега $Y_1$ бъде непрекъсната случайна величина с плътност
\[
f_{Y_1}(y) =
\begin{cases}
\tfrac{1}{2} e^{-y}, & y > 0 \\[2pt]
\tfrac{1}{4}, & y \in (-2, 0) \\[2pt]
0, & \text{иначе.}
\end{cases}
\]
Плътността е подчертано несиметрична — експоненциално отмаляване надясно, константа наляво — и въпреки това $\E[Y_1] = 0$.\footnote{Директно пресмятане дава $\E[Y_1] = \tfrac12 - \tfrac12 = 0$ и $\E[Y_1^2] = \tfrac53$, тоест $\Var[Y_1] = \tfrac53$. (В лекцията дисперсията беше спомената предположително като единица.) Стойността ѝ не влияе на извода: нормирането в ЦГТ се извършва с $\sigma$, каквото и да е то.}
Ако образуваме сумата $S_n' = \sum_{j=1}^n Y_j$, поведението на $\frac{S_n'}{\sqrt{n}}$ ще бъде абсолютно същото като това на дискретните $\pm 1$ стъпки от предния пример. То отново клони към $Z \sim N(0,1)$. В граничния случай изобщо няма разлика между бинарно дискретно разпределение и непрекъснато разпределение, когато става въпрос за суми от голям брой независими екземпляри.

\end{example}

\section{Функции на моментите (Moment Generating Functions)}

За да докажем строго (или поне идейно) Централната гранична теорема~\ref{thm:clt}, ни е необходим мощен математически инструмент — функциите на моментите.

\begin{defn}
Нека $X$ е случайна величина. Функцията на моментите на $X$, означена с $M_X(t)$, се дефинира като математическото очакване на експонентата от $tX$:
\[ M_X(t) = \E[e^{tX}] \]
за всички $t$, за които това очакване съществува (т.е. е крайно), обикновено в някакъв интервал $|t| < r_0$ ($r_0 > 0$). За дискретни случайни величини това се пресмята като:
\[ M_X(t) = \sum_{i} e^{t x_i} \mathbb{P}(X = x_i). \]
За непрекъснати случайни величини се пресмята чрез интеграл върху плътността:
\[ M_X(t) = \int_{-\infty}^{\infty} e^{tx} f_X(x) dx. \]
\end{defn}

\begin{example}[Равномерно разпределение в интервала от 0 до 1]\label{ex:11-3}
Случайната величина $X \sim U(0,1)$ има плътност $f_X(x) = 1$ за $x \in [0,1]$ и 0 извън този интервал. Функцията на моментите е:
\[ M_X(t) = \int_0^1 e^{tx} \cdot 1 \, dx = \left. \frac{e^{tx}}{t} \right|_{0}^{1} = \frac{e^t - 1}{t}. \]
Въпреки че функцията изглежда недефинирана за $t=0$, по правилото на Лопитал или развитие в ред на Тейлър, лесно се вижда, че $\lim_{t \to 0} \frac{e^t - 1}{t} = 1$, което е точно $\E[e^0] = 1$. Следователно $M_X(t)$ е добре дефинирана за всяко реално $t$ (т.е. $r_0 = \infty$).

\end{example}

\begin{example}[Експоненциално разпределение]\label{ex:11-4}
Нека $X \sim \Exponential(1)$ с плътност $f_X(x) = e^{-x}$ за $x \ge 0$.
\[ M_X(t) = \int_0^\infty e^{tx} e^{-x} dx = \int_0^\infty e^{-x(1-t)} dx. \]
Този интеграл съществува и е равен на:
\[ M_X(t) = \frac{1}{1-t}, \quad \text{за } t < 1. \]
Ако $t \ge 1$, експонентата не затихва и интегралът е безкрайност. В този случай функцията на моментите съществува сигурно в интервала $(-1, 1)$, тоест за $|t| < 1$.

\end{example}

\subsection{Моменти на случайна величина}
Преди да преминем към свойствата, нека дефинираме самите моменти, откъдето идва името на функцията:
\begin{itemize}
    \item Обикновен момент от ред $k$: $\E[X^k]$. За $k=1$ това е очакването $\mu$.
    \item Абсолютен момент от ред $k$: $\E[|X|^k]$.
    \item Централен момент от ред $k$: $\E[(X - \mu)^k]$. За $k=2$ това е дисперсията $\Var[X] = \sigma^2$. Центрираме величината, изваждайки нейното средно.
    \item Абсолютен централен момент от ред $k$: $\E[|X - \mu|^k]$.
\end{itemize}
В практиката най-често се използват моментите до четвърти ред (очакване, дисперсия, асиметрия/skewness и ексцес/kurtosis).

\subsection{Свойства на функциите на моментите}

\begin{enumerate}
    \item \textbf{Стойност в нулата:} $M_X(0) = \E[e^0] = 1$.
    \item \textbf{Връзка с моментите:} Използвайки развитието на експоненциалната функция в ред на Тейлър, $e^{tX} = \sum_{k=0}^\infty \frac{t^k X^k}{k!}$, и разменяйки очакването и сумата, получаваме:
    \[ M_X(t) = \E \left[ \sum_{k=0}^\infty \frac{t^k X^k}{k!} \right] = \sum_{k=0}^\infty \frac{t^k}{k!} \E[X^k]. \]
    Това представяне кодира моментите на $X$. $k$-тата производна на $M_X(t)$ в точката $t=0$ ни дава точно $k$-тия момент:
    \[ \left. \frac{d^k}{dt^k} M_X(t) \right|_{t=0} = \E[X^k]. \]
    \item \textbf{Независимост:} Ако $X$ и $Y$ са независими случайни величини, тогава:
    \[ M_{X+Y}(t) = \E[e^{t(X+Y)}] = \E[e^{tX} e^{tY}] = \E[e^{tX}] \E[e^{tY}] = M_X(t) M_Y(t). \]
    Това свойство прави функциите на моментите изключително удобни при работа със суми на независими величини.
    \item \textbf{Линейна трансформация:} Ако $Y = aX + b$, тогава:
    \[ M_Y(t) = \E[e^{t(aX+b)}] = e^{bt} \E[e^{atX}] = e^{bt} M_X(at). \]
    \item \textbf{Единственост:} Ако $M_X(t) = M_Y(t)$ за всички $t$ в някаква околност на нулата ($|t| < r_0$), то $X$ и $Y$ имат едно и също разпределение ($X \stackrel{d}{=} Y$).
    \item \textbf{Сходимост (Теорема на Леви-Крамер):} Ако редица от функции на моментите $M_{X_n}(t)$ се схожда поточково към функция на моментите $M_X(t)$ за всички $t \in (-\varepsilon, \varepsilon)$, то съответните случайни величини се схождат по разпределение: $X_n \xrightarrow{d} X$. Това е в основата на доказателството на ЦГТ.
\end{enumerate}

\begin{remark}[Когато функцията на моментите не съществува]\label{rem:charfun}
Съществуват случайни величини с крайни средно и дисперсия, за които функцията на
моментите не съществува. За тях същата работа върши \emph{характеристичната
функция} $\varphi_X(t) = \E[e^{itX}]$, която е дефинирана винаги, тъй като
$|e^{itX}| = 1$. Идеята е сходна, но апаратът, необходим за нея (комплексен
анализ), излиза извън рамките на този курс, затова тук работим само с $M_X(t)$.
\end{remark}

\subsection{Функция на моментите на нормалното разпределение}

Нека $Z \sim N(0,1)$. Плътността е $f_Z(y) = \frac{1}{\sqrt{2\pi}} e^{-y^2/2}$. Изчисляваме $M_Z(s)$:
\[ M_Z(s) = \frac{1}{\sqrt{2\pi}} \int_{-\infty}^{\infty} e^{sy} e^{-\frac{y^2}{2}} dy = \frac{1}{\sqrt{2\pi}} \int_{-\infty}^{\infty} e^{-\frac{1}{2}(y^2 - 2sy)} dy. \]
Допълваме израза в експонентата до точен квадрат:
\[ y^2 - 2sy = y^2 - 2sy + s^2 - s^2 = (y-s)^2 - s^2. \]
Заместваме обратно в интеграла:
\[ M_Z(s) = \frac{1}{\sqrt{2\pi}} \int_{-\infty}^{\infty} e^{-\frac{1}{2}((y-s)^2 - s^2)} dy = e^{\frac{s^2}{2}} \int_{-\infty}^{\infty} \frac{1}{\sqrt{2\pi}} e^{-\frac{(y-s)^2}{2}} dy. \]
Интегралът отдясно не е нищо друго освен лицето под плътността на нормално разпределение със средно $s$ и дисперсия 1. Тъй като това е валидна плътност, интегралът е точно 1. Следователно:
\[ M_Z(s) = e^{\frac{s^2}{2}}. \]

Сега, ако имаме произволна нормална случайна величина $X \sim N(\mu, \sigma^2)$, ние знаем от свойствата на нормалното разпределение (структурна теорема), че $X$ може да се представи като $X = \mu + \sigma Z$. Прилагайки свойството за линейна трансформация, получаваме:
\[ M_X(t) = e^{\mu t} M_Z(\sigma t) = e^{\mu t} e^{\frac{(\sigma t)^2}{2}} = e^{\mu t + \frac{\sigma^2}{2} t^2}. \]

\section{Доказателство на теорема~\ref{thm:clt}}

\textbf{Цел:} Да докажем, че $Z_n = \frac{S_n - n\mu}{\sigma \sqrt{n}} \xrightarrow[n \to \infty]{d} Z \sim N(0,1)$.

Ще преработим израза за $Z_n$, за да го приведем в по-удобен вид за работа с функции на моментите. Разделяме всяко от събираемите:
\[ Z_n = \frac{1}{\sqrt{n}} \frac{S_n - n\mu}{\sigma} = \frac{1}{\sqrt{n}} \sum_{j=1}^n \frac{X_j - \mu}{\sigma} = \frac{1}{\sqrt{n}} \sum_{j=1}^n Y_j \]
където сме дефинирали $Y_j = \frac{X_j - \mu}{\sigma}$.
Очевидно $Y_j$ са независими и еднакво разпределени случайни величини. За тяхното средно и дисперсия имаме:
\[ \E[Y_1] = \frac{\E[X_1] - \mu}{\sigma} = 0, \quad \Var[Y_1] = \frac{\Var[X_1]}{\sigma^2} = 1. \]
Ще допуснем за улеснение, че функцията на моментите $M_{X_1}(t)$ е добре дефинирана за всички $t$. Следователно, функцията на моментите $M_{Y_1}(t)$ също съществува.

Търсим функцията на моментите на $Z_n$:
\[ M_{Z_n}(t) = M_{\frac{1}{\sqrt{n}} \sum_{j=1}^n Y_j}(t) = M_{\sum_{j=1}^n Y_j}\left(\frac{t}{\sqrt{n}}\right). \]
Тъй като $Y_j$ са независими, функцията на моментите на сумата се разпада в произведение от функциите на моментите на отделните събираеми:
\[ M_{Z_n}(t) = \prod_{j=1}^n M_{Y_j}\left(\frac{t}{\sqrt{n}}\right) = \left( M_{Y_1}\left(\frac{t}{\sqrt{n}}\right) \right)^n. \]
Сега трябва да разберем поведението на $M_{Y_1}\left(\frac{t}{\sqrt{n}}\right)$. Развиваме $M_{Y_1}(s)$ в ред на Тейлър около нулата:
\[ M_{Y_1}(s) = 1 + s\E[Y_1] + \frac{s^2}{2}\E[Y_1^2] + o(s^2). \]
Знаем, че $\E[Y_1] = 0$. За втория момент, тъй като очакването е нула, той съвпада с дисперсията, тоест $\E[Y_1^2] = 1$. Следователно:
\[ M_{Y_1}(s) = 1 + \frac{s^2}{2} + o(s^2). \]
Заместваме $s = \frac{t}{\sqrt{n}}$:
\[ M_{Y_1}\left(\frac{t}{\sqrt{n}}\right) = 1 + \frac{t^2}{2n} + o\left(\frac{t^2}{n}\right). \]
Връщайки се обратно към $M_{Z_n}(t)$, получаваме:
\[ M_{Z_n}(t) = \left( 1 + \frac{t^2 / 2}{n} + o\left(\frac{1}{n}\right) \right)^n. \]
За фиксирано $t$, когато $n \to \infty$, това е добре познатата граница за експоненциалната функция $\lim_{n \to \infty} (1 + \frac{x}{n})^n = e^x$. В нашия случай $x = \frac{t^2}{2}$:
\[ M_{Z_n}(t) \xrightarrow[n \to \infty]{} e^{\frac{t^2}{2}}. \]
Но това е точно функцията на моментите $M_Z(t)$ на стандартното нормално разпределение $Z \sim N(0,1)$. По теоремата за сходимост, тъй като функциите на моментите се схождат, то $Z_n \xrightarrow{d} Z$. Доказателството е завършено.

\section{Приложения и оценка на грешката}

\subsection{ЦГТ за Биномно разпределение}
Нека $X_n \sim \Bin(n, p)$ е биномно разпределена случайна величина. Тя може да бъде представена като сума от независими Бернулиеви случайни величини:
\[ X_n = \sum_{j=1}^n Y_j, \quad Y_j \sim \Ber(p). \]
За Бернулиево разпределение имаме очакване $\mu = p$ и дисперсия $\sigma^2 = p(1-p)$. Следователно, прилагайки ЦГТ за тази сума, получаваме:
\[ \frac{X_n - np}{\sqrt{np(1-p)}} \xrightarrow[n \to \infty]{d} Z. \]
Това ни позволява да оценяваме лесно вероятности от вида $\mathbb{P}(X_n \ge k_n)$, където $k_n = np + a$:
\[ \mathbb{P}(X_n \ge k_n) = \mathbb{P}\left( \frac{X_n - np}{\sqrt{np(1-p)}} \ge \frac{k_n - np}{\sqrt{np(1-p)}} \right) \approx \mathbb{P}\left( Z \ge \frac{k_n - np}{\sqrt{np(1-p)}} \right). \]

\subsection{Оценка на грешката при статистически проучвания и Монте Карло}
Продължаваме \crosslecture{ex:monte-carlo}{приложението на закона за големите числа към Монте Карло от предходната лекция}. Ако интервюираме $N$ човека за кого ще гласуват, или пък използваме метод Монте Карло за намиране на лице на фигура чрез хвърляне на случайни точки, ние реално правим серия от експерименти на Бернули. Истинската вероятност (пропорция) за успех е неизвестно $p$.
Искаме да намерим колко експеримента са нужни, така че емпиричното средно $S_N/N$ да се отклонява от $p$ с повече от $\varepsilon$ с вероятност по-малка от $\delta$. Дефинираме грешката $E_N = \frac{S_N}{N} - p$.
\[ \mathbb{P}(|E_N| > \varepsilon) = \mathbb{P}\left( \frac{|S_N - Np|}{N} > \varepsilon \right) = \mathbb{P}\left( \frac{|S_N - Np|}{\sqrt{N p(1-p)}} > \frac{\varepsilon \sqrt{N}}{\sqrt{p(1-p)}} \right). \]
Тъй като $p$ е неизвестно, ние максимизираме дисперсията. Знаем, че $p(1-p) \le 1/4$ за $p \in (0,1)$ (достига се при $p=1/2$). Оттук $\sqrt{p(1-p)} \le 1/2$. Следователно изразът отдясно става:
\[ \frac{\varepsilon \sqrt{N}}{\sqrt{p(1-p)}} \ge \frac{\varepsilon \sqrt{N}}{1/2} = 2\varepsilon\sqrt{N}. \]
Следователно, оценявайки вероятността отгоре, заместваме с по-малко число в условието на нормалното разпределение:
\[ \mathbb{P}(|E_N| > \varepsilon) \le \mathbb{P}(|Z_N| > 2\varepsilon\sqrt{N}) \approx \mathbb{P}(|Z| > 2\varepsilon\sqrt{N}). \]
Това може да се изрази чрез интеграла от плътността на нормалното разпределение:
\[ \mathbb{P}(|Z| > 2\varepsilon\sqrt{N}) = \frac{1}{\sqrt{2\pi}} \int_{|y| > 2\varepsilon\sqrt{N}} e^{-\frac{y^2}{2}} dy. \]
Ако положим $2\varepsilon\sqrt{N} = A$, тоест $\varepsilon = \frac{A}{2\sqrt{N}}$, можем да запишем:
\[ \mathbb{P}\left(|E_N| > \frac{A}{2\sqrt{N}}\right) \le \frac{1}{\sqrt{2\pi}} \int_{|y|>A} e^{-\frac{y^2}{2}} dy \le \frac{2}{\sqrt{2\pi}} \frac{e^{-\frac{A^2}{2}}}{A} \le \delta. \]
Множителят $2$ идва от това, че интегрираме по двете опашки: $\int_A^\infty e^{-y^2/2}\,dy \le e^{-A^2/2}/A$.
По този начин, ако фиксираме нивото на достоверност $\delta$ (напр. $0{,}01$) и желаната грешка $\varepsilon$, можем да решим горното уравнение, за да изчислим необходимия обем извадка $N$. Това е математическата обосновка зад факта, че социологическите проучвания обикновено включват извадка от около 1000 до 2000 души, за да постигнат надеждни резултати.

Важно е да отбележим колко точно е нашето приближение. Знакът $\approx$ между функцията на разпределение на $Z_n$ и на стандартното нормално разпределение може да бъде строго оценен.

\begin{thm}[Неравенство на Бери\,--\,Есен]\label{thm:berry-esseen}
За всяко $x \in \mathbb{R}$
\[ \left| \mathbb{P}(Z_n < x) - \Phi(x) \right| \le \frac{C \cdot \E[|X_1 - \mu|^3]}{\sigma^3 \sqrt{n}}, \]
където $C$ е абсолютна константа (около $0{,}47$).
\end{thm}
 Това показва, че грешката на самата апроксимация по ЦГТ намалява със скорост $1/\sqrt{n}$ и зависи от третия централен момент на случайните величини.

\section{Функция на моментите на Гама разпределение}

Като приложение, демонстриращо гъвкавостта на този инструмент, ще изчислим функцията на моментите за Гама разпределение. Нека $X \sim \Gamma(\alpha, \beta)$. Плътността е дадена от:
\[ f_X(x) = \frac{\beta^\alpha}{\Gamma(\alpha)} x^{\alpha-1} e^{-\beta x}, \quad \text{за } x > 0. \]
Функцията на моментите се задава с интеграла:
\[ M_X(t) = \int_0^\infty e^{tx} \frac{\beta^\alpha}{\Gamma(\alpha)} x^{\alpha-1} e^{-\beta x} dx = \frac{\beta^\alpha}{\Gamma(\alpha)} \int_0^\infty x^{\alpha-1} e^{-x(\beta-t)} dx. \]
За да решим интеграла, полагаме нова променлива $y = x(\beta-t)$, което предполага, че $t < \beta$ (за да е сходящ интегралът). Тогава $x = \frac{y}{\beta-t}$ и $dx = \frac{dy}{\beta-t}$. Замествайки:
\[ M_X(t) = \frac{\beta^\alpha}{\Gamma(\alpha)} \int_0^\infty \left( \frac{y}{\beta-t} \right)^{\alpha-1} e^{-y} \frac{dy}{\beta-t} = \frac{\beta^\alpha}{\Gamma(\alpha)(\beta-t)^\alpha} \int_0^\infty y^{\alpha-1} e^{-y} dy. \]
Забелязваме, че подинтегралният израз заедно с $1/\Gamma(\alpha)$ формира плътността на Гама разпределение с параметри $\Gamma(\alpha, 1)$:
\[ \frac{1}{\Gamma(\alpha)} y^{\alpha-1} e^{-y} \]
Тъй като интегралът от всяка валидна плътност в граници от $0$ до $\infty$ е 1, то:
\[ \int_0^\infty y^{\alpha-1} e^{-y} dy = \Gamma(\alpha). \]
Тогава изразът се опростява до:
\[ M_X(t) = \frac{\beta^\alpha}{(\beta-t)^\alpha}, \quad \text{за } t < \beta. \]

С помощта на функцията на моментите лесно пресмятаме самите моменти. За да намерим очакването, диференцираме един път:
\[ M_X'(t) = \alpha \beta^\alpha (\beta-t)^{-\alpha-1} \]
Заместваме $t=0$:
\[ \E[X] = M_X'(0) = \alpha \beta^\alpha \beta^{-\alpha-1} = \frac{\alpha}{\beta}. \]
За втория момент намираме втората производна:
\[ M_X''(t) = \alpha(\alpha+1) \beta^\alpha (\beta-t)^{-\alpha-2} \]
При $t=0$:
\[ \E[X^2] = M_X''(0) = \alpha(\alpha+1) \beta^\alpha \beta^{-\alpha-2} = \frac{\alpha(\alpha+1)}{\beta^2}. \]
Тази техника спестява тежкото директно интегриране и е мощен похват в теорията на вероятностите.

\section{Задачи}

\begin{enumerate}
    \item За $Z \sim N(0,1)$ проверете чрез допълване до квадрат, смяна на променливата и свойството, че интегралът от плътността е единица, че
    \[
    M_Z(s)=\frac{1}{\sqrt{2\pi}}\int_{-\infty}^{\infty}e^{sy-y^2/2}\,dy=e^{s^2/2}.
    \]

    \item В модела с дърпане на въже от пример~\ref{ex:11-1} променете условието за победа: дясната страна печели само ако резултантната сила е по-голяма от $100$. Помислете как Централната гранична теорема се прилага към новото събитие.

    \item За $X\sim\Gamma(\alpha,\beta)$ използвайте
    \[
    M_X(t)=\frac{\beta^\alpha}{(\beta-t)^\alpha},\qquad t<\beta,
    \]
    и последователно диференциране в $t=0$, за да пресметнете моментите от по-висок ред.
\end{enumerate}


\chapter{Точкови оценки. Максимално правдоподобие и метод на моментите}
\section{Въведение в статистиката и концепцията за извадки}

Започваме нашето занимание със статистиката — наука, която лежи в основата на вероятностите, използвайки техния фундамент за изграждането на апарат, техники и модели, чрез които можем да изучаваме явления около нас, за които имаме непълна или липсваща информация.

Какво означава това? Представете си, че имате някаква генерална съвкупност — например всички хора в дадена държава, всички клетки в организма, молекули или гласоподаватели. Искаме да изучим някакво тяхно свойство, като например средната височина на хората в България или пропорцията на мъжете и жените. Оказва се обаче, че изследването на всеки един обект от тази съвкупност е твърде скъпо и непрактично. Вместо да разпитваме всеки отделен човек (което би било непосилно), ние правим оценка за свойството на базата на \emph{извадка} — подмножество от обектите, което избираме, за да разберем какво е разпределението на изследваната характеристика.

Самият избор на извадка не е тривиална задача. В историята има показателни примери за провали при съставянето ѝ. Един от тях е през 1936 г. за президентските избори в САЩ. Една компания прави опит да прогнозира изхода от изборите въз основа на много голяма извадка, набирайки я чрез телефонни указатели. По онова време обаче телефони са имали предимно по-богатите хора. Поради това извадката е имала систематично изкривяване (bias) към определена социална страта и прогнозата им се проваля. За разлика от тях, фирмата на Джордж Галъп с едва 4000 респонденти успява да прогнозира правилно, защото тяхната извадка е била представителна и подбрана напълно случайно, отразявайки реалната структура на обществото (включително както по-богати, така и по-бедни). В нашите разглеждания ние винаги ще предполагаме, че извадката се прави напълно случайно и е представителна за генералната съвкупност.

\subsection{Математическа постановка}

Имаме някаква случайна величина $X$, описваща интересуващото ни свойство (например височина, доход, глас за определен кандидат). Тя има функция на разпределение $F_X(x) = \mathbb{P}(X < x)$, а в случай на непрекъсната величина — и плътност $f_X(x)$.

Целта ни е, чрез извадка от $n$ на брой наблюдения, да оценим тази функция на разпределение или нейната плътност. Тъй като изборът е случаен, самите наблюдения $X_1, X_2, \dots, X_n$ формират вектор от \emph{независими и еднакво разпределени случайни величини} (н.е.р.с.в.), като всяко $X_j$ има същото разпределение като $X$:
\[
    \vec{X} = (X_1, X_2, \dots, X_n)
\]
Това е концептуалната постановка. В практиката ние провеждаме експеримента или запитването и получаваме конкретни числови стойности $x_1, x_2, \dots, x_n$, които представляват реализацията на случайната величина $\vec{X}$. От тези числа, ние трябва да направим обосновано предположение за функцията на разпределение $F_X(x)$ или за плътността $f_X(x)$.

\section{Точкови оценки}

Ако нямаме никакви предварителни допускания за функцията на разпределение на $X$, задачата е много трудна. В класическата параметрична статистика ние стесняваме задачата, като допускаме, че $X$ принадлежи на определен \emph{клас} от случайни величини, параметризиран с даден параметър $\theta$.

Например:
\begin{itemize}
    \item $X \sim \Ber(p)$ — Бернулиево разпределение, където $\theta = p$ (вероятността за успех, напр. гласуване за даден кандидат).
    \item $X \sim \text{Exp}(\lambda)$ — Експоненциално разпределение за време на живот на електронни компоненти, където $\theta = \lambda$.
    \item $X \sim N(\mu, \sigma^2)$ — Нормално разпределение, където параметърът е двумерен вектор $\theta = (\mu, \sigma^2)$.
\end{itemize}

По този начин, вместо да търсим произволна функция, нашата цел се свежда до това да намерим оценка $\hat{\theta}$ за параметъра $\theta$ въз основа на наблюденията. Това води до концепцията за \emph{точкова оценка} (Point Estimator).

Точковата оценка $\hat{\theta} = \hat{\theta}(\vec{X})$ е функция на случайния вектор от наблюденията, която ни служи като приближение за истинската, но неизвестна стойност на параметъра $\theta$. Тъй като всяка функция може формално да се разглежда като оценка, имаме нужда от методи за конструиране на оценки, които притежават оптимални свойства. Два от най-популярните метода са Методът на максималното правдоподобие (МПО) и Методът на моментите (ММО).

\section{Метод на максималното правдоподобие (МПО)}

Методът на максималното правдоподобие, или Maximum Likelihood Estimation (MLE), почива на една много интуитивна идея: от всички възможни стойности за параметъра $\theta$ ние трябва да изберем тази, при която вероятността (или съвместната плътност) да наблюдаваме данните, които реално сме получили, е възможно най-голяма.

\subsection{Дефиниция на функция на правдоподобие}
Нека $X$ е случайна величина, чиято функция на разпределение (или плътност) $f_X(x; \theta)$ е от клас, параметризиран с $\theta$. Имаме независимата извадка $\vec{X} = (X_1, X_2, \dots, X_n)$.

\begin{keydefn}[функция на правдоподобие]\label{def:likelihood}
\emph{Функция на правдоподобие} $L(\vec{X}, \theta)$ наричаме съвместната плътност (или вероятност в дискретния случай) на вектора $\vec{X}$. Тъй като наблюденията са независими, съвместната плътност се разлага в произведение от индивидуалните плътности:
\[
    L(\vec{X}, \theta) = \prod_{j=1}^n f_X(X_j; \theta) .
\]
\end{keydefn}

За конкретна реализация на извадката $\vec{x} = (x_1, \dots, x_n)$, функцията на правдоподобие от дефиниция~\ref{def:likelihood} е:
\[
    L(\vec{x}, \theta) = \prod_{j=1}^n f_X(x_j; \theta)
\]

\textbf{Максимално правдоподобна оценка} (МПО) $\hat{\theta}$ е тази стойност на $\theta$, която максимизира функцията на правдоподобие:
\[
    L(\vec{X}, \hat{\theta}) = \sup_{\theta} L(\vec{X}, \theta)
\]

За удобство при пресмятанията, понеже логаритъмът е строго растяща функция, максимизирането на $L(\vec{X}, \theta)$ е еквивалентно на максимизирането на логаритъма на функцията на правдоподобие:
\[
    \ln L(\vec{X}, \theta) = \sum_{j=1}^n \ln f_X(X_j; \theta)
\]
Ако $f_X(x; \theta)$ е диференцируема по $\theta$, можем да намерим екстремума, решавайки системата уравнения:
\[
    \frac{\partial \ln L(\vec{X}, \theta)}{\partial \theta} = 0
\]

\begin{example}[Равномерно разпределение $U(0, \theta)$]\label{ex:12-1}

Допускаме, че $X$ е равномерно разпределена в интервала $(0, \theta)$, т.е. $X \sim U(0, \theta)$. Параметърът $\theta > 0$. Плътността е дадена чрез индикаторната функция:
\[
    f_X(x; \theta) = \frac{1}{\theta} \ind_{(0, \theta)}(x)
\]
Функцията на правдоподобие за $n$ наблюдения е:
\[
    L(\vec{X}, \theta) = \prod_{j=1}^{n} \frac{1}{\theta} \ind_{(0, \theta)}(X_j) = \frac{1}{\theta^n} \ind_{(0, \theta)}(X^*) = \begin{cases} \frac{1}{\theta^n} & \text{ако } X^* \le \theta \\ 0 & \text{ако } X^* > \theta \end{cases}
\]
където $X^* = \max_{j \le n}(X_j)$ е максималната стойност сред наблюденията. Забележете, че произведението на индикаторните функции $\prod_{j=1}^n \ind_{(0, \theta)}(X_j)$ е равно на 1 тогава и само тогава, когато всички $X_j < \theta$, което е еквивалентно на това техният максимум $X^*$ да е по-малък от $\theta$.

Търсим $\hat{\theta}$, което максимизира $L(\vec{X}, \theta)$:
\[
    L(\vec{X}, \hat{\theta}) = \sup_{\theta > 0} L(\vec{X}, \theta)
\]
Тъй като функцията $\frac{1}{\theta^n}$ е строго намаляваща по $\theta$, нейната максимална стойност се достига в най-малката възможна точка за $\theta$, при която функцията не е нула. Според индикатора, трябва $\theta \ge X^*$. Следователно най-голямата стойност се постига, когато:
\[
    \hat{\theta}_1 = X^* = \max_{j \le n}(X_j)
\]
Ако например сме получили наблюденията 1, 2, 3, 4 и 5, максимално правдоподобната оценка за $\theta$ би била просто най-голямата стойност — 5.

\end{example}

\begin{example}[Нормално разпределение $N(\mu, \sigma^2)$]\label{ex:12-2}

Нека $X \sim N(\mu, \sigma^2)$. Имаме два параметъра за оценяване: $\theta = (\mu, \sigma^2)$.
Плътността на нормалното разпределение за едно наблюдение е:
\[
    f_X(x; \mu, \sigma^2) = \frac{1}{\sqrt{2\pi\sigma^2}} e^{-\frac{(x - \mu)^2}{2\sigma^2}}
\]
Съставяме функцията на правдоподобие за извадката $\vec{X}$:
\[
    L(\vec{X}; \theta) = \left(\frac{1}{\sqrt{2\pi}}\right)^n \sigma^{-n} e^{-\sum_{j=1}^n \frac{(X_j - \mu)^2}{2\sigma^2}}
\]
Логаритмуваме я, за да превърнем произведението в сума, което значително улеснява диференцирането:
\[
    \ln L(\vec{X}, \theta) = n \ln \frac{1}{\sqrt{2\pi}} - \frac{n}{2} \ln \sigma^2 - \sum_{j=1}^n \frac{(X_j - \mu)^2}{2\sigma^2}
\]
Намираме частната производна спрямо $\mu$ и я приравняваме на нула:
\[
    \frac{\partial \ln L}{\partial \mu} = \frac{1}{\sigma^2} \sum_{j=1}^n (X_j - \mu) = 0
\]
\[
    \sum_{j=1}^n X_j - n\mu = 0 \implies \hat{\mu} = \frac{1}{n} \sum_{j=1}^n X_j = \overline{X}_n
\]
Оценката за $\mu$ е средното аритметично на наблюденията. Забележете, че тя не зависи от това дали познаваме $\sigma^2$ или не.

Сега намираме частната производна спрямо $\sigma^2$ и я приравняваме на нула:
\[
    \frac{\partial \ln L}{\partial \sigma^2} = -\frac{n}{2\sigma^2} + \frac{1}{2\sigma^4} \sum_{j=1}^n (X_j - \mu)^2 = 0
\]
Ако параметърът $\mu$ ни е предварително известен, решението директно ни дава:
\[
    \hat{\sigma}^2 = \frac{1}{n} \sum_{j=1}^n (X_j - \mu)^2
\]
Ако $\mu$ е неизвестен, заместваме неговата оценка $\hat{\mu} = \overline{X}_n$ в уравнението, откъдето получаваме:
\[
    \hat{\sigma}^2 = \frac{1}{n} \sum_{j=1}^n (X_j - \overline{X}_n)^2
\]


\end{example}

\section{Метод на моментите (ММО)}

Методът на моментите почива на Закона за големите числа, според който емпиричните (извадкови) моменти се схождат почти сигурно към теоретичните моменти на разпределението при увеличаване на обема на извадката.

Нека въведем емпиричния момент от ред $j$:
\[
    \bar{X}_n^{(j)} = \frac{1}{n} \sum_{k=1}^n X_k^j, \quad j \ge 1
\]
Това е средното аритметично от наблюденията, повдигнати на $j$-та степен. (При $j=1$ получаваме обикновеното средно аритметично $\bar{X}_n$).

\begin{defn}[ММО]
Нека $X$ е случайна величина с функция на разпределение $F_X(x; \theta)$, където параметърът е $s$-мерен вектор $\theta = (\theta_1, \dots, \theta_s)$. Оценката по метода на моментите $\hat{\theta} = (\hat{\theta}_1, \dots, \hat{\theta}_s)$ намираме, решавайки следната система от $s$ уравнения:
\[
    \mu^{(j)}(\theta) = \bar{X}_n^{(j)}, \quad 1 \le j \le s
\]
където $\mu^{(j)}(\theta) = \E[X^j]$ са теоретичните моменти на случайната величина като функция на параметъра $\theta$.
\end{defn}

Идеята е да приравним първите $s$ на брой теоретични момента към съответните извадкови моменти и да решим получената система спрямо неизвестните параметри. Тъй като $\bar{X}_n^{(j)} \to \E[X^j]$ при $n \to \infty$, методът дава оценки, които при определени условия на обратимост клонят към истинските стойности на параметрите.

\begin{example}[Равномерно разпределение $U(0, \theta)$ (ММО)]\label{ex:12-3}

Разглеждаме отново случая $X \sim U(0, \theta)$. Тъй като имаме само един параметър, се нуждаем от едно уравнение — за първия момент.
Теоретичното очакване на $X$ е:
\[
    \E[X] = \frac{\theta}{2}
\]
Емпиричният първи момент е средното аритметично $\bar{X}_n$. Приравняваме ги:
\[
    \frac{\theta}{2} = \bar{X}_n
\]
Оттук директно изразяваме оценката по ММО за $\theta$:
\[
    \hat{\theta}_2 = 2 \bar{X}_n
\]
Виждаме, че за равномерното разпределение двата метода (МПО и ММО) дават коренно различни оценки ($\hat{\theta}_1 = \max(X_j)$ срещу $\hat{\theta}_2 = 2\bar{X}_n$).

\end{example}

\begin{example}[Нормално разпределение $N(\mu, \sigma^2)$ (ММО)]\label{ex:12-4}

Нека $X \sim N(\mu, \sigma^2)$. Имаме два параметъра ($\mu$ и $\sigma^2$), затова ни трябват първите два момента.
Теоретичните моменти са:
\[
    \E[X] = \mu
\]
\[
    \E[X^2] = \Var(X) + (\E[X])^2 = \sigma^2 + \mu^2
\]
Съставяме системата от уравнения, приравнявайки ги към емпиричните моменти:
\begin{align*}
    \mu &= \bar{X}_n^{(1)} \\
    \sigma^2 + \mu^2 &= \bar{X}_n^{(2)}
\end{align*}
От първото уравнение веднага получаваме:
\[
    \hat{\mu} = \bar{X}_n^{(1)} = \frac{1}{n} \sum_{j=1}^n X_j
\]
Замествайки $\hat{\mu}$ във второто уравнение, получаваме оценката за дисперсията:
\[
    \hat{\sigma}^2 = \bar{X}_n^{(2)} - \left(\bar{X}_n^{(1)}\right)^2 = \frac{1}{n} \sum_{j=1}^n X_j^2 - \left(\frac{1}{n} \sum_{j=1}^n X_j\right)^2 = \frac{1}{n} \sum_{j=1}^n \left(X_j - \bar{X}_n^{(1)}\right)^2
\]
В случая на нормалното разпределение, оценките по метода на моментите (ММО) и по метода на максималното правдоподобие (МПО) съвпадат напълно.

\end{example}

\section{Свойства на точковите оценки}

За да определим коя от получените оценки е по-добра, изследваме техните статистически свойства, като \emph{неизместеност} (unbiasedness) и \emph{състоятелност} (consistency).

\begin{defn}[неизместена оценка]\label{def:unbiased}
Една оценка $\hat{\theta}$ се нарича \emph{неизместена}, ако нейното математическо очакване съвпада с истинската стойност на параметъра:
\[ \E[\hat{\theta}] = \theta . \]
\end{defn}

\begin{supp}[оценка с минимална дисперсия; ефективна оценка]\label{supp:efficient}
Неизместената оценка $\hat{\theta}$ се нарича \emph{оценка с минимална дисперсия},
ако за всяка друга неизместена оценка $\tilde{\theta}$ е изпълнено
$\Var \hat{\theta} \le \Var \tilde{\theta}$. Неизместена оценка с минимална
дисперсия се нарича \emph{ефективна}. Ако за даден параметър съществува ефективна
оценка, тя е единствена.
\end{supp}

\subsection{Свойства на оценките при $U(0, \theta)$}

Сравняваме двете получени оценки: $\hat{\theta}_1 = X^*$ (от МПО) и $\hat{\theta}_2 = 2 \bar{X}_n$ (от ММО).

За оценката по метода на моментите $\hat{\theta}_2$:
\[
    \E[\hat{\theta}_2] = \E[2 \bar{X}_n] = \frac{2}{n} \sum_{j=1}^n \E[X_j] = \frac{2}{n} \cdot n \cdot \frac{\theta}{2} = \theta
\]
Тъй като $\E[\hat{\theta}_2] = \theta$, оценката $\hat{\theta}_2$ е \textbf{неизместена} по смисъла на дефиниция~\ref{def:unbiased}. Освен това, благодарение на Закона за големите числа, $\bar{X}_n \xrightarrow{\text{п.с.}} \frac{\theta}{2}$, следователно $\hat{\theta}_2$ е и \textbf{силно състоятелна} оценка.

За оценката по метода на максималното правдоподобие $\hat{\theta}_1 = X^*$:
За да намерим очакването на $X^*$, първо намираме нейната функция на разпределение $F_{X^*}(x)$ за $x \in (0, \theta)$:
\[
    F_{X^*}(x) = \mathbb{P}(X^* < x) = \mathbb{P}\left(\bigcap_{j=1}^n \{X_j < x\}\right)
\]
Тъй като наблюденията са независими и еднакво разпределени, вероятността от сечението се разпада в произведение от вероятностите:
\[
    F_{X^*}(x) \stackrel{\text{незав.}}{=} \prod_{j=1}^n \mathbb{P}(X_j < x) = \left( \mathbb{P}(X_1 < x) \right)^n = \left(\frac{x}{\theta}\right)^n, \quad x \in (0, \theta)
\]
Плътността се получава чрез диференциране:
\[
    f_{X^*}(x) = n \frac{x^{n-1}}{\theta^n}, \quad x \in (0, \theta)
\]
Сега пресмятаме математическото очакване на $X^*$:
\[
    \E[X^*] = \int_0^\theta x \cdot n \frac{x^{n-1}}{\theta^n} dx = \frac{n}{\theta^n} \left[ \frac{x^{n+1}}{n+1} \right]_0^\theta = \frac{n}{n+1} \theta
\]
Тъй като $\E[\hat{\theta}_1] = \frac{n}{n+1} \theta \neq \theta$, оценката по метода на максималното правдоподобие не изпълнява дефиниция~\ref{def:unbiased} и е \textbf{изместена}. Това е интуитивно логично — максималната стойност в извадката винаги ще бъде малко по-малка от същинската горна граница $\theta$. С увеличаването на извадката обаче, този максимум се доближава все повече до истинското $\theta$ (оценката е асимптотично неизместена). Ако желаем строго неизместена оценка, базирана на $X^*$, можем да я конструираме като:
\[
    \tilde{\theta}_1 = \frac{n+1}{n} X^*
\]
Въпреки че $X^*$ е изместена, тя все пак е \textbf{състоятелна}. Това лесно може да се докаже, като изследваме вероятността $X^*$ да остане отдалечена от $\theta$ на разстояние $\epsilon$:
\[
    \mathbb{P}(X^* < \theta - \epsilon) = \left(\frac{\theta - \epsilon}{\theta}\right)^n
\]
При $n \to \infty$, тъй като дробта в скобите е строго по-малка от 1, този израз клони към 0, което показва, че $X^*$ се схожда по вероятност към $\theta$.

\subsection{Изместеност на дисперсията при Нормално разпределение}

Нека обърнем внимание и на получената оценка за дисперсията при $N(\mu, \sigma^2)$:
\[
    \hat{\sigma}^2 = \frac{1}{n} \sum_{j=1}^n \left(X_j - \bar{X}_n^{(1)}\right)^2
\]
Може да се докаже, че $\E[\hat{\sigma}^2] = \frac{n-1}{n} \sigma^2 \neq \sigma^2$. Тоест, получената по методите МПО и ММО оценка $\hat{\sigma}^2$ е изместена (тя систематично подценява истинската дисперсия, понеже използва емпиричното средно вместо истинското). За да се коригира това изкривяване, в практиката се използва \emph{коригираната извадкова дисперсия}:
\[
    S^2 = \frac{n}{n-1} \hat{\sigma}^2 = \frac{1}{n-1} \sum_{j=1}^{n} (X_j - \bar{X}_n^{(1)})^2
\]
За нея вече е в сила:
\[
    \E[S^2] = \frac{n}{n-1} \E[\hat{\sigma}^2] = \frac{n}{n-1} \frac{n-1}{n} \sigma^2 = \sigma^2
\]
което я прави неизместена оценка за параметъра $\sigma^2$.

\section{Задачи}

\begin{enumerate}
    \item За нормалния модел от пример~\ref{ex:12-2} пресметнете частните производни на $\ln L(\vec{X};\mu,\sigma^2)$ по $\mu$ и по $\sigma^2$, приравнете ги на нула и решете получената система.

    \item Нека $X_1,\dots,X_n$ са независими и равномерно разпределени в $(0,\theta)$ и $X^*=\max_{j\le n}X_j$. Проверете състоятелността на оценката $\hat\theta_1=X^*$, като за $\varepsilon>0$ пресметнете
    \[
    \mathbb{P}(X^*<\theta-\varepsilon)
    \]
    и изследвате границата при $n\to\infty$.
\end{enumerate}


\chapter{Доверителни интервали и проверка на хипотези}
\section{Доверителни интервали}

\subsection{Постановка на задачата}
Нека имаме някаква случайна величина $X$ с функция на разпределение $F_X(x, \theta)$, която е от някакъв клас разпределения, параметризиран чрез параметъра $\theta$, където $\theta \in \mathbb{R}$. В тази лекция ще разглеждаме само едномерни модели, тоест параметърът, от който зависи разпределението, е едномерно число.

Разполагаме с извадка $\vec{X} = (X_1, \dots, X_n)$ от наблюдения. На базата на данните от тази извадка ние можем да конструираме точкова оценка $\hat{\theta} = \hat{\theta}(\vec{X})$, която е функция на наблюденията. За точковата оценка $\hat{\theta}$ вече знаем от предишни лекции какви свойства може да притежава: състоятелност, неизместеност, как е добита (по метода на максималното правдоподобие или по метода на моментите) и в какъв смисъл е оптимална.

Доверителният интервал се опитва да реши един друг проблем. Ако имаме някаква извадка, чрез точковата оценка ще получим някакво конкретно число, например $1{,}2546$, може би и с още знаци след десетичната запетая. Възниква въпросът: колко близо е това число до истинското $\theta$? Искаме да имаме някаква хубава гаранция, че с голяма вероятност това число е близо до истинския параметър, който стои зад модела.

Следователно, целта ни е въз основа на извадката $\vec{X}$ да дефинираме един интервал с две граници — лява граница $I_1 = I_1(\vec{X})$ и дясна граница $I_2 = I_2(\vec{X})$, така че вероятността истинският параметър $\theta$ да попадне между тях да бъде поне $\gamma$.

\begin{keydefn}[доверителен интервал и ниво на доверие]\label{def:ci}
Интервалът $(I_1, I_2)$, за който
\[
\mathbb{P}(I_1 < \theta < I_2) \ge \gamma ,
\]
се нарича \emph{доверителен интервал} за $\theta$ с \emph{ниво на доверие} $\gamma$.
\end{keydefn}

\noindent
Навсякъде по-нататък централните статистики са с непрекъснати разпределения, така че
$\mathbb{P}(T = q) = 0$ и няма значение дали неравенствата са строги или не. Затова в
пресмятанията ще ги пишем както е по-удобно.

Често използвани стойности за $\gamma$ са $0{,}95$, $0{,}99$ или $0{,}90$, в зависимост от конкретния проблем, който решаваме.

Много е важно да се отбележи къде точно се крие стохастичността (случайността) в това неравенство. Истинското $\theta$ не е случайно — то е някакъв фиксиран, обективно съществуващ параметър на модела, който ние искаме да открием. Случайни са границите $I_1$ и $I_2$, които се определят от извадката. Концептуално те са случайни величини, а при конкретна реализация на извадката (когато измерим данните), те ще бъдат просто две числа.

\subsection{Централна статистика}
Как можем да намерим такива доверителни интервали? Ще ги дефинираме строго посредством така наречената \emph{централна статистика}.

\begin{defn}[централна статистика]\label{def:pivot}
Нека $T = T(\vec{X}, \theta)$ е функция, която зависи както от данните $\vec{X}$, така и от неизвестния параметър $\theta$. Казваме, че $T$ е \emph{централна статистика}, ако са изпълнени следните две свойства:
\begin{itemize}
    \item[А)] Като функция на $\theta$, $T$ е монотонна (разглеждаме само едномерни модели, където $\theta \in \mathbb{R}$).
    \item[Б)] Функцията на разпределение на $T$, която ще означим с $H(x) = \mathbb{P}(T < x)$, не зависи от параметъра $\theta$.
\end{itemize}
\end{defn}

Тоест, въпреки че самата функция $T$ зависи от две променливи (вектора от наблюдения $\vec{X}$ и параметъра $\theta$), разпределението ѝ не зависи от $\theta$.

\subsection{Конструиране на интервала}
Как централната статистика от дефиниция~\ref{def:pivot} ни позволява да конструираме доверителен интервал с ниво на доверие $\gamma$ по дефиниция~\ref{def:ci}?

Нека приемем за момента, че функцията на разпределение $H(x)$ има плътност $h(x) = H'(x)$, за да можем да си представим нещата нагледно. Целта ни е да подберем две квантилни стойности $q_1$ и $q_2$, такива че вероятностната маса между тях да е точно $\gamma$:
\[
H(q_2) - H(q_1) = \gamma
\]
Тъй като търсим само разликата да е $\gamma$, имаме два свободни параметъра $q_1$ и $q_2$ и много голямо разнообразие от начини, по които да ги изберем. Можем да изместваме $q_1$ и $q_2$ наляво или надясно, стига площта под графиката на плътността между тях да остава $\gamma$.

Най-често $q_1$ и $q_2$ се избират така, че грешката да е симетрична — тоест остатъкът от вероятността $1 - \gamma$ да се раздели поравно в двата края (в двете „опашки“ на разпределението). Така в лявата опашка оставяме маса $\frac{1-\gamma}{2}$, а в дясната опашка оставяме също $\frac{1-\gamma}{2}$:
\[
H(q_1) = \frac{1-\gamma}{2}
\]
\[
H(q_2) = 1 - \frac{1-\gamma}{2} = \frac{1+\gamma}{2}
\]
\begin{lecfigure}[Симетричен избор на квантилите. Защрихованата площ между $q_1$ и $q_2$
е нивото на доверие $\gamma$; в двете опашки остава по $\frac{1-\gamma}{2}$. При
симетрично около нулата разпределение $q_1 = -q$ и $q_2 = q$.\label{fig:quantiles}]
\begin{tikzpicture}[x=1.15cm, y=1.75cm]
  \def\q{1.75}
  % tails
  \fill[exrule!25] (-3.2,0) -- plot[domain=-3.2:-\q, samples=60] (\x,{exp(-\x*\x/2)})
      -- (-\q,0) -- cycle;
  \fill[exrule!25] (\q,0) -- plot[domain=\q:3.2, samples=60] (\x,{exp(-\x*\x/2)})
      -- (3.2,0) -- cycle;
  % central mass
  \fill[thmframe!18] (-\q,0) -- plot[domain=-\q:\q, samples=100] (\x,{exp(-\x*\x/2)})
      -- (\q,0) -- cycle;
  % density
  \draw[cdfstep] plot[domain=-3.2:3.2, samples=140, smooth] (\x,{exp(-\x*\x/2)});
  \draw[axisline] (-3.5,0) -- (3.6,0) node[right, font=\small] {$x$};
  % markers
  \draw[guide] (-\q,0) -- (-\q,{exp(-\q*\q/2)});
  \draw[guide] (\q,0)  -- (\q,{exp(-\q*\q/2)});
  \node[below, font=\small] at (-\q,-0.03) {$q_1$};
  \node[below, font=\small] at (\q,-0.03) {$q_2$};
  \node[font=\small] at (0,0.42) {$\gamma$};
  \node[font=\small] at (-2.42,0.1) {$\frac{1-\gamma}{2}$};
  \node[font=\small] at (2.42,0.1) {$\frac{1-\gamma}{2}$};
\end{tikzpicture}
\end{lecfigure}

Това означава, че ако бъркаме, ще бъркаме еднакво и от ляво, и от дясно. Разбира се, има ситуации (в зависимост от модела), при които може да сложим $q_1 = -\infty$ и да оставим цялата маса $1-\gamma$ в дясната опашка. Ние обаче ще се фокусираме върху симетричния случай.

Досега говорихме за $q_1$ и $q_2$ като за „точки, между които остава маса $\gamma$“. Нека
дадем на това име.

\begin{defn}[Квантил]\label{def:quantile}
Нека $F$ е функция на разпределение и $\gamma \in (0,1)$. \emph{$\gamma$-квантилът}
$q_\gamma$ е числото, за което
\[
  \mathbb{P}(X \le q_\gamma) = \gamma, \qquad X \sim F,
\]
тоест точката, вляво от която остава вероятностна маса $\gamma$. При строго растяща
непрекъсната $F$ такова число съществува и е единствено: $q_\gamma = F^{-1}(\gamma)$.
\end{defn}

\noindent
Записът $q_\gamma$ винаги се отнася към конкретно разпределение, което се уточнява от
контекста: $q_{0{,}975} = 1{,}96$ за $N(0,1)$, но $q_{0{,}975} = 2{,}776$ за $t(4)$.
Числените стойности са дадени в
\crosslecture{app:tables}{приложението със статистическите таблици}.

Ако разпределението на централната статистика е симетрично около нулата (както ще видим при нормалното разпределение), тогава можем да изберем $q_1 = -q$ и $q_2 = q$, където $q = q_{\frac{1+\gamma}{2}}$ е съответният квантил. Така двете граници се отличават само по знак.

След като сме избрали $q_1$ и $q_2$, имаме:
\[
\mathbb{P}(q_1 < T(\vec{X}, \theta) < q_2) = \gamma
\]
Тъй като $T$ е монотонна функция по $\theta$ (нека допуснем, че е монотонно намаляваща), можем да обърнем неравенствата спрямо $\theta$ и да получим:
\[
I_1(\vec{X}) < \theta < I_2(\vec{X})
\]
където $I_1$ и $I_2$ са търсените граници на доверителния интервал.

\section{Примери за доверителни интервали}

\subsection{Известна дисперсия}\label{sec:ci-known-var}
Нека имаме извадка от $n$ независими и еднакво разпределени наблюдения $\vec{X} = (X_1, \dots, X_n)$ от нормално разпределение:
\[
X \sim N(\mu, \sigma^2)
\]
В този първи пример предполагаме, че дисперсията $\sigma^2$ е \textbf{известна}. Неизвестният параметър, за който искаме да конструираме доверителен интервал с ниво на доверие $\gamma = 0{,}95$, е математическото очакване $\mu$.

Знаем, че точковата оценка за $\mu$, която е оптимална (както по метода на максималното правдоподобие, така и по метода на моментите), е средното аритметично на извадката:
\[
\overline{X}_n = \frac{1}{n} \sum_{j=1}^n X_j
\]
Тъй като сума от нормални случайни величини е отново нормална случайна величина, разпределението на средното аритметично е:
\[
\overline{X}_n \sim N\left(\mu, \frac{\sigma^2}{n}\right)
\]
За да конструираме централна статистика, центрираме $\overline{X}_n$ като извадим средното $\mu$ и нормираме, като разделим на стандартното отклонение $\sigma / \sqrt{n}$:
\[
T = \frac{\overline{X}_n - \mu}{\sigma / \sqrt{n}} \sim \underline{\underline{N(0, 1)}}
\]
Статистиката $T$ е изключително удобна. Тя е монотонно намаляваща по $\mu$ (това е линейна функция на $\mu$ с отрицателен коефициент) и разпределението ѝ е стандартно нормално $N(0, 1)$, което изобщо не зависи от $\mu$. Следователно $T$ е централна статистика, а нейната функция на разпределение $H(x) = \Phi(x)$ е добре позната.

Тъй като $N(0,1)$ е симетрично разпределение, можем да вземем квантили $q_1 = -q$ и $q_2 = q$. За ниво на доверие $\gamma = 0{,}95$, търсим $q$ такова, че:
\[
\gamma = 0{,}95 = \mathbb{P}(-q < T < q) = \mathbb{P}\left(-q < \frac{\overline{X}_n - \mu}{\sigma / \sqrt{n}} < q\right)
\]
Това означава, че в двете опашки извън интервала $[-q, q]$ трябва да остане общо вероятност $0{,}05$, или по $0{,}025$ във всяка. От таблиците за стандартното нормално разпределение знаем, че квантилът, съответстващ на площ $0{,}975$ вляво от него, е точно $q = 1{,}96$.

Решаваме неравенствата вътре във вероятността спрямо $\mu$:
\[
-q \frac{\sigma}{\sqrt{n}} < \overline{X}_n - \mu < q \frac{\sigma}{\sqrt{n}}
\]
\[
\mathbb{P}\left( \underbrace{\overline{X}_n - \frac{q\sigma}{\sqrt{n}}}_{T_1} < \mu < \underbrace{\overline{X}_n + \frac{q\sigma}{\sqrt{n}}}_{T_2} \right) = \gamma
\]
Получихме явния вид на левия и десния край на доверителния интервал:
\[
T_1 = \overline{X}_n - \frac{q\sigma}{\sqrt{n}}
\]
\[
T_2 = \overline{X}_n + \frac{q\sigma}{\sqrt{n}}
\]
Тези граници зависят от извадката единствено чрез нейното средно аритметично $\overline{X}_n$.

Дължината на този доверителен интервал е:
\[
T_2 - T_1 = \frac{2q\sigma}{\sqrt{n}}
\]
Тук ясно се вижда едно съвсем естествено и желателно свойство: когато размерът на извадката $n$ нараства (когато имаме повече наблюдения), интервалът се свива. Оценката ни става все по-точна и по-близка до истинското $\mu$.

\subsection{Доверителен интервал за дисперсията при известно $\mu$}

Ако $X \sim N(\mu, \sigma^2)$ и \emph{средното $\mu$ е известно}, а търсим интервал за
$\sigma^2$, за централна статистика служи
\[
V_n = \frac{n\hat{\sigma}^2_{\mu}}{\sigma^2}
    = \sum_{j=1}^{n} \left( \frac{X_j - \mu}{\sigma} \right)^{\!2}
    \sim \chi^2(n),
\]
където $\hat{\sigma}^2_{\mu} = \frac1n \sum_{j=1}^n (X_j - \mu)^2$. Всяко събираемо е квадрат
на стандартно нормална величина, а събираемите са независими, откъдето и разпределението
$\chi^2(n)$. Разпределението не зависи от $\sigma^2$, а $V_n$ е монотонна по $\sigma^2$ —
следователно $V_n$ е централна статистика по дефиниция~\ref{def:pivot} и интервалът се
получава по същата схема като по-горе.\footnote{Този раздел е възстановен от записа на
дъската (кадър в 61-ата минута); аудиозаписът на лекцията прекъсва по-рано.}

\subsection{Неизвестна дисперсия. Разпределение на Стюдънт}
Ситуацията става малко по-сложна, ако дисперсията $\sigma^2$ е \textbf{неизвестна}. Имаме един параметър $\mu$, който ни интересува и за който искаме да конструираме доверителен интервал (например с $\gamma = 0{,}95$), но имаме и втори неизвестен параметър $\sigma^2$.

Точковите оценки за двата параметъра са:
\[
\hat{\mu} = \overline{X}_n
\]
\[
\hat{\sigma}^2 = \frac{\sum_{j=1}^n (X_j - \overline{X}_n)^2}{n} = \frac{n-1}{n} S^2
\]
където $S^2 = \frac{1}{n-1} \sum_{j=1}^n (X_j - \overline{X}_n)^2$ е неизместената точкова оценка за дисперсията.

Как можем да конструираме централна статистика за $\mu$ на базата на $S^2$, без да знаем истинското $\sigma^2$? Ключът към решаването на тази задача се крие в следното много интересно и нетривиално твърдение:
\begin{prop}
Нека $X_1, \dots, X_n$ са независими и \textbf{нормално} разпределени, $X_j \sim N(\mu, \sigma^2)$.
Тогава средното аритметично $\hat{\mu} = \overline{X}_n$ е \textbf{независимо} от извадковата дисперсия $S^2$:
\[
\hat{\mu} \perp\!\!\!\perp S^2
\]
и освен това:
\[
\frac{(n-1)S^2}{\sigma^2} \sim \chi^2(n-1)
\]
Това е дълбок факт, защото $\overline{X}_n$ участва директно във формулата за $S^2$, и въпреки това те са независими случайни величини.
\end{prop}

Нормалността тук е съществена, а не удобно допълнително предположение: може да се докаже, че
$\overline{X}_n$ и $S^2$ са независими \emph{само} при нормално разпределение. Затова цялата
конструкция по-долу — и разпределението на Стюдънт, което следва от нея — важи за нормални
извадки, а за други разпределения е в най-добрия случай приближение.

\emph{Интуитивно обяснение на разпределението на $S^2$:}
Нека разгледаме стандартизираните случайни величини $Z_j = \frac{X_j - \mu}{\sigma} \sim N(0,1)$. Сумата от техните квадрати дава разпределение хи-квадрат ($\chi^2$) с $n$ степени на свобода:
\[
\sum_{j=1}^n Z_j^2 = \sum_{j=1}^n \left( \frac{X_j - \mu}{\sigma} \right)^2 =: U \sim \chi^2(n)
\]
След чисто алгебрични преобразувания, прибавяйки и изваждайки $\overline{X}_n$ в числителя, изразът $U$ може да бъде разбит на две части:
\[
U = (n-1)\frac{\sum_{j=1}^n (X_j - \overline{X}_n)^2}{\sigma^2(n-1)} + n\frac{(\overline{X}_n - \mu)^2}{\sigma^2}
\]
Което може да се запише чрез $S^2$:
\[
U = \underbrace{\frac{(n-1)S^2}{\sigma^2}}_{U_n} + \underbrace{\frac{n(\overline{X}_n - \mu)^2}{\sigma^2}}_{Z_n^2}
\]
Вече знаем, че $Z_n = \frac{\overline{X}_n - \mu}{\sigma/\sqrt{n}}$ е стандартно нормално $N(0,1)$, следователно квадратът му $Z_n^2$ е хи-квадрат разпределен с 1 степен на свобода, $\chi^2(1)$. Цялата сума $U$ е с $n$ степени на свобода. Благодарение на факта, че двете събираеми са независими, техните степени на свобода се сумират: $n-1$ степени от $U_n$ плюс $1$ степен от $Z_n^2$ дават общо $n$.

Да се върнем на конструирането на централната статистика. Ние искаме да използваме точковата оценка $\overline{X}_n$, да я центрираме изваждайки $\mu$, но тъй като не знаем $\sigma$, ще я разделим на нейната оценка $S / \sqrt{n}$:
\[
T_n = \frac{\overline{X}_n - \mu}{S / \sqrt{n}}
\]
За да намерим разпределението на $T_n$, можем да извършим следната еквилибристика — да умножим и разделим едновременно на неизвестното $\sigma$:
\[
T_n = \frac{\frac{\overline{X}_n - \mu}{\sigma / \sqrt{n}}}{\frac{S}{\sigma}} = \frac{Z_n}{\sqrt{\frac{S^2}{\sigma^2}}} = \frac{Z_n}{\sqrt{\frac{U_n}{n-1}}}
\]
В числителя имаме $Z_n \sim N(0,1)$, а в знаменателя имаме корен квадратен от $U_n \sim \chi^2(n-1)$, разделено на неговите степени на свобода. Тъй като $\overline{X}_n$ и $S^2$ са независими, $Z_n$ и $U_n$ също са независими. По дефиниция, такова отношение дефинира \textbf{разпределение на Стюдънт} (или t-разпределение) с $n-1$ степени на свобода:
\[
T_n = \frac{Z_n}{\sqrt{\frac{U_n}{n-1}}} = \frac{\overline{X}_n - \mu}{S/\sqrt{n}} \sim t(n-1)
\]
Получихме красив резултат: $T_n$ е монотонно намаляваща функция по $\mu$, а нейното разпределение $t(n-1)$ не зависи от $\mu$ и не съдържа неизвестното $\sigma$. Следователно, $T_n$ е централна статистика за $\mu$.

Тъй като t-разпределението също е симетрично (подобно на нормалното, но с малко по-тежки опашки), за да намерим доверителен интервал с ниво на доверие $\gamma$, избираме симетрични квантили $-q$ и $q$, където $q = q_{\frac{1+\gamma}{2}}$ е квантилът на $t(n-1)$ разпределението.
\[
\gamma = \mathbb{P}(I_1 < \mu < I_2) = \mathbb{P}(-q < T_n < q)
\]
По същия начин както в раздел~\ref{sec:ci-known-var}, решаваме спрямо $\mu$ и получаваме доверителния интервал:
\[
I_1 = \overline{X}_n - q \frac{S}{\sqrt{n}}, \quad I_2 = \overline{X}_n + q \frac{S}{\sqrt{n}}
\]

\subsection{Връзка с Централната гранична теорема (ЦГТ)}
Един важен коментар за разпределението на Стюдънт: когато обемът на извадката $n$ е достатъчно голям (в практиката често се приема $n \ge 30$), $T_n$ може да се приближи много добре чрез стандартно нормално разпределение:
\[
T_n \approx Z_n \sim N(0,1)
\]
Това е така, защото по Закона за големите числа, знаменателят $\sqrt{\frac{U_n}{n-1}} \approx 1$, което означава, че за голямо $n$ той почти не влияе на статистиката. Квантилите на t-разпределението при големи степени на свобода на практика съвпадат с квантилите на нормалното разпределение, което значително улеснява пресмятанията.

Но по-важното следствие е във връзка с \textbf{Централната гранична теорема (ЦГТ)}. Ако имате извадка от \emph{произволна} случайна величина $X$ (не задължително нормална), която има някакво математическо очакване $\E X = \mu$ и дисперсия $\Var X = \sigma^2$, то при голямо $n$, статистиката $T_n = \frac{\overline{X}_n - \mu}{\sigma/\sqrt{n}}$ е приблизително разпределена като $N(0,1)$.
Това е силата на математиката — не е нужно да знаем точното разпределение на данните (дали е експоненциално, поасоново или някакво друго), стига да разполагаме с достатъчно голяма извадка (обикновено $n \ge 30$ е достатъчно), за да третираме $T_n$ като стандартно нормално разпределение и да конструираме доверителни интервали за $\mu$. Дори когато $\sigma$ е неизвестно, можем да го заместим с оценката $S$, и за голямо $n$ приближението продължава да работи отлично.

\section{Проверка на хипотези}
\subsection{Основни понятия}
При проверката на статистически хипотези формулираме твърдения относно неизвестните параметри на изследвания модел. Разглеждаме две съревноваващи се хипотези:
\begin{itemize}
    \item \textbf{Нулева хипотеза} $H_0: \theta = \theta_0$
    \item \textbf{Алтернативна хипотеза} $H_1: \theta = \theta_1$
\end{itemize}
\begin{defn}[критична област]\label{def:critical-region}
Подмножеството на пространството на извадките $W \subseteq \mathbb{R}^n$, чрез което вземаме решение коя хипотеза да приемем на базата на извадката $\vec{X}$, наричаме \emph{критична област}.
\end{defn}

Правилото за вземане на решение е следното:
\begin{itemize}
    \item Ако реализацията на извадката попадне в критичната област, $\vec{X} \in W$, ние \textbf{отхвърляме} нулевата хипотеза $H_0$ и приемаме алтернативата $H_1$.
    \item Ако реализацията не попадне в критичната област, $\vec{X} \notin W$, ние \textbf{приемаме} (по-точно, не можем да отхвърлим) нулевата хипотеза $H_0$.
\end{itemize}

Тъй като данните са случайни, винаги има риск да вземем грешно решение. Възможни са два вида грешки:

\begin{defn}[грешки от първи и втори род]\label{def:errors}
\begin{itemize}
    \item \emph{Грешка от първи род:} да отхвърлим $H_0$, въпреки че тя е вярна. Вероятността за тази грешка бележим с $\alpha$ и наричаме \emph{ниво на значимост}:
    \[
    \mathbb{P}(\vec{X} \in W \given H_0) = \alpha_W = \alpha
    \]
    \item \emph{Грешка от втори род:} да приемем $H_0$, въпреки че тя е грешна (т.е. вярна е $H_1$). Вероятността за тази грешка бележим с $\beta$:
    \[
    \mathbb{P}(\vec{X} \notin W \given H_1) = \beta_W = \beta
    \]
\end{itemize}
\end{defn}

Целта в теорията на проверката на хипотези е да фиксираме грешката от първи род от дефиниция~\ref{def:errors}, тоест нивото на значимост $\alpha$ (например на $0{,}05$ или $0{,}01$), и при това фиксирано ограничение да изберем критична област по дефиниция~\ref{def:critical-region}, която минимизира вероятността за грешка от втори род $\beta_W$.

\begin{defn}[най-мощна критична област]\label{def:most-powerful}
Критичната област $W^*$ се нарича \emph{оптимална критична област} (о.к.о.), или \emph{най-мощна} (н.м.о.), при фиксирано $\alpha = \alpha_{W^*}$, ако за всяка друга критична област $W$ със същата грешка от първи род ($\alpha_W = \alpha_{W^*} = \alpha$) е изпълнено:
\[
\beta_{W^*} = \min_{\alpha_W = \alpha_{W^*} = \alpha} \beta_W .
\]
\end{defn}

Това е еквивалентно на максимизиране на мощността на критерия $1 - \beta_W = \mathbb{P}(\vec{X} \in W \given H_1)$.

\subsection{Лема на Нейман-Пиърсън}
Тази лема ни дава мощен инструмент за намиране на най-мощната критична област по дефиниция~\ref{def:most-powerful} при проверка на две прости хипотези.

Нека имаме случайна величина $X$ с плътност $f_X(x; \theta)$. Тогава съвместната плътност (функцията на правдоподобие) за извадката $\vec{X}$ е $L(\vec{x}, \theta)$. Означаваме функцията на правдоподобие при вярна нулева хипотеза с $L_0(\vec{x}) = L(\vec{x}, \theta_0)$, а при вярна алтернативна хипотеза с $L_1(\vec{x}) = L(\vec{x}, \theta_1)$.

\begin{keylem}[Нейман—Пиърсън]\label{lem:neyman-pearson}
Ако съществуват константа $k > 0$ и критична област $W^* \subseteq \mathbb{R}^n$ с грешка от първи род $\alpha$, такива че
\[
W^* \subseteq \{ \vec{x} \in \mathbb{R}^n : L_1(\vec{x}) \ge k L_0(\vec{x}) \}
\]
и за нейното допълнение
\[
(W^*)^c \subseteq \{ \vec{x} \in \mathbb{R}^n : L_1(\vec{x}) \le k L_0(\vec{x}) \} ,
\]
то $W^*$ е най-мощната критична област (н.м.о.) за проверка на хипотезата $H_0$ срещу алтернативата $H_1$ на ниво на значимост $\alpha$.
\end{keylem}

\begin{example}[Тест за средното на нормално разпределение]
Ще приложим\label{ex:13-1} лема~\ref{lem:neyman-pearson} в един класически пример. Нека имаме извадка $\vec{X} = (X_1, \dots, X_n)$ от нормално разпределение $X \sim N(\mu, \sigma^2)$, при което дисперсията $\sigma^2$ е \textbf{известна}.

Искаме да проверим нулевата хипотеза, че средното $\mu$ е равно на някаква стойност $\mu_0$, срещу алтернативата, че е равно на по-голяма стойност $\mu_1$:
\begin{itemize}
    \item $H_0: \mu = \mu_0$
    \item $H_1: \mu = \mu_1 \quad (\text{като приемаме, че } \mu_1 > \mu_0)$
\end{itemize}
Нивото на значимост (грешката от първи род) е фиксирано, например $\alpha = 0{,}05$.

Започваме като запишем функциите на правдоподобие за двете хипотези:
\[
L_0(\vec{x}) = \left(\frac{1}{\sqrt{2\pi}}\right)^n \frac{1}{\sigma^n} \prod_{j=1}^n e^{-\frac{(x_j - \mu_0)^2}{2\sigma^2}}
\]
\[
L_1(\vec{x}) = \left(\frac{1}{\sqrt{2\pi}}\right)^n \frac{1}{\sigma^n} \prod_{j=1}^n e^{-\frac{(x_j - \mu_1)^2}{2\sigma^2}}
\]
По лема~\ref{lem:neyman-pearson} търсим критична област от вида $W^* = \{ \vec{x} \in \mathbb{R}^n : L_1(\vec{x}) \ge k L_0(\vec{x}) \}$. Да видим какво представлява това неравенство, като заместим функциите на правдоподобие. Константите пред произведението се съкращават и неравенството придобива вида:
\[
W^* = \left\{ \vec{x} \in \mathbb{R}^n : e^{-\frac{\sum_{j=1}^n (x_j - \mu_1)^2}{2\sigma^2}} \ge k e^{-\frac{\sum_{j=1}^n (x_j - \mu_0)^2}{2\sigma^2}} \right\}
\]
Тъй като експоненциалната функция е монотонна, логаритмуваме двете страни:
\[
-\frac{1}{2\sigma^2} \sum_{j=1}^n (x_j - \mu_1)^2 \ge \ln k - \frac{1}{2\sigma^2} \sum_{j=1}^n (x_j - \mu_0)^2
\]
Умножаваме по $2\sigma^2$:
\[
- \sum_{j=1}^n (x_j - \mu_1)^2 \ge 2\sigma^2 \ln k - \sum_{j=1}^n (x_j - \mu_0)^2
\]
Разкриваме скобите във всяка от сумите:
\[
- \sum_{j=1}^n x_j^2 + 2\mu_1 \sum_{j=1}^n x_j - n\mu_1^2 \ge 2\sigma^2 \ln k - \sum_{j=1}^n x_j^2 + 2\mu_0 \sum_{j=1}^n x_j - n\mu_0^2
\]
Членовете $-\sum_{j=1}^n x_j^2$ от двете страни се съкращават. Групираме членовете, които съдържат $\sum_{j=1}^n x_j$ отляво, а останалите константи прехвърляме отдясно:
\[
2(\mu_1 - \mu_0) \sum_{j=1}^n x_j \ge 2\sigma^2 \ln k + n\mu_1^2 - n\mu_0^2
\]
От условието $\mu_1 > \mu_0$ знаем, че $2(\mu_1 - \mu_0) > 0$. Следователно, можем да разделим на този израз, без да се обръща знакът на неравенството:
\[
\sum_{j=1}^n x_j \ge \frac{2\sigma^2 \ln k + n\mu_1^2 - n\mu_0^2}{2(\mu_1 - \mu_0)}
\]
Цялата дясна страна на това неравенство е просто някаква нова константа, тъй като всички величини в нея ($\sigma, \mu_0, \mu_1, n, k$) са фиксирани числа. Нека я означим с $\tilde{K}$.
Така, формата на най-мощната критична област значително се опростява:
\[
W^* = \left\{ \vec{x} \in \mathbb{R}^n : \sum_{j=1}^n x_j \ge \underbrace{\frac{2\sigma^2 \ln k + n\mu_1^2 - n\mu_0^2}{2(\mu_1 - \mu_0)}}_{\tilde{K}} \right\}
\]
Това правило е много интуитивно: отхвърляме хипотезата за по-малкото средно $\mu_0$ в полза на по-голямото средно $\mu_1$, когато сумата от наблюденията надвиши определен праг $\tilde{K}$.

Остава единствено да определим стойността на константата $\tilde{K}$. Тя се намира от условието вероятността за грешка от първи род да бъде точно $\alpha$:
\[
\alpha = \mathbb{P}(\vec{X} \in W^* \given H_0) = \mathbb{P}\left( \sum_{j=1}^n X_j \ge \tilde{K} \given H_0 \right)
\]
Тъй като пресмятаме вероятността при условие, че $H_0$ е вярна, ние знаем, че всяко $X_j \sim N(\mu_0, \sigma^2)$. Следователно можем да центрираме и нормираме сумата, за да получим стандартно нормално разпределение:
\[
\mathbb{P}\left( \frac{\sum_{j=1}^n X_j - n\mu_0}{\sigma\sqrt{n}} \ge \frac{\tilde{K} - n\mu_0}{\sigma\sqrt{n}} \given H_0 \right) = \alpha
\]
Лявата страна в неравенството вече е стандартна нормална величина $Z \sim N(0,1)$. Нека означим дясната страна с $q_{1-\alpha}$, което е квантилът на нормалното разпределение, оставящ площ $\alpha$ в дясната опашка (или $1-\alpha$ в лявата). Имаме:
\[
\mathbb{P}\left( Z \ge \frac{\tilde{K} - n\mu_0}{\sigma\sqrt{n}} \right) = \alpha
\]
Следователно приравняваме границата на квантила:
\[
\frac{\tilde{K} - n\mu_0}{\sigma\sqrt{n}} = q_{1-\alpha}
\]
Квантилът $q_{1-\alpha}$ (например за $\alpha = 0{,}05$ от таблиците се намира стойността му) може да бъде намерен от статистическите таблици. От това уравнение еднозначно се намира константата $\tilde{K}$, с което най-мощната критична област $W^*$ е напълно определена.

\end{example}

\section{Задачи}

\begin{enumerate}
    \item При неизвестна дисперсия използвайте централната статистика
    \[
    T_n=\frac{\bar X_n-\mu}{S/\sqrt n}\sim t(n-1)
    \]
    и квантила $q=q_{(1+\gamma)/2}$, за да решите неравенството
    \[
    -q<T_n<q
    \]
    спрямо $\mu$ и да получите границите $I_1$ и $I_2$ на доверителния интервал с ниво на доверие $\gamma$.
\end{enumerate}


\chapter{Линейна регресия}
\section{Въведение в линейната регресия}

В тази лекция ще се запознаем с линейната регресия. Тя е един от най-основните инструменти в статистиката. Линейната регресия стои зад много съвременни инструменти и технологии, въпреки че често пъти е скрита вътре в самия модел, като например в невронните мрежи, рейтинговите системи и други.

Разбира се, както при всеки модел, ако предпоставките или допусканията не са изпълнени, лесно можете да стигнете до грешни изводи. Създаването на модели е изкуство, което изисква познаване на теорията и се овладява най-вече чрез практиката.

\subsection{Исторически контекст: Регрес към средното}
\begin{example}[Регрес към средното при ръста]\label{ex:14-1}
Началото на линейната регресия е свързано с Франсис Галтън, който се опитал да изучи зависимостта между ръста на бащите и ръста на техните синове. В своето изследване Галтън се опитал да предскаже височината на сина като функция на ръста на бащата. Забелязал е, че един добър модел може да се запише във вида:
\[
\text{Ръст на сина} = \text{Среден ръст} + \beta (\text{Ръст на бащата} - \text{Среден ръст}) + \varepsilon
\]
В изчисленията на Галтън коефициентът $\beta$ е бил приблизително $0{,}6$. Ако бащата е с 10 см по-висок от средния ръст за популацията, моделът предвижда отклонение
\[
0{,}6\cdot 10\text{ см}=6\text{ см}
\]
за сина: средно около 6 см над средния ръст, а не цели 10 см.

Това явление се нарича \textbf{регрес към средното}. То се забелязва нерядко и при други човешки характеристики, като например интелекта. Изключително умни хора имат тенденцията да имат наследници, чиито показатели са по-близо до средното ниво на популацията. Името „регресия“ произхожда именно от това „връщане“ към средната стойност.
\end{example}

Най-общо, един такъв модел може да се запише като:
\[
Y_k = \beta_0 + \beta_1 X_k + \varepsilon_k
\]
където $X_k$ е предикторът (например ръстът на бащата или дозата на дадено лекарство), $Y_k$ е откликът (ръстът на сина или реакцията към лекарството), $\beta_0$ и $\beta_1$ са линейни коефициенти, а $\varepsilon_k$ е стохастична грешка.

\section{Детерминистичен модел. Метод на най-малките квадрати}

Нека разгледаме детерминистичната оптимизационна задача, която стои зад линейната регресия. Имаме дадени $n$ на брой наблюдения, представляващи точки в равнината $(X_k, Y_k)$ за $k=1, \dots, n$. Търсим права линия с уравнение $y_k = b_0 + b_1 X_k$, която да моделира данните по възможно най-добрия начин. Грешката за дадено наблюдение $k$ е разликата между истинската стойност $Y_k$ и предвидената стойност от нашата права $\hat{Y}_k = b_0 + b_1 X_k$.

\begin{defn}[метод на най-малките квадрати]\label{def:ols}
\emph{Оценки по метода на най-малките квадрати} наричаме онези стойности на $b_0$ и $b_1$, които минимизират сумата от квадратите на грешките:
\[
\sum_{k=1}^n \varepsilon_k^2 = \sum_{k=1}^n (Y_k - b_0 - b_1 X_k)^2 \to \min .
\]
\end{defn}

Изборът на квадрата на грешката (вместо модул) исторически е бил мотивиран от изчислително удобство, тъй като функцията е гладка и лесно се диференцира, а също така съответства на Евклидовото разстояние. Освен това, от статистическа гледна точка, този подход води до добри оценки.

\subsection{Намиране на коефициентите}
За да минимизираме функцията, намираме нейните частни производни спрямо $b_0$ и $b_1$ и ги приравняваме на нула. Ще използваме стандартните означения $\bar{X} = \frac{1}{n}\sum X_k$ и $\bar{Y} = \frac{1}{n}\sum Y_k$ за средните аритметични стойности.

Диференцирайки по $b_0$ и разделяйки на $-2$, получаваме първото уравнение (нормално уравнение):
\[
\sum_{k=1}^n (Y_k - b_0 - b_1 X_k) = 0
\]
\[
\sum_{k=1}^n Y_k - n b_0 - b_1 \sum_{k=1}^n X_k = 0
\]
Изразявайки сумите чрез средните аритметични ($n \bar{Y} - n b_0 - b_1 n \bar{X} = 0$), след съкращаване на $n$ стигаме до връзката за $b_0$:
\[
b_0 = \bar{Y} - b_1 \bar{X}
\]

Диференцирайки по $b_1$, получаваме второто уравнение:
\[
\sum_{k=1}^n X_k (Y_k - b_0 - b_1 X_k) = 0
\]
\[
\sum_{k=1}^n X_k Y_k - b_0 n \bar{X} - b_1 \sum_{k=1}^n X_k^2 = 0
\]

Ако заместим $b_0$ във второто уравнение и извършим алгебричните преобразувания, можем да изведем явна формула за наклона $b_1$. Формулата може да се запише в няколко еквивалентни и много полезни вида:
\[
b_1 = \frac{\sum_{k=1}^n X_k Y_k - n \bar{X}\bar{Y}}{\sum_{k=1}^n X_k^2 - n \bar{X}^2} = \frac{\sum_{k=1}^n (X_k - \bar{X})(Y_k - \bar{Y})}{\sum_{k=1}^n (X_k - \bar{X})^2}
\]
Забележете, че горните изрази изключително много приличат на формулата за ковариация и дисперсия, откъдето идва и връзката с коефициента на корелация.

\subsection{Удобен запис за статистическия анализ}
За да изследваме свойствата на $b_0$ и $b_1$ по-лесно, ще въведем следното означение. Нека $A$ е сумата от квадратите на отклоненията на предикторите от тяхното средно:
\[
A = \sum_{k=1}^n (X_k - \bar{X})^2
\]
Въвеждаме и теглата $c_k$:
\[
c_k = \frac{X_k - \bar{X}}{A}
\]
Тогава оценката $b_1$ може елегантно да се запише като линейна комбинация на отклиците $Y_k$:
\[
\hat{\beta}_1 = b_1 = \sum_{k=1}^n c_k Y_k
\]
Оценката $\hat{\beta}_0$ може да се запише по аналогичен начин, замествайки $b_1$ в $b_0 = \bar{Y} - b_1 \bar{X}$:
\[
\hat{\beta}_0 = b_0 = \sum_{k=1}^n \left( \frac{1}{n} - \bar{X} c_k \right) Y_k
\]
Много важно свойство на константите $c_k$, което ще използваме често, е, че сумата им е нула:
\[
\sum_{k=1}^n c_k = \frac{\sum_{k=1}^n (X_k - \bar{X})}{A} = \frac{n\bar{X} - n\bar{X}}{A} = 0
\]

\section{Стохастичен модел}

Дотук разгледахме задачата детерминистично. По-общо, моделът се разглежда статистически. Предикторите $X_k$ се считат за детерминистични величини (например дозата от дадено лекарство, която ние контролираме и избираме), докато отклиците $Y_k$ са случайни величини. Стохастиката в $Y_k$ идва единствено от случайната грешка $\varepsilon_k$.

\begin{defn}[стохастичен модел на простата линейна регресия]\label{def:slr}
Истинският модел е
\[
Y_k = \beta_0 + \beta_1 X_k + \varepsilon_k ,
\]
при следните допускания за грешките $\varepsilon_k$:
\begin{enumerate}
    \item \textbf{Независимост:} случайните величини $(\varepsilon_1, \dots, \varepsilon_n)$ са взаимно независими.
    \item \textbf{Липса на систематична грешка:} $\E[\varepsilon_k] = 0$ за всяко $k$, откъдето $\E[Y_k] = \beta_0 + \beta_1 X_k$.
    \item \textbf{Хомоскедастичност (еднаква дисперсия):} $\Var(\varepsilon_k) = \sigma^2$ за всяко $k$.
    \item \textbf{Нормалност:} грешките са нормално разпределени, $\varepsilon_k \sim N(0, \sigma^2)$.
\end{enumerate}
\end{defn}

Тези допускания са силни. Второто означава, че средно моделът е точен; третото — че грешката има еднаква вариативност, независимо от стойността на предиктора $X_k$.

\textbf{Забележка относно допусканията:} Обобщения на този модел има в най-различни посоки (напр. хетероскедастичност). Винаги обаче трябва да наложите някаква структура на грешката, за да можете да правите изводи. Ако човек просто използва модела без да проверява дали тези допускания са валидни (както се е случвало често в практиката преди финансовата криза от 2008-2009 г.), той може да получи напълно погрешни и катастрофални резултати, защото моделът ще върне число, но то може да е напълно невалидно.

При допусканията на модела от дефиниция~\ref{def:slr} се оказва, че оценките $\hat{\beta}_0 = b_0$ и $\hat{\beta}_1 = b_1$, получени по метода от дефиниция~\ref{def:ols}, съвпадат с \textbf{максимално правдоподобните оценки} (MLE).

\section{Свойства на оценките}

Използвайки допусканията, ще изследваме свойствата на нашите оценки. Тъй като $b_1$ и $b_0$ са линейни комбинации на нормално разпределени случайни величини ($Y_k$), то самите те също ще имат нормално разпределение. Трябва само да намерим тяхното математическо очакване и дисперсия.

\subsection{Оценка на наклона $\hat{\beta}_1$ ($b_1$)}
Проверяваме дали оценката $b_1$ е неизместена (unbiased). Математическото очакване е:
\[
\E[b_1] = \E\left[ \sum_{k=1}^n c_k Y_k \right] = \sum_{k=1}^n c_k \E[Y_k] = \sum_{k=1}^n c_k (\beta_0 + \beta_1 X_k)
\]
Тъй като $\sum_{k=1}^n c_k = 0$, членът с $\beta_0$ отпада:
\[
\E[b_1] = \beta_1 \sum_{k=1}^n c_k X_k
\]
Сега нека пресметнем сумата $\sum c_k X_k$. Тъй като можем да добавим $-\bar{X}$ (което е константа) без да променим стойността на сумата (защото $\bar{X} \sum c_k = 0$):
\[
\sum_{k=1}^n c_k X_k = \sum_{k=1}^n \frac{X_k - \bar{X}}{A} X_k = \sum_{k=1}^n \frac{X_k - \bar{X}}{A} (X_k - \bar{X}) = \frac{\sum (X_k - \bar{X})^2}{A} = \frac{A}{A} = 1
\]
Следователно $\E[b_1] = \beta_1$, което доказва, че оценката е \textbf{неизместена}.

Дисперсията на $b_1$ намираме използвайки факта, че $Y_k$ са независими (тъй като $\varepsilon_k$ са независими) и всяко има дисперсия $\sigma^2$:
\[
\Var(b_1) = \Var\left( \sum_{k=1}^n c_k Y_k \right) = \sum_{k=1}^n \Var(c_k Y_k) = \sum_{k=1}^n c_k^2 \Var(Y_k) = \sum_{k=1}^n c_k^2 \sigma^2
\]
Заместваме $c_k$:
\[
\Var(b_1) = \sigma^2 \sum_{k=1}^n \left( \frac{X_k - \bar{X}}{A} \right)^2 = \sigma^2 \frac{\sum (X_k - \bar{X})^2}{A^2} = \frac{\sigma^2 A}{A^2} = \frac{\sigma^2}{A}
\]
В крайна сметка получаваме, че $b_1$ има следното нормално разпределение:
\[
b_1 \sim N\left(\beta_1, \frac{\sigma^2}{A}\right)
\]

\subsection{Оценка на пресечната точка $\hat{\beta}_0$ ($b_0$)}
Аналогично пресмятаме очакването на $b_0$:
\[
\E[b_0] = \sum_{k=1}^n \left( \frac{1}{n} - \bar{X} c_k \right) \E[Y_k] = \sum_{k=1}^n \left( \frac{1}{n} - \bar{X} c_k \right) (\beta_0 + \beta_1 X_k)
\]
Ако се разкрият скобите и се използват вече доказаните свойства ($\sum c_k = 0$ и $\sum c_k X_k = 1$), се получава, че $\E[b_0] = \beta_0$. Оценката $b_0$ също е неизместена.

Дисперсията на $b_0$ е:
\[
\Var(b_0) = \sum_{k=1}^n \left( \frac{1}{n} - \bar{X} c_k \right)^2 \Var(Y_k) = \sigma^2 \sum_{k=1}^n \left( \frac{1}{n} - \bar{X} c_k \right)^2
\]
След алгебрични преработки, които ще спестим за краткост, се получава:
\[
\Var(b_0) = \sigma^2 \left( \frac{1}{n} + \frac{\bar{X}^2}{A} \right)
\]
Следователно, разпределението на $b_0$ е:
\[
b_0 \sim N\left(\beta_0, \sigma^2 \left(\frac{1}{n} + \frac{\bar{X}^2}{A}\right)\right)
\]

\section{Оценяване на дисперсията $\sigma^2$}

Ако дисперсията $\sigma^2$ на грешките ни е предварително известна, то тогава можем веднага да използваме горните разпределения. Често пъти обаче $\sigma^2$ не е известна и следва да я оценим от данните.

Нека разгледаме стандартизираните грешки:
\[
Z_k = \frac{Y_k - \beta_0 - \beta_1 X_k}{\sigma} \sim N(0, 1)
\]
Тъй като тези величини са независими стандартни нормални, сумата от техните квадрати има хи-квадрат разпределение с $n$ степени на свобода:
\[
\sum_{k=1}^n Z_k^2 = \frac{1}{\sigma^2} \sum_{k=1}^n (Y_k - \beta_0 - \beta_1 X_k)^2 \sim \chi^2(n)
\]
Проблемът е, че ние не знаем истинските параметри $\beta_0$ и $\beta_1$. Можем да ги заместим с техните оценки $b_0$ и $b_1$, които сме пресметнали. Тъй като използваме две оценки (губим две степени на свобода), разпределението се променя на $\chi^2(n-2)$:
\[
\frac{1}{\sigma^2} \sum_{k=1}^n (Y_k - b_0 - b_1 X_k)^2 \sim \chi^2(n-2)
\]

Според Закона за големите числа (ЗГЧ), ако разделим случайна величина с $\chi^2(n-2)$ разпределение на нейните степени на свобода, то при $n \to \infty$ тази дроб клони към 1 (математическото очакване на $Z_1^2$). Това ни подсказва, че можем да конструираме \textbf{силна асимптотично състоятелна оценка} за дисперсията $\sigma^2$:
\[
\hat{\sigma}^2 = \frac{1}{n-2} \sum_{k=1}^n (Y_k - b_0 - b_1 X_k)^2
\]
В ситуации, в които $\sigma^2$ е неизвестна, просто я заместваме с $\hat{\sigma}^2$ във всички по-нататъшни разсъждения.

\section{Проверка на хипотези}

Често се налага да проверяваме хипотези относно линейните коефициенти. Най-често нулевата хипотеза е, че предикторът не влияе на отклика, т.е. че наклонът $\beta_1$ е нула: $H_0: \beta_1 = 0$. По-общо, можем да тестваме дали параметърът приема конкретна стойност $H_0: \beta_1 = \beta$, срещу алтернативата $H_1: \beta_1 \neq \beta$.

\subsection{Случай на известно $\sigma^2$}
Когато $\sigma^2$ е известно, използваме факта, че $b_1 \sim N(\beta_1, \sigma^2/A)$. При валидност на нулевата хипотеза $H_0: \beta_1 = \beta$, статистиката $Z$ е стандартно нормално разпределена:
\[
Z = \frac{b_1 - \beta}{\sqrt{\frac{\sigma^2}{A}}} \sim N(0, 1)
\]
Ако изберем ниво на значимост (грешка от първи род) $\alpha$, критичната област $W$, в която отхвърляме нулевата хипотеза, е симетрична:
\[
W = \left\{ |Z| \geq q_{1-\frac{\alpha}{2}} \right\}
\]
където $q_{1-\frac{\alpha}{2}}$ е квантилът на стандартното нормално разпределение.
Оттук лесно можем да построим и доверителен интервал за $\beta_1$ с вероятност $1-\alpha$:
\[
\beta_1 \in \left( b_1 - q_{1-\frac{\alpha}{2}}\sqrt{\frac{\sigma^2}{A}} \; , \; b_1 + q_{1-\frac{\alpha}{2}}\sqrt{\frac{\sigma^2}{A}} \right)
\]

\subsection{Случай на неизвестно $\sigma^2$}
Когато $\sigma^2$ не е известно, ние заместваме $\sigma^2$ с нейната оценка $\hat{\sigma}^2$. Разпределението на тестовата статистика се променя от нормално на \textbf{t-разпределение на Стюдънт} с $n-2$ степени на свобода:
\[
T = \frac{b_1 - \beta}{\sqrt{\frac{\hat{\sigma}^2}{A}}} \sim t(n-2)
\]
При конструиране на критична област и доверителни интервали, вместо квантилите на нормалното разпределение $q_{1-\alpha/2}$, трябва да използваме квантилите на $t$-разпределението $t_{1-\alpha/2, n-2}$, които имат по-тежки опашки.

\textbf{Важна бележка:} Когато $n$ е достатъчно голямо (обикновено $n \geq 32$, така че $n-2 \geq 30$), $t$-разпределението е много близко до стандартното нормално разпределение $N(0,1)$. В такива случаи можем безопасно да използваме нормалните квантили $q_{1-\alpha/2}$ като апроксимация, третирайки статистиката като нормална, просто замествайки $\sigma^2$ с $\hat{\sigma}^2$.

\subsection{Тестване на пресечната точка $\beta_0$}
Абсолютно същата логика може да се приложи и за тестване на хипотеза за пресечната точка: $H_0: \beta_0 = \beta^*$.
При известно $\sigma^2$ статистиката е:
\[
Z = \frac{b_0 - \beta^*}{\sqrt{\sigma^2\left(\frac{1}{n} + \frac{\bar{X}^2}{A}\right)}} \sim N(0, 1)
\]
Ако $\sigma^2$ не е известно, отново използваме $T$-статистика с $\hat{\sigma}^2$ и $t$-разпределение.


\renewcommand{\chaptername}{Упражнение}
\chapter{Комбинаторика и просто случайно блуждаене}
\renewcommand{\chaptername}{Лекция}
\section{Преговор по комбинаторика}
Тази среща е преговорно занятие върху материал, който се предполага познат от курса по дискретна математика. Основните понятия се припомнят чрез поредица от задачи; част от решенията се извеждат подробно, а останалите са оставени за самостоятелна работа. Сред припомнените инструменти е принципът за включване и изключване, който преподавателят свързва и със задачите от изпита по дискретна математика при доц. Марков.

\subsection{Правило за умножение и пермутации}
Нека разгледаме крайно множество от $n$ елемента. Броят на възможните наредби на тези елемента (пермутации без повторение) се дава от факториела на $n$:
\[ n! = n(n-1) \dots 2 \cdot 1 \]

\begin{supp}[формула на Стирлинг]\label{supp:stirling}
С нарастването на $n$ числата $n!$ растат толкова бързо, че прякото им пресмятане
става практически невъзможно. Формулата на Стирлинг дава удобно приближение:
\[ n! \sim \sqrt{2\pi n}\left(\frac{n}{e}\right)^{n}, \qquad n \to \infty . \]
\end{supp}
По дефиниция приемаме, че $0! = 1$.

\begin{keythm}[Принцип за включване и изключване]\label{thm:incl-excl}
За произволни крайни множества $A_1, \dots, A_n$:
\[
\left| \bigcup_{i=1}^{n} A_i \right|
 = \sum_{i} |A_i|
 - \sum_{i<j} |A_i \cap A_j|
 + \sum_{i<j<k} |A_i \cap A_j \cap A_k|
 - \dots
 + (-1)^{n-1} \left| \bigcap_{i=1}^{n} A_i \right| .
\]
\end{keythm}

За две множества това е познатото $|A \cup B| = |A| + |B| - |A \cap B|$: изваждаме
сечението, защото елементите в него са преброени два пъти. За три множества след като
извадим трите двойни сечения сме извадили тройното сечение един път в повече, затова го
добавяме обратно. Оттук и редуването на знаците.

Доказателството е по индукция: пишем $\bigcup_{i=1}^{n} A_i = \left(\bigcup_{i=1}^{n-1} A_i\right) \cup A_n$,
прилагаме формулата за две множества и след това индукционното предположение — веднъж за
набора $A_1, \dots, A_{n-1}$ и веднъж за набора от сеченията $A_i \cap A_n$. Записът е
тромав, но идеята е тази; важното е да се помни защо знаците се редуват.

\subsection{Вариации без повторение}
Ако от $n$ елемента избираме $k$ и \emph{редът има значение}, но повторения не се допускат,
броят на възможностите е
\[
V_n^k = n(n-1)\dots(n-k+1) = \frac{n!}{(n-k)!} .
\]
Разсъждението е последователно: за първата позиция има $n$ възможности, за втората $n-1$
(един елемент вече е употребен) и така до $k$-тата позиция, за която остават $n-k+1$
възможности. Такива наредени $k$-торки се наричат \emph{вариации} от клас $k$ без повторение.
Забележете, че $V_n^k$ е точно числителят в биномния коефициент по-долу: разликата е
дали редът се брои, или не.

\subsection{Биномни коефициенти (Комбинации)}
\begin{defn}[биномен коефициент]\label{def:binom}
Броят на начините да изберем $k$ елемента от множество с $n$ елемента, като редът на избиране няма значение (комбинации без повторение), се означава с
\[ \binom{n}{k} = C_n^k = \frac{n!}{k!(n-k)!} = \frac{n(n-1)\dots(n-k+1)}{k!} . \]
\end{defn}

Числата от дефиниция~\ref{def:binom} се наричат още биномни коефициенти, тъй като участват в теорема~\ref{thm:newton} за Нютоновия бином.

\begin{prop}[Основни свойства на биномните коефициенти]\label{prop:binom-props}
\begin{enumerate}
    \item \emph{Симетрия:} $\displaystyle \binom{n}{k} = \binom{n}{n-k}$.
    \item \emph{Рекурентна връзка (Правило на Паскал):} $\displaystyle \binom{n}{k} = \binom{n-1}{k-1} + \binom{n-1}{k}$.
    \item \emph{Сума:} $\displaystyle \sum_{k=0}^{n} \binom{n}{k} = 2^n$.
    \item \emph{Знакопроменлива сума:} $\displaystyle \sum_{k=0}^{n} (-1)^k \binom{n}{k} = 0$.
\end{enumerate}
\end{prop}

Свойствата в твърдение~\ref{prop:binom-props}, и в частност правилото на Паскал, позволяват лесно пресмятане на биномните коефициенти и конструирането на Триъгълника на Паскал. Сумата $2^n$ представлява броя на всички подмножества на множество с $n$ елемента; знакопроменливата сума (получена от Нютоновия бином при $a=1$, $b=-1$) означава, че броят на подмножествата с четен брой елементи е равен на броя на подмножествата с нечетен брой елементи.

\subsection{Нютонов бином}

\begin{thm}[Нютонов бином]\label{thm:newton}
\[ (a+b)^n = \sum_{k=0}^{n} \binom{n}{k} a^k b^{n-k} = \binom{n}{0} b^n + \binom{n}{1} ab^{n-1} + \dots + \binom{n}{n} a^n \]
\end{thm}

\subsection{Разпределяне на частици по клетки}
Често в комбинаториката се срещат задачи за разпределяне на частици по клетки.
Ако имаме $k$ различими частици и $n$ клетки, броят на възможните разпределения е:
\[ n^k \]
Това е така, защото всяка от $k$-те частици може да бъде поставена в коя да е от $n$-те клетки напълно независимо от останалите.

\begin{example}[Разпределяне без празна клетка]\label{ex:15-surj}
Нека $k \ge n$ и разпределяме $k$ различими частици в $n$ различими клетки така, че
\emph{никоя клетка да не остане празна}. Броят на такива разпределения е
\[
\sum_{m=0}^{n} (-1)^m \binom{n}{m} (n-m)^k .
\]

Пряко преброяване не работи: условието „няма празна клетка“ не се разлага на независими
избори. Затова броим допълнението чрез принцип~\ref{thm:incl-excl}. Нека
$A_i$ е множеството от разпределения, при които $i$-тата клетка е празна. Тогава
разпределенията с поне една празна клетка са $\left| \bigcup_{i=1}^n A_i \right|$.

Ако фиксираме $m$ конкретни клетки и искаме всички те да са празни, остават $n-m$ клетки за
$k$-те частици, тоест $(n-m)^k$ разпределения; а $m$-те клетки се избират по
$\binom{n}{m}$ начина. По принципа за включване и изключване
\[
\left| \bigcup_{i=1}^n A_i \right| = \sum_{m=1}^{n} (-1)^{m-1} \binom{n}{m} (n-m)^k ,
\]
и като извадим това от общия брой $n^k$, получаваме твърдението (членът $m=0$ дава точно
$n^k$).
\end{example}

Този пример е добра илюстрация защо изобщо съществува принципът за включване и изключване:
когато налагаме условие от вида „нищо не е празно“, отделните случаи се припокриват и
трябва да се коригират последователно.

\subsection{Уравнения с цели неотрицателни числа (Метод „Звезди и прегради“)}
Разглеждаме уравнението:
\[ x_1 + x_2 + \dots + x_k = n \]
където $x_i$ са цели неотрицателни числа ($x_i \geq 0$ за всяко $i=1,\dots,k$) и $n \geq 0$ е цяло число.
Търсим броя на решенията на това уравнение. Можем да си представим, че имаме $n$ неразличими обекта (единици или „звезди“), които искаме да разпределим в $k$ клетки. За да разделим обектите на $k$ групи, ни трябват $k-1$ „прегради“ (или „черти“).

Общият брой на позициите за звездите и преградите е $n + k - 1$. От тези позиции трябва да изберем $k-1$ места, на които да поставим преградите (или еквивалентно, да изберем $n$ места за звездите).

\begin{prop}[Звезди и прегради]\label{prop:stars-bars}
Броят на решенията на уравнението $x_1 + x_2 + \dots + x_k = n$ в цели неотрицателни числа е
\[ \binom{n+k-1}{k-1} = \binom{n+k-1}{n} . \]
\end{prop}

\begin{cor}[Комбинации с повторение]\label{cor:multiset}
Броят на начините да изберем $k$ елемента от $n$-елементно множество, когато редът няма
значение, но \emph{повторенията са позволени} (тоест броят на $k$-елементните
мултимножества), е
\[ \binom{n+k-1}{k} . \]
\end{cor}

\begin{proof}
Да изберем мултимножество означава да кажем колко пъти участва всеки от $n$-те елемента:
това е избор на цели неотрицателни числа $x_1, \dots, x_n$ с $x_1 + \dots + x_n = k$.
По твърдение~\ref{prop:stars-bars} с $k$ на брой звезди и $n$ клетки броят е
$\binom{k+n-1}{n-1} = \binom{n+k-1}{k}$.
\end{proof}

Сравнете двете съседни формули: $\binom{n}{k}$ брои подмножествата (без повторение), а
$\binom{n+k-1}{k}$ — мултимножествата (с повторение).

Твърдение~\ref{prop:stars-bars} е изключително важно и се прилага в много контексти, включително за намиране на броя на разпределенията на $n$ неразличими частици в $k$ различими клетки (където редът на частиците в дадена клетка няма значение).

\section{Просто случайно блуждаене}

Ще разгледаме един класически модел в теорията на вероятностите — просто случайно блуждаене (Simple Random Walk).
Нека една частица се движи по целочислените точки на реалната права. Тя стартира от началото на координатната система (точка 0) в момент $t=0$. Във всеки следващ дискретен момент от времето $t=1, 2, \dots$ частицата прави стъпка надясно (към $+1$) или наляво (към $-1$).
Нека $\xi_i \in \{+1, -1\}$ е стойността на $i$-тата стъпка. Тогава позицията на частицата след $n$ стъпки е:
\[ S_n = \xi_1 + \xi_2 + \dots + \xi_n \]
като по дефиниция приемаме $S_0 = 0$.

Можем да представим блуждаенето графично в координатната равнина $(t, S_t)$, където по хоризонталната ос нанасяме времето (броя стъпки), а по вертикалната — позицията на частицата. Тогава блуждаенето представлява начупена линия, започваща от $(0,0)$. Всяка отсечка от тази линия свързва точките $(i-1, S_{i-1})$ и $(i, S_i)$.

\subsection{Брой на всички пътища}
Ако разгледаме блуждаене с дължина $n$, то броят на всички възможни пътища (траектории) е $2^n$, тъй като на всяка от $n$-те стъпки имаме точно 2 възможни избора ($+1$ или $-1$).

За да стигнем от точка $(0,0)$ до точка $(n, x)$, трябва броят на стъпките надясно ($+1$) и наляво ($-1$) да са съответно $p$ и $q$, такива че да са изпълнени следните две условия:
\begin{align*}
p + q &= n \\
p - q &= x
\end{align*}
Като съберем двете уравнения, намираме:
\[ 2p = n + x \implies p = \frac{n+x}{2} \]
За да съществува такъв път, е необходимо и достатъчно $n+x$ да бъде четно число и $-n \leq x \leq n$.
\begin{prop}[Брой пътища до дадена точка]\label{prop:paths-count}
Ако $n+x$ е четно и $-n \le x \le n$, броят на пътищата от $(0,0)$ до $(n,x)$ е равен на броя на начините да изберем $p = \frac{n+x}{2}$ стъпки надясно от общо $n$ налични стъпки:
\[ N_n(x) = \binom{n}{p} = \binom{n}{\frac{n+x}{2}} . \]
\end{prop}

\section{Принцип на отражението}
Често ни интересуват пътища, които удовлетворяват определени условия, например пътища, които остават неотрицателни ($S_i \geq 0$ за всяко $i$). За да преброим такива пътища, използваме Принципа на отражението, въведен от Дезире Андре.

\begin{keythm}[Принцип на отражението]\label{thm:reflection}
Нека $x, y > 0$. Броят на пътищата от $(0,x)$ до $(n,y)$, които докосват или пресичат хоризонталната ос (т.е. $S_i = 0$ за някое $i$), е равен на общия брой пътища от $(0,-x)$ до $(n,y)$.
\end{keythm}

\begin{lecfigure}[Принципът на отражението. Плътната линия е „лош“ път от $(0,x)$ до
$(n,y)$; $t_0$ е първият момент, в който той докосва оста. Отразявайки началната част до момента $t_0$ спрямо оста (прекъснатата линия), получаваме път от $(0,-x)$ до $(n,y)$.
Съответствието е биекция.\label{fig:reflection}]
\begin{tikzpicture}[x=0.62cm, y=0.52cm]
  % axes
  \draw[axisline] (-0.4,0) -- (9.4,0) node[right, font=\small] {$t$};
  \draw[axisline] (0,-3.6) -- (0,3.9) node[above, font=\small] {$S$};
  % the "bad" path: (0,2) -> touches 0 at t=4 -> (8,2)
  \draw[walk] (0,2) -- (1,3) -- (2,2) -- (3,1) -- (4,0) -- (5,1) -- (6,2) -- (7,1) -- (8,2);
  % reflection of the initial segment across S = 0
  \draw[walkrefl] (0,-2) -- (1,-3) -- (2,-2) -- (3,-1) -- (4,0);
  % guides
  \draw[guide] (4,0) -- (4,-3.4);
  \node[font=\small] at (4,-3.9) {$t_0$};
  % endpoints
  \fill[thmframe] (0,2) circle (2.1pt) node[left, font=\small, black] {$x$};
  \fill[thmframe] (8,2) circle (2.1pt) node[right, font=\small, black] {$y$};
  \fill[keyframe] (0,-2) circle (2.1pt) node[left, font=\small, black] {$-x$};
  \fill[gray!70] (4,0) circle (1.8pt);
\end{tikzpicture}
\end{lecfigure}

Нека наречем пътищата, които докосват оста, „лоши“ (фигура~\ref{fig:reflection}). Вземаме един произволен такъв лош път. Нека $t_0$ е първият момент, в който пътят докосва нулата (т.е. $S_{t_0} = 0$, но $S_i > 0$ за $i < t_0$). Принципът на отражението гласи, че ако отразим симетрично началната част от пътя (от $t=0$ до $t=t_0$) спрямо абсцисната ос $S=0$, ще получим нов път, който започва от $(0,-x)$, докосва нулата в $t_0$ и продължава по оригиналния път до $(n,y)$.

Съществено е, че отразяваме спрямо \emph{първия} момент на докосване. Ако избирахме
произволен момент, в който пътят е на нулата, един и същ лош път би давал по няколко
различни образа и съответствието нямаше да е еднозначно. Първият момент е определен
еднозначно от самия път, затова изображението е коректно дефинирано; обратно, при даден
път от $(0,-x)$ до $(n,y)$ неговият пръв момент на докосване възстановява еднозначно
оригинала.

Тъй като всеки път от $(0,-x)$ до $(n,y)$ (при $x,y > 0$) задължително пресича оста $S=0$ в някакъв момент (започва под нея и завършва над нея), съществува взаимно еднозначно съответствие (биекция) между лошите пътища от $(0,x)$ до $(n,y)$ и множеството на всички възможни пътища от $(0,-x)$ до $(n,y)$. Това доказва твърдението.

\begin{supp}[Отражение спрямо произволно ниво]\label{supp:reflection-level}
Теорема~\ref{thm:reflection} е формулирана за ниво $0$ и изисква $x, y > 0$. Същият
аргумент работи спрямо кое да е ниво $a$, стига началната и крайната точка да лежат от
една и съща страна на него: ако $x, y > a$, то броят на пътищата от $(0,x)$ до $(n,y)$,
които докосват ниво $a$, е равен на общия брой пътища от $(0,\, 2a - x)$ до $(n,y)$.
Наистина, $2a-x$ е огледалният образ на $x$ спрямо $a$, и отразяването на началната част
до първия момент на докосване на ниво $a$ дава същата биекция.

Тази форма ни е нужна веднага по-долу: там началната и крайната точка са на самата ос
($x = y = 0$), така че теорема~\ref{thm:reflection} не се прилага пряко, а работим с
ниво $a = -1$.
\end{supp}

\subsection{Пътища, които остават неотрицателни}
Нека разгледаме пътищата, стартиращи от $(0,0)$ и завършващи в $(2n, 0)$. Според предходното твърдение общият им брой е $N_{2n}(0) = \binom{2n}{n}$.

Интересуваме се от броя на пътищата от $(0,0)$ до $(2n, 0)$, които остават неотрицателни през цялото време, т.е. $S_i \geq 0$ за $i=1, \dots, 2n$.
За да намерим техния брой, ще преброим „лошите“ пътища — тези, за които $S_i = -1$ за поне едно $i$ — и ще ги извадим от общия брой пътища.

Тук началната и крайната точка лежат на самата ос, затова теорема~\ref{thm:reflection}
не може да се приложи буквално. Използваме вариантa от допълнение~\ref{supp:reflection-level}
с ниво $a = -1$: началото $0$ и краят $0$ са от една и съща страна на $-1$, а „лоши“ са
точно пътищата, които докосват ниво $-1$. Огледалният образ на началната точка спрямо
$-1$ е $2(-1) - 0 = -2$, следователно лошите пътища са в биекция с \emph{всички} пътища
от $(0,-2)$ до $(2n,0)$.

Броят на пътищата от $(0,-2)$ до $(2n,0)$ се намира като решим:
\begin{align*}
p + q &= 2n \\
p - q &= 2
\end{align*}
Оттук $2p = 2n + 2 \implies p = n+1$. Следователно броят им е $\binom{2n}{n+1}$.

Тогава броят на „добрите“ пътища (които никога не достигат $-1$, следователно $S_i \geq 0$ за всички $i$) от $(0,0)$ до $(2n,0)$ е:
\[ \binom{2n}{n} - \binom{2n}{n+1} \]
Тази разлика може да се опрости алгебрично по следния начин:
\begin{align*}
\binom{2n}{n} - \binom{2n}{n+1} &= \frac{(2n)!}{n!n!} - \frac{(2n)!}{(n+1)!(n-1)!} \\
&= \frac{(2n)!}{n!(n-1)!} \left( \frac{1}{n} - \frac{1}{n+1} \right) \\
&= \frac{(2n)!}{n!(n-1)!} \left( \frac{n+1 - n}{n(n+1)} \right) \\
&= \frac{(2n)!}{n!(n-1)!} \frac{1}{n(n+1)} \\
&= \frac{1}{n+1} \frac{(2n)!}{n!n!} \\
&= \frac{1}{n+1} \binom{2n}{n}
\end{align*}
Този забележителен резултат представлява $n$-тото число на Каталан ($C_n$). Числата на Каталан имат огромно значение в комбинаториката, тъй като те броят множество различни комбинаторни структури (например правилни скобови изрази, триангулации на изпъкнали многоъгълници и дървета).

\subsection{Пътища с условие за строга положителност}
Ако разгледаме по-строго условие — да намерим броя на пътищата от $(0,0)$ до $(2n,0)$, които се връщат в нулата за първи път чак на последната стъпка $2n$ (т.е. $S_i > 0$ за $1 \leq i \leq 2n-1$ и $S_{2n} = 0$).

Тогава първата стъпка на такъв път задължително е нагоре към $(1,1)$, а предпоследната позиция задължително е $(2n-1, 1)$, откъдето прави стъпка надолу към $(2n,0)$.
Броят на тези пътища съвпада с броя на пътищата от $(1,1)$ до $(2n-1, 1)$, които никога не докосват нулата (остават строго положителни, $S_i > 0$).

Ако транслираме координатната система с вектор $(-1,-1)$, това е равносилно на намирането на броя на пътищата от $(0,0)$ до $(2n-2, 0)$, които остават неотрицателни (т.е. не слизат под новото ниво 0).
Според предходната точка, този брой е равен на съответното число на Каталан за $2n-2$ стъпки:
\[ \frac{1}{n} \binom{2n-2}{n-1} = C_{n-1} \]
което е $(n-1)$-тото число на Каталан.

\section{Задачи}

Следващите условия възстановяват задачите, които са формулирани устно или са записани на дъската. Запазена е номерацията от занятието.

Условията на задачи от 6 до 9 и 11 не са произнесени или записани в наличните материали: преподавателят се позовава на предварително раздаден лист със задачи и разработва само задача 10. Затова липсващите условия не могат да бъдат възстановени достоверно.

\begin{enumerate}
    \item[1.] Припомнете принципа за включване и изключване за $n$ крайни множества. Доказателството е по индукция, като индукционното предположение трябва да се приложи и към подходящото семейство от сечения; за този курс е важно преди всичко да се помни кога и как се използва принципът.

    \item[2.] Нека $M$ е множество с $n$ елемента. Намерете броя на:
    \begin{enumerate}
        \item $k$-елементните подмножества на $M$;
        \item наредените $k$-торки от различни елементи на $M$;
        \item наредените $k$-торки от елементи на $M$, когато повторенията са позволени.
    \end{enumerate}

    \item[3.] Намерете броя на решенията на
    \[
    x_1+\dots+x_k=n
    \]
    първо при $x_1,\dots,x_k\in\mathbb{N}$, а след това при $x_1,\dots,x_k\in\mathbb{N}_0$. Използвайте метода „звезди и прегради“.

    \item[4.] Разпределяме $k$ различими частици в $n$ различими клетки. Намерете броя на разпределенията, когато:
    \begin{enumerate}
        \item всяка клетка съдържа най-много една частица;
        \item няма ограничение за броя частици в клетка;
        \item няма празна клетка, като $k\ge n$. За последната точка използвайте принципа за включване и изключване.
    \end{enumerate}

    \item[5.] Разпределяме $k$ неразличими частици в $n$ различими клетки. Намерете броя на разпределенията, когато:
    \begin{enumerate}
        \item всяка клетка съдържа най-много една частица;
        \item няма ограничение за броя частици в клетка;
        \item няма празна клетка. Последната точка е оставена за домашна работа.
    \end{enumerate}

    \item[10.] Разгледайте просто случайно блуждаене, което започва от $(0,0)$ и на всяка стъпка се изменя с $+1$ или $-1$. Пребройте всички пътища до $(2n,0)$, а след това пътищата, които остават неотрицателни. За второто преброяване приложете принципа на отражението.
\end{enumerate}


\end{document}
